\chapter{Ñeembucú}\label{dept:Ñeembucú}
\mapaparaguai{12}{Ñeembucú}
\vfill\clearpage
\section{Neembucu /\allowbreak  pacote 14}

O pacote 14 cobre o departamento , com 16 localidades. Neembucu e
planicie aluvial no sudoeste paraguaio, marcada por humedales, o Rio
Paraguai e fronteira com a Argentina.

\begin{itemize}
\tightlist
\item
  Pilar e a capital, com Porto fluvial de aguas profundas operavel em
  secas extremas e polo textil.
\item
  Vocacao de pecuaria extensiva, industria textil e servicos portuarios.
\item
  Lacuna de solar/\allowbreak pluviometria da capital fechada com dados NASA POWER
  (2001-2020).
\end{itemize}

\subsection{Ficha climática de Pilar (capital
departamental)}

Fonte: NASA POWER Climatology API, período 2001-2020, coordenadas
-26.87°S /\allowbreak  -58.31°W (acesso em 2026-03-20).

\textbf{Irradiância solar global --- ALLSKY\_SFC\_SW\_DWN (kWh/\allowbreak m²/\allowbreak dia):}

\begin{longtable}[]{@{}
  >{\raggedright\arraybackslash}p{(\linewidth - 22\tabcolsep) * \real{0.0833}}
  >{\raggedright\arraybackslash}p{(\linewidth - 22\tabcolsep) * \real{0.0833}}
  >{\raggedright\arraybackslash}p{(\linewidth - 22\tabcolsep) * \real{0.0833}}
  >{\raggedright\arraybackslash}p{(\linewidth - 22\tabcolsep) * \real{0.0833}}
  >{\raggedright\arraybackslash}p{(\linewidth - 22\tabcolsep) * \real{0.0833}}
  >{\raggedright\arraybackslash}p{(\linewidth - 22\tabcolsep) * \real{0.0833}}
  >{\raggedright\arraybackslash}p{(\linewidth - 22\tabcolsep) * \real{0.0833}}
  >{\raggedright\arraybackslash}p{(\linewidth - 22\tabcolsep) * \real{0.0833}}
  >{\raggedright\arraybackslash}p{(\linewidth - 22\tabcolsep) * \real{0.0833}}
  >{\raggedright\arraybackslash}p{(\linewidth - 22\tabcolsep) * \real{0.0833}}
  >{\raggedright\arraybackslash}p{(\linewidth - 22\tabcolsep) * \real{0.0833}}
  >{\raggedright\arraybackslash}p{(\linewidth - 22\tabcolsep) * \real{0.0833}}@{}}
\toprule\noalign{}
\begin{minipage}[b]{\linewidth}\raggedright
Jan
\end{minipage} & \begin{minipage}[b]{\linewidth}\raggedright
Fev
\end{minipage} & \begin{minipage}[b]{\linewidth}\raggedright
Mar
\end{minipage} & \begin{minipage}[b]{\linewidth}\raggedright
Abr
\end{minipage} & \begin{minipage}[b]{\linewidth}\raggedright
Mai
\end{minipage} & \begin{minipage}[b]{\linewidth}\raggedright
Jun
\end{minipage} & \begin{minipage}[b]{\linewidth}\raggedright
Jul
\end{minipage} & \begin{minipage}[b]{\linewidth}\raggedright
Ago
\end{minipage} & \begin{minipage}[b]{\linewidth}\raggedright
Set
\end{minipage} & \begin{minipage}[b]{\linewidth}\raggedright
Out
\end{minipage} & \begin{minipage}[b]{\linewidth}\raggedright
Nov
\end{minipage} & \begin{minipage}[b]{\linewidth}\raggedright
Dez
\end{minipage} \\
\midrule\noalign{}
\endhead
\bottomrule\noalign{}
\endlastfoot
6.76 & 6.22 & 5.34 & 4.20 & 3.30 & 2.73 & 3.14 & 3.88 & 4.59 & 5.48 &
6.49 & 6.76 \\
\end{longtable}

Média anual: 4.90 kWh/\allowbreak m²/\allowbreak dia. Inclinação recomendada: 27° Norte (anual);
37° inverno; 17° verão.

\textbf{Precipitação --- PRECTOTCORR (mm/\allowbreak dia):}

\begin{longtable}[]{@{}
  >{\raggedright\arraybackslash}p{(\linewidth - 22\tabcolsep) * \real{0.0833}}
  >{\raggedright\arraybackslash}p{(\linewidth - 22\tabcolsep) * \real{0.0833}}
  >{\raggedright\arraybackslash}p{(\linewidth - 22\tabcolsep) * \real{0.0833}}
  >{\raggedright\arraybackslash}p{(\linewidth - 22\tabcolsep) * \real{0.0833}}
  >{\raggedright\arraybackslash}p{(\linewidth - 22\tabcolsep) * \real{0.0833}}
  >{\raggedright\arraybackslash}p{(\linewidth - 22\tabcolsep) * \real{0.0833}}
  >{\raggedright\arraybackslash}p{(\linewidth - 22\tabcolsep) * \real{0.0833}}
  >{\raggedright\arraybackslash}p{(\linewidth - 22\tabcolsep) * \real{0.0833}}
  >{\raggedright\arraybackslash}p{(\linewidth - 22\tabcolsep) * \real{0.0833}}
  >{\raggedright\arraybackslash}p{(\linewidth - 22\tabcolsep) * \real{0.0833}}
  >{\raggedright\arraybackslash}p{(\linewidth - 22\tabcolsep) * \real{0.0833}}
  >{\raggedright\arraybackslash}p{(\linewidth - 22\tabcolsep) * \real{0.0833}}@{}}
\toprule\noalign{}
\begin{minipage}[b]{\linewidth}\raggedright
Jan
\end{minipage} & \begin{minipage}[b]{\linewidth}\raggedright
Fev
\end{minipage} & \begin{minipage}[b]{\linewidth}\raggedright
Mar
\end{minipage} & \begin{minipage}[b]{\linewidth}\raggedright
Abr
\end{minipage} & \begin{minipage}[b]{\linewidth}\raggedright
Mai
\end{minipage} & \begin{minipage}[b]{\linewidth}\raggedright
Jun
\end{minipage} & \begin{minipage}[b]{\linewidth}\raggedright
Jul
\end{minipage} & \begin{minipage}[b]{\linewidth}\raggedright
Ago
\end{minipage} & \begin{minipage}[b]{\linewidth}\raggedright
Set
\end{minipage} & \begin{minipage}[b]{\linewidth}\raggedright
Out
\end{minipage} & \begin{minipage}[b]{\linewidth}\raggedright
Nov
\end{minipage} & \begin{minipage}[b]{\linewidth}\raggedright
Dez
\end{minipage} \\
\midrule\noalign{}
\endhead
\bottomrule\noalign{}
\endlastfoot
4.96 & 4.29 & 4.68 & 5.56 & 3.49 & 2.37 & 1.55 & 1.29 & 2.25 & 5.02 &
5.91 & 5.24 \\
\end{longtable}

Média anual: \textasciitilde1.416 mm/\allowbreak ano. Estação chuvosa Out-Jan; seca
Jul-Ago.
\clearpage
\subsection*{Dados Departamentais}
\subsubsection{Desenvolvimento Humano e
Social}

\subsection{IDH Departamental}

\begin{longtable}[]{@{}ll@{}}
\toprule\noalign{}
Indicador & Valor \\
\midrule\noalign{}
\endhead
\bottomrule\noalign{}
\endlastfoot
IDH & 0,710 \\
Classificação IDH & alto \\
Ranking nacional (1=melhor) & 8/\allowbreak 18 \\
Esperança de vida & 73,0 anos \\
Anos de escolaridade média & 8,2 anos \\
RNB per capita (USD PPA) & 8.800 \\
\end{longtable}

\subsection{Indicadores de Pobreza}

\begin{longtable}[]{@{}ll@{}}
\toprule\noalign{}
Indicador & Valor \\
\midrule\noalign{}
\endhead
\bottomrule\noalign{}
\endlastfoot
Taxa de pobreza total (\%) & 28,8\% \\
Taxa de pobreza extrema (\%) & 6,5\% \\
Índice de Gini & 0,445 \\
\end{longtable}

\subsection{Infraestrutura Básica}

\begin{longtable}[]{@{}ll@{}}
\toprule\noalign{}
Indicador & Valor \\
\midrule\noalign{}
\endhead
\bottomrule\noalign{}
\endlastfoot
Acesso a água potável (\%) & 79,1\% \\
Acesso a saneamento (\%) & 77,4\% \\
\end{longtable}

\subsubsection{Segurança Pública}

\subsection{Indicadores}

\begin{longtable}[]{@{}
  >{\raggedright\arraybackslash}p{(\linewidth - 2\tabcolsep) * \real{0.6111}}
  >{\raggedright\arraybackslash}p{(\linewidth - 2\tabcolsep) * \real{0.3889}}@{}}
\toprule\noalign{}
\begin{minipage}[b]{\linewidth}\raggedright
Indicador
\end{minipage} & \begin{minipage}[b]{\linewidth}\raggedright
Valor
\end{minipage} \\
\midrule\noalign{}
\endhead
\bottomrule\noalign{}
\endlastfoot
Taxa de homicídios (por 100k hab) & 0,2--0,4 (2019--2020, fonte INE\footnote{Instituto Nacional de Estadística (Instituto Nacional de Estatística), órgão oficial de dados demográficos e censitários.} ---
mais baixa do país) \\
Crimes registrados (total/\allowbreak ano) & muito baixo (população pequena e
ambiente rural) \\
Número de comisarías & \textasciitilde6 (estimativa departamental) \\
Índice segurança relativo & alto \\
\end{longtable}

\subsubsection{Mercado de Terra Rural}

\subsection{Preços por Tipo de Terra}

\begin{longtable}[]{@{}llll@{}}
\toprule\noalign{}
Tipo & Preço (USD/\allowbreak ha) & Faixa & Tendência \\
\midrule\noalign{}
\endhead
\bottomrule\noalign{}
\endlastfoot
Agrícola alta produtividade & 3.000 & 1.800 -- 4.500 & ↑ \\
Pastagem & 1.500 & 900 -- 2.200 & → \\
Mata /\allowbreak  reserva & 600 & 300 -- 1.000 & → \\
\end{longtable}

\subsection{Características do
Mercado}

\begin{longtable}[]{@{}ll@{}}
\toprule\noalign{}
Parâmetro & Valor \\
\midrule\noalign{}
\endhead
\bottomrule\noalign{}
\endlastfoot
Valorização anual média (\%) & 4--6\% \\
Lote mínimo rural (ha) & 20 \\
Imposto territorial rural (\%) & 1\% sobre valor fiscal \\
Liquidez do mercado & baixa \\
\end{longtable}
\clearpage
\newpage
\secaoDiagnostico{Alberdi}{\mapaConteudo{-26.1666}{-58.1333}}{Alberdi é um enclave estratégico de resiliência no sul do Paraguai.
Oferece um dos melhores níveis de isolamento demográfico e tático do
país, protegida por vastos humedais e longe dos eixos de conflito
primários. Sua força reside na abundância de proteína animal e na
infraestrutura de suprimentos importados. O desafio crítico é a gestão
hídrica; a sobrevivência em Alberdi exige a escolha de terrenos em cotas
seguras e a manutenção de autonomia energética, aproveitando a nova rota
asfáltica que garante logística eficiente em tempos de paz.

\subsubsection{Por que sim}

\begin{itemize}
\tightlist
\item
  Isolamento geográfico superior e baixa visibilidade estratégica.
\item
  Abundância de recursos alimentares (gado e pesca).
\item
  Comunidade pacífica e resiliente.
\item
  Conectividade asfáltica moderna (Ruta PY19).
\end{itemize}

\subsubsection{Por que não}

\begin{itemize}
\tightlist
\item
  Risco crítico de inundações fluviais extraordinárias.
\item
  Dependência de infraestrutura de defesa costeira.
\item
  Restrições da Lei de Fronteira para investidores estrangeiros.
\end{itemize}

\subsubsection{Recomendação
Técnica}

Altamente indicado para residências de isolamento tático ``low profile''
ou projetos de pecuária/\allowbreak rizicultura. Exige planejamento de construção
contra cheias e autonomia de recursos. O GSS de 6.9 reflete o poder do
isolamento tático equilibrado pela vulnerabilidade hídrica. Referencia
atual de score: GSS 6.9 em 2026-03-05.

\subsubsection{}

\begin{itemize}
\tightlist
\item
  \begin{enumerate}
  \def\labelenumi{\arabic{enumi})}
  \tightlist
  \item
    Dados censitários validados (INE\footnote{Instituto Nacional de Estadística (Instituto Nacional de Estatística), órgão oficial de dados demográficos e censitários.} 2022).
  \end{enumerate}
\item
  \begin{enumerate}
  \def\labelenumi{\arabic{enumi})}
  \setcounter{enumi}{1}
  \tightlist
  \item
    Relatórios de engenharia do muro de defesa costeira (MOPC).
  \end{enumerate}
\item
  \begin{enumerate}
  \def\labelenumi{\arabic{enumi})}
  \setcounter{enumi}{2}
  \tightlist
  \item
    Mapa de solos e bacias hidrográficas de Ñeembucú (MAG).
  \end{enumerate}
\item
  \begin{enumerate}
  \def\labelenumi{\arabic{enumi})}
  \setcounter{enumi}{3}
  \tightlist
  \item
    Registro de segurança pública e comércio de fronteira.
  \end{enumerate}
\end{itemize}}
\begin{table}[ht!]
\small \begin{tabular}{p{3.0cm}p{0.8cm}p{6.0cm}} \toprule
\textbf{Critério} & \textbf{Nota} & \textbf{Justificativa} \\ \midrule
A - Ameaças Estratégicas & 7.5 & Isolamento geográfico estratégico; longe de alvos militares primários e eixos de fallout \\
B - Risco Social & 7.5 & População pequena e gerenciável; baixa densidade demográfica total \\
C - Riscos Naturais & 5.0 & Risco crítico de inundação; nota penalizada pela dependência de muros de defesa \\
D - Autossuficiência & 7.0 & Excelentes recursos de proteína animal e água; hub de suprimentos importados \\
E - Institucional & 6.5 & Comunidade resiliente e segura; forte presença institucional de fronteira \\
\midrule \textbf{GSS FINAL} & \textbf{6.9} & \textbf{Moderadamente Seguro (Recomendado para quem busca isolamento tático, desde que mitigados os riscos hídricos).} \\ \bottomrule \end{tabular}
\end{table}
\subsubsection{Vulnerabilidades e
Mitigação}\label{vuln-Alberdi-vulnerabilidades-e-mitigauxe7uxe3o}

\begin{itemize}
\tightlist
\item
  \textbf{Vulnerabilidade Hídrica Extrema:} Dependência vital do muro de
  contenção (Mitigação: residência construída com sistema de palafitas
  ou em zonas rurais mais elevadas fora do perímetro de inundação
  direta).
\item
  \textbf{Dependência Comercial:} Economia sensível à política cambial
  argentina (Mitigação: diversificação para pecuária ou agricultura de
  arroz).
\item
  \textbf{Isolamento de Redes:} Fragilidade elétrica (Mitigação:
  instalação de sistemas solares fotovoltaicos e geradores).
\end{itemize}
\subsubsection{Dossiê de Campo}\label{dossie-Alberdi-dossiuxea-de-campo}
\begin{description}[style=nextline,leftmargin=0.5cm,font=\bfseries]
\item[Coordenadas]  26°10'S 58°08'W.
\item[Topografia]  Relevo predominantemente plano, área de baixada (humedales) margeada pelo Rio Paraguai, situada estrategicamente em frente à cidade de Formosa (Argentina). Altitude \textasciitilde 60m. Área de \textasciitilde 222 km². Geologicamente estável, sem histórico de sismicidade relevante.
\item[Alvos Estratégicos]  Muro de Defesa Costeira (Infraestrutura crítica vital contra inundações); Porto fluvial de comércio transfronteiriço; Conexão logística via balsa com a Argentina. Localização estratégica como entreposto comercial, mas isolada dos eixos militares primários do Paraguai.
\item[Fallout]  Ventos predominantes do Norte (NE) e Leste. Baixíssimo risco de fallout de alvos nacionais; risco residual de incidentes industriais na vizinha Formosa (Arg).
\item[População]  7.536 habitantes (Censo 2022 Final). População com forte tradição comercial e resiliência hídrica.
\item[Densidade]  \textasciitilde 33,9 hab/\allowbreak km² (Densidade baixa, embora concentrada no núcleo urbano orla-rio). Oferece excelente isolamento social em relação aos grandes centros urbanos paraguaios.
\item[Vias de Acesso]  A nova Ruta PY19 (pavimentada) rompeu o isolamento histórico, conectando Alberdi a Villeta/\allowbreak Assunção ao norte e a Pilar ao sul. Acesso fluvial contínuo para a Argentina.
\item[Serviços]  Hospital Distrital de Alberdi. Infraestrutura comercial focada em eletrônicos e vestuário (polo de compras para argentinos). Serviços básicos funcionais, mas com dependência de centros maiores para alta complexidade.
\item[Custo de Vida]  Baixo; economia dolarizada pelo comércio fronteiriço, com acesso a produtos importados a preços competitivos.
\item[Preço da Terra]  US\$ 1.500-2.500/\allowbreak ha (campos para pecuária/\allowbreak arroz); US\$ 35-55/\allowbreak m² (Urbano).
\item[Hidrologia]  Risco Crítico de Inundação. Localidade historicamente vulnerável às cheias extraordinárias do Rio Paraguai. A segurança urbana depende 100\% da integridade do muro de defesa e do sistema de bombas de drenagem pluvial.
\item[Clima]  Subtropical Úmido. Verões intensos e úmidos; invernos suaves. Média de precipitação de 1.500 mm/\allowbreak ano.
\item[Energia]  Rede nacional (Itaipu/\allowbreak Yacyretá); vulnerável a interrupções em tempestades devido à extensão das linhas de transmissão rurais.
\item[Água]  Captação fluvial tratada e poços artesianos. Abundância hídrica superficial, mas com risco de contaminação em grandes cheias.
\item[Qualidade do Solo]  Solo argiloso e hidromórfico (campos baixos); excelente aptidão para pecuária de corte e rizicultura (arroz).
\item[Resources Locais]  Grande excedente de gado bovino e peixes fluviais. O hub comercial garante acesso a estoques de mercadorias importadas. Elevado potencial para autossuficiência de proteína animal.
\item[Segurança]  Cidade comercial muito pacífica e segura. Baixíssima criminalidade urbana violenta. Estabilidade social forjada pela cooperação comunitária em emergências hídricas. Forte presença da Armada Paraguaia e Polícia de Fronteira.
\item[Leis Local: Município consolidado com regime de comércio de fronteira. Restrição de Fronteira (Lei 2532/\allowbreak 05)]  Aplicável a terras rurais na faixa de 50km da fronteira argentina.
\end{description}
\paragraph{Solo (SoilGrids 2.0, média ponderada 0--30
cm)}

\begin{longtable}[]{@{}ll@{}}
\toprule\noalign{}
Parâmetro & Valor \\
\midrule\noalign{}
\endhead
\bottomrule\noalign{}
\endlastfoot
pH (H$_2$O) & 6.0 \\
Carbono orgânico (SOC) & 15.93 g/\allowbreak kg \\
Argila & 19.03 \% \\
Areia & 53.33 \% \\
Silte (calc.) & 27.6 \% \\
Densidade aparente & 1.34 g/\allowbreak cm³ \\
\textbf{Aptidão agrícola} & \textbf{Alta} \\
\end{longtable}

Fonte: ISRIC SoilGrids 2.0 via WCS (acesso em 2026-03-21). Coords:
-26.1667°, -58.1333°. Média ponderada camadas 0-5, 5-15, 15-30 cm.
\subsubsection{Indicadores Sociais}

\textbf{Fontes:} PNUD 2020, INE\footnote{Instituto Nacional de Estadística (Instituto Nacional de Estatística), órgão oficial de dados demográficos e censitários.} Censo 2022, INE\footnote{Instituto Nacional de Estadística (Instituto Nacional de Estatística), órgão oficial de dados demográficos e censitários.} EPHC 2023, Ministerio
Público 2024\\
\textbf{Nota:} Valores marcados como \emph{(dept.)} referem-se ao
departamento de Ñeembucú; valores \emph{(dist.)} são específicos deste
distrito. Dados marcados \emph{(est.)} são estimativas.

\begin{longtable}[]{@{}
  >{\raggedright\arraybackslash}p{(\linewidth - 4\tabcolsep) * \real{0.4231}}
  >{\raggedright\arraybackslash}p{(\linewidth - 4\tabcolsep) * \real{0.2692}}
  >{\raggedright\arraybackslash}p{(\linewidth - 4\tabcolsep) * \real{0.3077}}@{}}
\toprule\noalign{}
\begin{minipage}[b]{\linewidth}\raggedright
Indicador
\end{minipage} & \begin{minipage}[b]{\linewidth}\raggedright
Valor
\end{minipage} & \begin{minipage}[b]{\linewidth}\raggedright
Âmbito
\end{minipage} \\
\midrule\noalign{}
\endhead
\bottomrule\noalign{}
\endlastfoot
IDH (2020) & 0,710 & dept. \\
Ranking IDH nacional & 8/\allowbreak 18 & dept. \\
Esperança de vida & 73,0 anos & dept. \\
Escolaridade média & 9.2 anos & dist. \\
RNB per capita (USD PPA) & 8.800 & dept. \\
Pobreza monetária (\%) & 28,8\% & dept. \\
Pobreza extrema (\%) & 6,5\% & dept. \\
Índice de Gini & 0,445 & dept. \\
Acesso a água potável (\%) & 79,1\% & dept. \\
Acesso a saneamento (\%) & 77,4\% & dept. \\
Taxa de homicídios (est., /\allowbreak 100k hab) & 0,2--0,4 (2019--2020, fonte INE
--- mais baixa do país) & dept. \\
Índice de segurança & alto & dept. \\
População (Censo 2022) & 7.536 habitantes & dist. \\
Idade mediana & N/\allowbreak D & dist. \\
Taxa de fecundidade & N/\allowbreak D & dist. \\
\end{longtable}
\subsubsection{Combustível}

\textbf{Referência:} PETROPAR /\allowbreak  postos locais (2024) \textbf{Tipo de
localidade:} Interior

\begin{longtable}[]{@{}lll@{}}
\toprule\noalign{}
Combustível & USD/\allowbreak litro & Gs/\allowbreak litro (aprox.) \\
\midrule\noalign{}
\endhead
\bottomrule\noalign{}
\endlastfoot
Gasolina 93 oct & 0.99 & 7,326 \\
Gasolina 97 oct (premium) & 1.09 & 8,066 \\
Diesel & 0.92 & 6,808 \\
\end{longtable}

\begin{quote}
Preços podem variar ±5\% conforme posto e sazonalidade. Chaco e interior
remoto apresentam maior variação.
\end{quote}

\subsubsection{Cobertura Celular}

\textbf{Fonte:} CONATEL PY /\allowbreak  operadoras (2024)

\begin{longtable}[]{@{}ll@{}}
\toprule\noalign{}
Parâmetro & Valor \\
\midrule\noalign{}
\endhead
\bottomrule\noalign{}
\endlastfoot
Cobertura 4G população (dept.) & 75\% \\
Cobertura 4G área rural & 50\% \\
Melhor operadora & Tigo \\
Qualidade rural & moderada \\
\end{longtable}

\begin{quote}
Para áreas rurais fora do núcleo urbano, recomenda-se chip Tigo como
principal e Personal como backup.
\end{quote}

\subsubsection{Internet}

\textbf{Fonte:} CONATEL /\allowbreak  Speedtest Ookla (2024)

\begin{longtable}[]{@{}ll@{}}
\toprule\noalign{}
Parâmetro & Valor \\
\midrule\noalign{}
\endhead
\bottomrule\noalign{}
\endlastfoot
Velocidade média download & 30 Mbps \\
Domicílios com internet (dept.) & 40\% \\
Tecnologia predominante & rádio \\
Opção rural & Starlink disponível (\textasciitilde USD 44/\allowbreak mês) \\
\end{longtable}

\subsubsection{Mercado Imobiliário e Terra
Rural}

\textbf{Fonte:} INDERT /\allowbreak  Clasificados.com.py (2024)

\begin{longtable}[]{@{}ll@{}}
\toprule\noalign{}
Tipo & Referência \\
\midrule\noalign{}
\endhead
\bottomrule\noalign{}
\endlastfoot
Terra agrícola alta prod. (USD/\allowbreak ha) & 3,000 \\
Imóvel urbano (USD/\allowbreak m²) & 500 \\
Aluguel 2 quartos (USD/\allowbreak mês) & 200 \\
\end{longtable}

\begin{quote}
Valores de referência departamental. Localidades menores podem ter
preços 20--40\% abaixo da capital departamental. \#\#\# Saúde
\end{quote}

\textbf{Fonte:} MSPBS /\allowbreak  IPS Paraguay (2024-2026), consolidação
departamental e proxy local

\begin{longtable}[]{@{}ll@{}}
\toprule\noalign{}
Serviço & Disponibilidade \\
\midrule\noalign{}
\endhead
\bottomrule\noalign{}
\endlastfoot
USF /\allowbreak  Posto de Saúde & sim \\
Hospital Regional & sim \\
IPS (seguro social) & não \\
Farmácia & sim \\
Distância ao hospital de referência & \textasciitilde95 km (Pilar) \\
\end{longtable}

\textbf{Principais estabelecimentos:} USF local; referencia hospitalar
em Pilar

\textbf{Observação para imigrantes:} Atendimento primario local ou em
raio curto; casos de maior complexidade seguem para Pilar. Cobertura
privada continua recomendada para especialidades e urgências de maior
complexidade.
\newpage
\secaoDiagnostico{Cerrito}{\mapaConteudo{-27.3333}{-57.6666}}{Cerrito é o refúgio ``anfíbio'' do sul paraguaio. Oferece um isolamento
geográfico e demográfico que transforma a localidade em uma zona
virtualmente invisível a crises globais, protegida pela geografia dos
grandes rios e pela própria elevação que dá nome ao distrito. Sua força
absoluta reside na abundância inesgotável de recursos hídricos e
pesqueiros, além da segurança social de uma comunidade integrada e
pacífica. O desafio principal é a logística terrestre, que exige do
ocupante autonomia de transporte fluvial e estoque tático de
suprimentos. É o local ideal para quem busca segurança através da
inacessibilidade e abundância natural.

\subsubsection{Por que sim}

\begin{itemize}
\tightlist
\item
  Máximo isolamento geográfico e demográfico (anonimato absoluto).
\item
  Abundância ilimitada de água e recursos alimentares (peixes/\allowbreak gado).
\item
  Ambiente social extremamente seguro sob vigilância naval.
\item
  Longe de qualquer alvo militar terrestre primário ou eixo de fallout.
\end{itemize}

\subsubsection{Por que não}

\begin{itemize}
\tightlist
\item
  Risco de isolamento físico total por terra em eventos climáticos
  extremos.
\item
  Dependência de logística fluvial para suprimentos técnicos complexos.
\item
  Restrições da Lei de Fronteira para investidores estrangeiros.
\end{itemize}

\subsubsection{Recomendação
Técnica}

Altamente indicado para residências de isolamento tático que priorizem a
invisibilidade e possuam infraestrutura própria de transporte fluvial.
Requer autonomia energética e hídrica individual. O GSS de 7.1 reflete a
solidez do isolamento geográfico e a riqueza de recursos naturais da
localidade. Referencia atual de score: GSS 7.1 em 2026-03-05.

\subsubsection{}

\begin{itemize}
\tightlist
\item
  \begin{enumerate}
  \def\labelenumi{\arabic{enumi})}
  \tightlist
  \item
    Dados censitários validados (INE\footnote{Instituto Nacional de Estadística (Instituto Nacional de Estatística), órgão oficial de dados demográficos e censitários.} 2022).
  \end{enumerate}
\item
  \begin{enumerate}
  \def\labelenumi{\arabic{enumi})}
  \setcounter{enumi}{1}
  \tightlist
  \item
    Relatórios de monitoramento fluvial do Rio Paraná (EBY/\allowbreak ANNP).
  \end{enumerate}
\item
  \begin{enumerate}
  \def\labelenumi{\arabic{enumi})}
  \setcounter{enumi}{2}
  \tightlist
  \item
    Mapa de solos e bacias hídricas de Ñeembucú (Portal Geoestadístico).
  \end{enumerate}
\item
  \begin{enumerate}
  \def\labelenumi{\arabic{enumi})}
  \setcounter{enumi}{3}
  \tightlist
  \item
    Registro de segurança pública e governança municipal.
  \end{enumerate}
\end{itemize}}
\begin{table}[ht!]
\small \begin{tabular}{p{3.0cm}p{0.8cm}p{6.0cm}} \toprule
\textbf{Critério} & \textbf{Nota} & \textbf{Justificativa} \\ \midrule
A - Ameaças Estratégicas & 7.5 & Posição tática protegida por colina; margem do rio e isolamento geográfico superior; baixa visibilidade de alvos \\
B - Risco Social & 8.5 & Densidade populacional extremamente baixa; isolamento social de alto nível e privacidade absoluta \\
C - Riscos Naturais & 6.0 & Risco hidrológico elevado; nota penalizada pela vulnerabilidade a inundações sazonais severas \\
D - Autossuficiência & 7.5 & Excelentes recursos pesqueiros e de proteína animal; água fluvial ilimitada \\
E - Institucional & 6.0 & Ambiente social seguro e estável, mas com suporte institucional limitado regionalmente \\
\midrule \textbf{GSS FINAL} & \textbf{7.1} & \textbf{Seguro (Altamente recomendado para quem busca isolamento tático fluvial e anonimato em ambiente natural seguro).} \\ \bottomrule \end{tabular}
\end{table}
\subsubsection{Vulnerabilidades e
Mitigação}\label{vuln-Cerrito-vulnerabilidades-e-mitigauxe7uxe3o}

\begin{itemize}
\tightlist
\item
  \textbf{Isolamento Terrestre Recorrente:} Risco de ficar sem acesso
  por terra em grandes cheias (Mitigação: posse de embarcação a motor
  como via de transporte primária; manutenção de estoque de suprimentos
  para 90 dias).
\item
  \textbf{Fragilidade de Suporte Médico:} Distância crítica de hospitais
  terrestres (Mitigação: manutenção de estoque médico básico e kit
  trauma local; aproveitar proximidade fluvial com Ayolas).
\item
  \textbf{Dependência de Redes:} Vulnerabilidade elétrica rural
  (Mitigação: instalação de sistemas solares fotovoltaicos).
\end{itemize}
\subsubsection{Dossiê de Campo}\label{dossie-Cerrito-dossiuxea-de-campo}
\begin{description}[style=nextline,leftmargin=0.5cm,font=\bfseries]
\item[Coordenadas]  27°20'S 57°40'W.
\item[Topografia]  Relevo plano com a presença marcante de uma pequena colina ("Cerrito") que oferece visão tática sobre o Rio Paraná. Margeado ao sul pelo Rio Paraná (fronteira com Argentina). Altitude \textasciitilde 65m. Área massiva de \textasciitilde 682 km². Geologicamente estável.
\item[Alvos Estratégicos]  Polo de turismo de pesca internacional; Porto fluvial secundário; Conexão tática com a ilha argentina de Apipé. Ausência total de alvos militares de grande porte, oferecendo excelente invisibilidade estratégica e proteção por isolamento geográfico deliberado no extremo sul.
\item[Fallout]  Ventos predominantes do Norte (NE) e Leste. Baixíssimo risco de fallout de alvos continentais primários devido ao isolamento profundo.
\item[População]  4.340 habitantes (Censo 2022 Final). População com forte ligação com a atividade pesqueira e pecuária.
\item[Densidade]  \textasciitilde 6,4 hab/\allowbreak km² (Densidade extremamente baixa, garantindo privacidade absoluta e ausência de riscos sociais de massa urbana).
\item[Vias de Acesso]  Acesso via ramais de terra a partir da Ruta PY04. Localização remota que favorece o anonimato estratégico, mas sujeita a períodos de isolamento físico devido à precariedade das vias de terra durante cheias extraordinárias. Acesso fluvial contínuo e privilegiado.
\item[Serviços]  Unidade de Saúde da Família (USF) local. Dependência de Ayolas (via rio) ou Pilar (\textasciitilde 100 km) para atendimentos de alta complexidade. Infraestrutura comercial focada no turismo de pesca.
\item[Custo de Vida]  Baixo; economia baseada na pesca, pecuária extensiva e ecoturismo.
\item[Preço da Terra]  US\$ 2.000-4.000/\allowbreak ha (áreas com costa de rio); US\$ 800-1.200/\allowbreak ha (campos baixos).
\item[Hidrologia]  Risco de Inundação Moderado. Embora o núcleo urbano esteja em cota elevada (Cerrito), vastas áreas rurais e acessos terrestres são vulneráveis às cheias do Rio Paraná. Abundância hídrica superficial inesgotável.
\item[Clima]  Subtropical Úmido. Verões intensos; invernos amenos com alta umidade e ventos fluviais.
\item[Energia]  Rede nacional (Itaipu/\allowbreak Yacyretá); instável em áreas rurais remotas.
\item[Água]  Abundância hídrica absoluta via Rio Paraná e poços artesianos produtivos. Água superficial exige tratamento.
\item[Qualidade do Solo]  Solo arenoso e aluvial; excelente para pastagens naturais e pecuária de corte.
\item[Resources Locais]  Grande produção de recursos pesqueiros e gado bovino. Elevadíssimo potencial para autossuficiência de proteína animal para a pequena população local.
\item[Segurança]  Localidade excepcionalmente pacífica e segura. Criminalidade urbana comum praticamente inexistente. Estabilidade institucional básica sustentada pela presença de destacamentos da Armada Paraguaia para controle fluvial. Coesão social baseada na cultura da pesca.
\item[Leis Local: Município consolidado. Restrição de Fronteira (Lei 2532/\allowbreak 05)]  Aplicável a terras rurais na faixa de 50km da fronteira argentina (limite fluvial).
\end{description}
\paragraph{Solo (SoilGrids 2.0, média ponderada 0--30
cm)}

\begin{longtable}[]{@{}ll@{}}
\toprule\noalign{}
Parâmetro & Valor \\
\midrule\noalign{}
\endhead
\bottomrule\noalign{}
\endlastfoot
pH (H$_2$O) & 5.7 \\
Carbono orgânico (SOC) & 18.4 g/\allowbreak kg \\
Argila & 15.05 \% \\
Areia & 68.82 \% \\
Silte (calc.) & 16.1 \% \\
Densidade aparente & 1.35 g/\allowbreak cm³ \\
\textbf{Aptidão agrícola} & \textbf{Média} \\
\end{longtable}

Fonte: ISRIC SoilGrids 2.0 via WCS (acesso em 2026-03-21). Coords:
-27.3333°, -57.6667°. Média ponderada camadas 0-5, 5-15, 15-30 cm.
\subsubsection{Indicadores Sociais}

\textbf{Fontes:} PNUD 2020, INE\footnote{Instituto Nacional de Estadística (Instituto Nacional de Estatística), órgão oficial de dados demográficos e censitários.} Censo 2022, INE\footnote{Instituto Nacional de Estadística (Instituto Nacional de Estatística), órgão oficial de dados demográficos e censitários.} EPHC 2023, Ministerio
Público 2024\\
\textbf{Nota:} Valores marcados como \emph{(dept.)} referem-se ao
departamento de Ñeembucú; valores \emph{(dist.)} são específicos deste
distrito. Dados marcados \emph{(est.)} são estimativas.

\begin{longtable}[]{@{}
  >{\raggedright\arraybackslash}p{(\linewidth - 4\tabcolsep) * \real{0.4231}}
  >{\raggedright\arraybackslash}p{(\linewidth - 4\tabcolsep) * \real{0.2692}}
  >{\raggedright\arraybackslash}p{(\linewidth - 4\tabcolsep) * \real{0.3077}}@{}}
\toprule\noalign{}
\begin{minipage}[b]{\linewidth}\raggedright
Indicador
\end{minipage} & \begin{minipage}[b]{\linewidth}\raggedright
Valor
\end{minipage} & \begin{minipage}[b]{\linewidth}\raggedright
Âmbito
\end{minipage} \\
\midrule\noalign{}
\endhead
\bottomrule\noalign{}
\endlastfoot
IDH (2020) & 0,710 & dept. \\
Ranking IDH nacional & 8/\allowbreak 18 & dept. \\
Esperança de vida & 73,0 anos & dept. \\
Escolaridade média & 6.7 anos & dist. \\
RNB per capita (USD PPA) & 8.800 & dept. \\
Pobreza monetária (\%) & 28,8\% & dept. \\
Pobreza extrema (\%) & 6,5\% & dept. \\
Índice de Gini & 0,445 & dept. \\
Acesso a água potável (\%) & 79,1\% & dept. \\
Acesso a saneamento (\%) & 77,4\% & dept. \\
Taxa de homicídios (est., /\allowbreak 100k hab) & 0,2--0,4 (2019--2020, fonte INE
--- mais baixa do país) & dept. \\
Índice de segurança & alto & dept. \\
População (Censo 2022) & 4.340 habitantes & dist. \\
Idade mediana & N/\allowbreak D & dist. \\
Taxa de fecundidade & N/\allowbreak D & dist. \\
\end{longtable}
\subsubsection{Combustível}

\textbf{Referência:} PETROPAR /\allowbreak  postos locais (2024) \textbf{Tipo de
localidade:} Interior

\begin{longtable}[]{@{}lll@{}}
\toprule\noalign{}
Combustível & USD/\allowbreak litro & Gs/\allowbreak litro (aprox.) \\
\midrule\noalign{}
\endhead
\bottomrule\noalign{}
\endlastfoot
Gasolina 93 oct & 0.99 & 7,326 \\
Gasolina 97 oct (premium) & 1.09 & 8,066 \\
Diesel & 0.92 & 6,808 \\
\end{longtable}

\begin{quote}
Preços podem variar ±5\% conforme posto e sazonalidade. Chaco e interior
remoto apresentam maior variação.
\end{quote}

\subsubsection{Cobertura Celular}

\textbf{Fonte:} CONATEL PY /\allowbreak  operadoras (2024)

\begin{longtable}[]{@{}ll@{}}
\toprule\noalign{}
Parâmetro & Valor \\
\midrule\noalign{}
\endhead
\bottomrule\noalign{}
\endlastfoot
Cobertura 4G população (dept.) & 75\% \\
Cobertura 4G área rural & 50\% \\
Melhor operadora & Tigo \\
Qualidade rural & moderada \\
\end{longtable}

\begin{quote}
Para áreas rurais fora do núcleo urbano, recomenda-se chip Tigo como
principal e Personal como backup.
\end{quote}

\subsubsection{Internet}

\textbf{Fonte:} CONATEL /\allowbreak  Speedtest Ookla (2024)

\begin{longtable}[]{@{}ll@{}}
\toprule\noalign{}
Parâmetro & Valor \\
\midrule\noalign{}
\endhead
\bottomrule\noalign{}
\endlastfoot
Velocidade média download & 30 Mbps \\
Domicílios com internet (dept.) & 40\% \\
Tecnologia predominante & rádio \\
Opção rural & Starlink disponível (\textasciitilde USD 44/\allowbreak mês) \\
\end{longtable}

\subsubsection{Mercado Imobiliário e Terra
Rural}

\textbf{Fonte:} INDERT /\allowbreak  Clasificados.com.py (2024)

\begin{longtable}[]{@{}ll@{}}
\toprule\noalign{}
Tipo & Referência \\
\midrule\noalign{}
\endhead
\bottomrule\noalign{}
\endlastfoot
Terra agrícola alta prod. (USD/\allowbreak ha) & 3,000 \\
Imóvel urbano (USD/\allowbreak m²) & 500 \\
Aluguel 2 quartos (USD/\allowbreak mês) & 200 \\
\end{longtable}

\begin{quote}
Valores de referência departamental. Localidades menores podem ter
preços 20--40\% abaixo da capital departamental. \#\#\# Saúde
\end{quote}

\textbf{Fonte:} MSPBS /\allowbreak  IPS Paraguay (2024-2026), consolidação
departamental e proxy local

\begin{longtable}[]{@{}ll@{}}
\toprule\noalign{}
Serviço & Disponibilidade \\
\midrule\noalign{}
\endhead
\bottomrule\noalign{}
\endlastfoot
USF /\allowbreak  Posto de Saúde & sim \\
Hospital Regional & não \\
IPS (seguro social) & não \\
Farmácia & sim \\
Distância ao hospital de referência & \textasciitilde106 km (Pilar) \\
\end{longtable}

\textbf{Principais estabelecimentos:} USF local; referencia hospitalar
em Pilar

\textbf{Observação para imigrantes:} Atendimento primario local ou em
raio curto; casos de maior complexidade seguem para Pilar. Cobertura
privada continua recomendada para especialidades e urgências de maior
complexidade.
\newpage
\secaoDiagnostico{Desmochados}{\mapaConteudo{-27.0333}{-58.0333}}{Desmochados é o refúgio do ``vazio absoluto'' no sul paraguaio. Oferece
um nível de isolamento demográfico que torna o distrito virtualmente
invisível a crises externas e conflitos de massa. Sua força absoluta
reside na abundância de proteína animal e na segurança social imbatível
de uma comunidade minúscula e resiliente. O desafio existencial é a
dinâmica hídrica dos esteros, que exige do ocupante autonomia total para
operar em isolamento físico sazonal. É o local ideal para quem busca
segurança através da inacessibilidade natural e deseja viver em um
ambiente de paz social imperturbável.

\subsubsection{Por que sim}

\begin{itemize}
\tightlist
\item
  Máximo isolamento geográfico e demográfico (anonimato absoluto).
\item
  Ambiente social mais seguro do departamento pelo distanciamento
  físico.
\item
  Abundância massiva de proteína animal (gado).
\item
  Longe de qualquer alvo militar ou eixo de fallout tático.
\end{itemize}

\subsubsection{Por que não}

\begin{itemize}
\tightlist
\item
  Risco de isolamento físico total por terra em cheias sazonais.
\item
  Inexistência de serviços de saúde local e dependência de Pilar.
\item
  Solo inadequado para agricultura diversificada mecanizada.
\end{itemize}

\subsubsection{Recomendação
Técnica}

Altamente indicado para estabelecimentos de refúgio tático que priorizem
o isolamento e a invisibilidade total. Requer autonomia energética e
logística de suprimentos radical. O GSS de 7.5 reflete a superioridade
do isolamento protegida pela barreira natural dos humedales. Referencia
atual de score: GSS 7.5 em 2026-03-05.

\subsubsection{}

\begin{itemize}
\tightlist
\item
  \begin{enumerate}
  \def\labelenumi{\arabic{enumi})}
  \tightlist
  \item
    Dados censitários validados (INE\footnote{Instituto Nacional de Estadística (Instituto Nacional de Estatística), órgão oficial de dados demográficos e censitários.} 2022).
  \end{enumerate}
\item
  \begin{enumerate}
  \def\labelenumi{\arabic{enumi})}
  \setcounter{enumi}{1}
  \tightlist
  \item
    Relatórios de inundações sazonais da SEN\footnote{Secretaría de Emergencia Nacional (Secretaria de Emergência Nacional), órgão governamental de resposta a desastres.} em Ñeembucú.
  \end{enumerate}
\item
  \begin{enumerate}
  \def\labelenumi{\arabic{enumi})}
  \setcounter{enumi}{2}
  \tightlist
  \item
    Mapa de solos e bacias hídricas de Ñeembucú (Portal Geoestadístico).
  \end{enumerate}
\item
  \begin{enumerate}
  \def\labelenumi{\arabic{enumi})}
  \setcounter{enumi}{3}
  \tightlist
  \item
    Registro de segurança pública departamental.
  \end{enumerate}
\end{itemize}}
\begin{table}[ht!]
\small \begin{tabular}{p{3.0cm}p{0.8cm}p{6.0cm}} \toprule
\textbf{Critério} & \textbf{Nota} & \textbf{Justificativa} \\ \midrule
A - Ameaças Estratégicas & 8.0 & Isolamento tático profundo; baixíssima visibilidade estratégica e ausência de alvos \\
B - Risco Social & 8.5 & Densidade populacional próxima de zero; segurança contra riscos de massa absoluta \\
C - Riscos Naturais & 6.0 & Risco hidrológico elevado; nota penalizada pela vulnerabilidade a inundações sazonais \\
D - Autossuficiência & 7.5 & Excelentes recursos de proteína animal e água; limitado pela inacessibilidade comercial \\
E - Institucional & 7.0 & Ambiente social estável e pacífico; governança comunitária funcional para o isolamento \\
\midrule \textbf{GSS FINAL} & \textbf{7.5} & \textbf{Seguro (Altamente recomendado para quem busca isolamento tático radical e anonimato absoluto, aceitando os desafios do isolamento hídrico).} \\ \bottomrule \end{tabular}
\end{table}
\subsubsection{Vulnerabilidades e
Mitigação}\label{vuln-Desmochados-vulnerabilidades-e-mitigauxe7uxe3o}

\begin{itemize}
\tightlist
\item
  \textbf{Isolamento Geográfico Sazonal:} Risco de interrupção de
  acessos terrestres (Mitigação: posse de veículo 4x4 robusto ou trator;
  manutenção de estoque de suprimentos para 90 dias).
\item
  \textbf{Fragilidade das Redes:} Dependência de energia externa em área
  remota (Mitigação: instalação de sistemas solares fotovoltaicos
  potentes off-grid).
\item
  \textbf{Ausência de Saúde Local:} Distância crítica de hospitais
  (Mitigação: manutenção de estoque médico avançado e autonomia médica
  básica).
\end{itemize}
\subsubsection{Dossiê de Campo}\label{dossie-Desmochados-dossiuxea-de-campo}
\begin{description}[style=nextline,leftmargin=0.5cm,font=\bfseries]
\item[Coordenadas]  27°02'S 58°02'W.
\item[Topografia]  Relevo predominantemente plano, característico do Bajo Ñeembucú, com vasta presença de humedales e esteros. Altitude \textasciitilde 55m. Área de \textasciitilde 280 km². Geologicamente estável.
\item[Alvos Estratégicos]  Áreas de produção pecuária tradicional; Local de isolamento geográfico deliberado no extremo sul. Ausência total de alvos militares ou infraestrutura industrial massiva, oferecendo o nível máximo de invisibilidade estratégica e proteção por anonimato rural.
\item[Fallout]  Ventos predominantes do Norte (NE) e Leste. Baixíssimo risco de fallout de alvos continentais primários devido ao isolamento tático profundo.
\item[População]  1.171 habitantes (Censo 2022 Final). Uma das menores populações municipais do país, composta por comunidades rurais resilientes.
\item[Densidade]  \textasciitilde 4,2 hab/\allowbreak km² (Densidade extremamente baixa, garantindo privacidade absoluta e ausência total de riscos sociais de massa urbana).
\item[Vias de Acesso]  Acesso via ramais de terra que conectam à Ruta PY04. Localização que favorece a discrição, funcionando como um enclave de tranquilidade absoluta fora do radar nacional. Sujeita a isolamento físico durante períodos de chuvas torrenciais que saturam as vias vicinais.
\item[Serviços]  Infraestrutura básica mínima (USF local). Dependência absoluta de Pilar (\textasciitilde 45 km) para serviços de saúde, logística e suporte administrativo.
\item[Custo de Vida]  Muito baixo; economia baseada na pecuária extensiva e autoconsumo.
\item[Preço da Terra]  US\$ 700-1.100/\allowbreak ha (campos baixos para pecuária). Valores de entrada entre os menores do Paraguai oriental.
\item[Hidrologia]  Risco de Inundação Crítico. Situado em zona de baixada. Vulnerável a cheias sazonais extraordinárias que podem isolar a localidade por terra por semanas. Solo hidromórfico com baixíssima capacidade de infiltração.
\item[Clima]  Subtropical Úmido. Verões intensos; invernos amenos com alta umidade e nevoeiros frequentes.
\item[Energia]  Rede nacional (Itaipu/\allowbreak Yacyretá); instável em áreas rurais isoladas.
\item[Água]  Abundância hídrica superficial (esteros) e disponibilidade de poços artesianos. Águas superficiais exigem tratamento.
\item[Qualidade do Solo]  Solo hidromórfico e arenoso; aptidão exclusiva para pastagens naturais e gado de corte. Agricultura mecanizada inviável sem grandes investimentos.
\item[Resources Locais]  Grande excedente de gado bovino per capita. Elevadíssimo potencial para autossuficiência de proteína animal para a ínfima população local.
\item[Segurança]  Localidade excepcionalmente pacífica e segura. Criminalidade urbana inexistente. Estabilidade social máxima sustentada pela resiliência da população ao isolamento. Localidade "invisível" para conflitos sociopolíticos nacionais.
\item[Leis Local: Município consolidado. Restrição de Fronteira (Lei 2532/\allowbreak 05)]  Aplicável a terras rurais na faixa de 50km da fronteira argentina.
\end{description}
\paragraph{Solo (SoilGrids 2.0, média ponderada 0--30
cm)}

\begin{longtable}[]{@{}ll@{}}
\toprule\noalign{}
Parâmetro & Valor \\
\midrule\noalign{}
\endhead
\bottomrule\noalign{}
\endlastfoot
pH (H$_2$O) & 6.05 \\
Carbono orgânico (SOC) & 15.33 g/\allowbreak kg \\
Argila & 16.82 \% \\
Areia & 67.4 \% \\
Silte (calc.) & 15.8 \% \\
Densidade aparente & 1.37 g/\allowbreak cm³ \\
\textbf{Aptidão agrícola} & \textbf{Alta} \\
\end{longtable}

Fonte: ISRIC SoilGrids 2.0 via WCS (acesso em 2026-03-21). Coords:
-27.0333°, -58.0333°. Média ponderada camadas 0-5, 5-15, 15-30 cm.

\begin{itemize}
\tightlist
\item
  \textbf{Homicidios/\allowbreak 100k:} \textasciitilde0,0 (Histórico de paz
  absoluta).
\end{itemize}

\subsection{10. Analise de Riscos}

\begin{itemize}
\tightlist
\item
  \textbf{Cenario curto prazo:} Manutenção da tranquilidade rural
  absoluta e isolamento físico cíclico.
\item
  \textbf{Cenario medio prazo:} Impacto de mudanças no regime de chuvas
  sobre a sustentabilidade da pecuária extensiva.
\end{itemize}

\subsection{11. Redacao Final}

\subsubsection{Diagnostico Integrado}

Desmochados é o refúgio do ``vazio absoluto'' no sul paraguaio. Oferece
um nível de isolamento demográfico que torna o distrito virtualmente
invisível a crises externas e conflitos de massa. Sua força absoluta
reside na abundância de proteína animal e na segurança social imbatível
de uma comunidade minúscula e resiliente. O desafio existencial é a
dinâmica hídrica dos esteros, que exige do ocupante autonomia total para
operar em isolamento físico sazonal. É o local ideal para quem busca
segurança através da inacessibilidade natural e deseja viver em um
ambiente de paz social imperturbável.

\subsubsection{Por que sim}

\begin{itemize}
\tightlist
\item
  Máximo isolamento geográfico e demográfico (anonimato absoluto).
\item
  Ambiente social mais seguro do departamento pelo distanciamento
  físico.
\item
  Abundância massiva de proteína animal (gado).
\item
  Longe de qualquer alvo militar ou eixo de fallout tático.
\end{itemize}

\subsubsection{Por que nao}

\begin{itemize}
\tightlist
\item
  Risco de isolamento físico total por terra em cheias sazonais.
\item
  Inexistência de serviços de saúde local e dependência de Pilar.
\item
  Solo inadequado para agricultura diversificada mecanizada.
\end{itemize}

\subsubsection{Recomendacao Tecnica}

Altamente indicado para estabelecimentos de refúgio tático que priorizem
o isolamento e a invisibilidade total. Requer autonomia energética e
logística de suprimentos radical. O GSS de 7.5 reflete a superioridade
do isolamento protegida pela barreira natural dos humedales. Referencia
atual de score: GSS 7.5 em 2026-03-05.

\subsection{13.}

\begin{itemize}
\tightlist
\item
  \begin{enumerate}
  \def\labelenumi{\arabic{enumi})}
  \tightlist
  \item
    Dados censitários validados (INE\footnote{Instituto Nacional de Estadística (Instituto Nacional de Estatística), órgão oficial de dados demográficos e censitários.} 2022).
  \end{enumerate}
\item
  \begin{enumerate}
  \def\labelenumi{\arabic{enumi})}
  \setcounter{enumi}{1}
  \tightlist
  \item
    Relatórios de inundações sazonais da SEN\footnote{Secretaría de Emergencia Nacional (Secretaria de Emergência Nacional), órgão governamental de resposta a desastres.} em Ñeembucú.
  \end{enumerate}
\item
  \begin{enumerate}
  \def\labelenumi{\arabic{enumi})}
  \setcounter{enumi}{2}
  \tightlist
  \item
    Mapa de solos e bacias hídricas de Ñeembucú (Portal Geoestadístico).
  \end{enumerate}
\item
  \begin{enumerate}
  \def\labelenumi{\arabic{enumi})}
  \setcounter{enumi}{3}
  \tightlist
  \item
    Registro de segurança pública departamental.
  \end{enumerate}
\end{itemize}

\subsubsection{Combustível}

\textbf{Referência:} PETROPAR /\allowbreak  postos locais (2024) \textbf{Tipo de
localidade:} Interior

\begin{longtable}[]{@{}lll@{}}
\toprule\noalign{}
Combustível & USD/\allowbreak litro & Gs/\allowbreak litro (aprox.) \\
\midrule\noalign{}
\endhead
\bottomrule\noalign{}
\endlastfoot
Gasolina 93 oct & 0.99 & 7,326 \\
Gasolina 97 oct (premium) & 1.09 & 8,066 \\
Diesel & 0.92 & 6,808 \\
\end{longtable}

\begin{quote}
Preços podem variar ±5\% conforme posto e sazonalidade. Chaco e interior
remoto apresentam maior variação.
\end{quote}

\subsubsection{Cobertura Celular}

\textbf{Fonte:} CONATEL PY /\allowbreak  operadoras (2024)

\begin{longtable}[]{@{}ll@{}}
\toprule\noalign{}
Parâmetro & Valor \\
\midrule\noalign{}
\endhead
\bottomrule\noalign{}
\endlastfoot
Cobertura 4G população (dept.) & 75\% \\
Cobertura 4G área rural & 50\% \\
Melhor operadora & Tigo \\
Qualidade rural & moderada \\
\end{longtable}

\begin{quote}
Para áreas rurais fora do núcleo urbano, recomenda-se chip Tigo como
principal e Personal como backup.
\end{quote}

\subsubsection{Internet}

\textbf{Fonte:} CONATEL /\allowbreak  Speedtest Ookla (2024)

\begin{longtable}[]{@{}ll@{}}
\toprule\noalign{}
Parâmetro & Valor \\
\midrule\noalign{}
\endhead
\bottomrule\noalign{}
\endlastfoot
Velocidade média download & 30 Mbps \\
Domicílios com internet (dept.) & 40\% \\
Tecnologia predominante & rádio \\
Opção rural & Starlink disponível (\textasciitilde USD 44/\allowbreak mês) \\
\end{longtable}

\subsubsection{Mercado Imobiliário e Terra
Rural}

\textbf{Fonte:} INDERT /\allowbreak  Clasificados.com.py (2024)

\begin{longtable}[]{@{}ll@{}}
\toprule\noalign{}
Tipo & Referência \\
\midrule\noalign{}
\endhead
\bottomrule\noalign{}
\endlastfoot
Terra agrícola alta prod. (USD/\allowbreak ha) & 3,000 \\
Imóvel urbano (USD/\allowbreak m²) & 500 \\
Aluguel 2 quartos (USD/\allowbreak mês) & 200 \\
\end{longtable}

\begin{quote}
Valores de referência departamental. Localidades menores podem ter
preços 20--40\% abaixo da capital departamental. \#\#\# Saúde
\end{quote}

\textbf{Fonte:} MSPBS /\allowbreak  IPS Paraguay (2024-2026), consolidação
departamental e proxy local

\begin{longtable}[]{@{}ll@{}}
\toprule\noalign{}
Serviço & Disponibilidade \\
\midrule\noalign{}
\endhead
\bottomrule\noalign{}
\endlastfoot
USF /\allowbreak  Posto de Saúde & sim \\
Hospital Regional & não \\
IPS (seguro social) & não \\
Farmácia & sim \\
Distância ao hospital de referência & \textasciitilde45 km (Pilar) \\
\end{longtable}

\textbf{Principais estabelecimentos:} USF local; referencia hospitalar
em Pilar

\textbf{Observação para imigrantes:} Atendimento primario local ou em
raio curto; casos de maior complexidade seguem para Pilar. Cobertura
privada continua recomendada para especialidades e urgências de maior
complexidade.
\subsubsection{Indicadores Sociais}

\textbf{Fontes:} PNUD 2020, INE\footnote{Instituto Nacional de Estadística (Instituto Nacional de Estatística), órgão oficial de dados demográficos e censitários.} Censo 2022, INE\footnote{Instituto Nacional de Estadística (Instituto Nacional de Estatística), órgão oficial de dados demográficos e censitários.} EPHC 2023, Ministerio
Público 2024\\
\textbf{Nota:} Valores marcados como \emph{(dept.)} referem-se ao
departamento de Ñeembucú; valores \emph{(dist.)} são específicos deste
distrito. Dados marcados \emph{(est.)} são estimativas.

\begin{longtable}[]{@{}
  >{\raggedright\arraybackslash}p{(\linewidth - 4\tabcolsep) * \real{0.4231}}
  >{\raggedright\arraybackslash}p{(\linewidth - 4\tabcolsep) * \real{0.2692}}
  >{\raggedright\arraybackslash}p{(\linewidth - 4\tabcolsep) * \real{0.3077}}@{}}
\toprule\noalign{}
\begin{minipage}[b]{\linewidth}\raggedright
Indicador
\end{minipage} & \begin{minipage}[b]{\linewidth}\raggedright
Valor
\end{minipage} & \begin{minipage}[b]{\linewidth}\raggedright
Âmbito
\end{minipage} \\
\midrule\noalign{}
\endhead
\bottomrule\noalign{}
\endlastfoot
IDH (2020) & 0,710 & dept. \\
Ranking IDH nacional & 8/\allowbreak 18 & dept. \\
Esperança de vida & 73,0 anos & dept. \\
Escolaridade média & 7.3 anos & dist. \\
RNB per capita (USD PPA) & 8.800 & dept. \\
Pobreza monetária (\%) & 28,8\% & dept. \\
Pobreza extrema (\%) & 6,5\% & dept. \\
Índice de Gini & 0,445 & dept. \\
Acesso a água potável (\%) & 79,1\% & dept. \\
Acesso a saneamento (\%) & 77,4\% & dept. \\
Taxa de homicídios (est., /\allowbreak 100k hab) & 0,2--0,4 (2019--2020, fonte INE
--- mais baixa do país) & dept. \\
Índice de segurança & alto & dept. \\
População (Censo 2022) & 1.171 habitantes & dist. \\
Idade mediana & N/\allowbreak D & dist. \\
Taxa de fecundidade & N/\allowbreak D & dist. \\
\end{longtable}
\subsubsection{Combustível}

\textbf{Referência:} PETROPAR /\allowbreak  postos locais (2024) \textbf{Tipo de
localidade:} Interior

\begin{longtable}[]{@{}lll@{}}
\toprule\noalign{}
Combustível & USD/\allowbreak litro & Gs/\allowbreak litro (aprox.) \\
\midrule\noalign{}
\endhead
\bottomrule\noalign{}
\endlastfoot
Gasolina 93 oct & 0.99 & 7,326 \\
Gasolina 97 oct (premium) & 1.09 & 8,066 \\
Diesel & 0.92 & 6,808 \\
\end{longtable}

\begin{quote}
Preços podem variar ±5\% conforme posto e sazonalidade. Chaco e interior
remoto apresentam maior variação.
\end{quote}

\subsubsection{Cobertura Celular}

\textbf{Fonte:} CONATEL PY /\allowbreak  operadoras (2024)

\begin{longtable}[]{@{}ll@{}}
\toprule\noalign{}
Parâmetro & Valor \\
\midrule\noalign{}
\endhead
\bottomrule\noalign{}
\endlastfoot
Cobertura 4G população (dept.) & 75\% \\
Cobertura 4G área rural & 50\% \\
Melhor operadora & Tigo \\
Qualidade rural & moderada \\
\end{longtable}

\begin{quote}
Para áreas rurais fora do núcleo urbano, recomenda-se chip Tigo como
principal e Personal como backup.
\end{quote}

\subsubsection{Internet}

\textbf{Fonte:} CONATEL /\allowbreak  Speedtest Ookla (2024)

\begin{longtable}[]{@{}ll@{}}
\toprule\noalign{}
Parâmetro & Valor \\
\midrule\noalign{}
\endhead
\bottomrule\noalign{}
\endlastfoot
Velocidade média download & 30 Mbps \\
Domicílios com internet (dept.) & 40\% \\
Tecnologia predominante & rádio \\
Opção rural & Starlink disponível (\textasciitilde USD 44/\allowbreak mês) \\
\end{longtable}

\subsubsection{Mercado Imobiliário e Terra
Rural}

\textbf{Fonte:} INDERT /\allowbreak  Clasificados.com.py (2024)

\begin{longtable}[]{@{}ll@{}}
\toprule\noalign{}
Tipo & Referência \\
\midrule\noalign{}
\endhead
\bottomrule\noalign{}
\endlastfoot
Terra agrícola alta prod. (USD/\allowbreak ha) & 3,000 \\
Imóvel urbano (USD/\allowbreak m²) & 500 \\
Aluguel 2 quartos (USD/\allowbreak mês) & 200 \\
\end{longtable}

\begin{quote}
Valores de referência departamental. Localidades menores podem ter
preços 20--40\% abaixo da capital departamental. \#\#\# Saúde
\end{quote}

\textbf{Fonte:} MSPBS /\allowbreak  IPS Paraguay (2024-2026), consolidação
departamental e proxy local

\begin{longtable}[]{@{}ll@{}}
\toprule\noalign{}
Serviço & Disponibilidade \\
\midrule\noalign{}
\endhead
\bottomrule\noalign{}
\endlastfoot
USF /\allowbreak  Posto de Saúde & sim \\
Hospital Regional & não \\
IPS (seguro social) & não \\
Farmácia & sim \\
Distância ao hospital de referência & \textasciitilde45 km (Pilar) \\
\end{longtable}

\textbf{Principais estabelecimentos:} USF local; referencia hospitalar
em Pilar

\textbf{Observação para imigrantes:} Atendimento primario local ou em
raio curto; casos de maior complexidade seguem para Pilar. Cobertura
privada continua recomendada para especialidades e urgências de maior
complexidade.
\newpage
\secaoDiagnostico{General Diaz}{\mapaConteudo{-27.1666}{-58.4166}}{General José Eduvigis Díaz é o arquétipo do isolamento tático no sul
paraguaio. Oferece o máximo em privacidade e segurança social, sendo uma
localidade ``fora do radar'' das grandes crises demográficas e
militares. Sua força reside na abundância de recursos de proteína animal
e no anonimato geográfico. O custo desse isolamento é a vulnerabilidade
hídrica e a dependência de Pilar para serviços vitais. É o local ideal
para quem busca viver com autonomia total de suprimentos em um ambiente
de paz absoluta, desde que esteja preparado para períodos de isolamento
físico causados pela natureza dos humedales.

\subsubsection{Por que sim}

\begin{itemize}
\tightlist
\item
  Densidade populacional extremamente baixa (privacidade absoluta).
\item
  Ambiente social seguro, ordeiro e pacífico.
\item
  Abundância de proteína animal (gado) e águas superficiais.
\item
  Baixíssima visibilidade estratégica.
\end{itemize}

\subsubsection{Por que não}

\begin{itemize}
\tightlist
\item
  Risco crítico de inundações e intransitabilidade de acessos de terra.
\item
  Infraestrutura de saúde e comercial local mínima.
\item
  Restrições da Lei de Fronteira.
\end{itemize}

\subsubsection{Recomendação
Técnica}

Altamente indicado para estabelecimentos de autossuficiência pecuária
que busquem o máximo isolamento demográfico. Requer investimento em
autonomia logística, energética e de saúde. O GSS de 6.5 reflete a
segurança do isolamento em contraste com a fragilidade da infraestrutura
e riscos naturais. Referencia atual de score: GSS 6.5 em 2026-03-05.

\subsubsection{}

\begin{itemize}
\tightlist
\item
  \begin{enumerate}
  \def\labelenumi{\arabic{enumi})}
  \tightlist
  \item
    Dados censitários validados (INE\footnote{Instituto Nacional de Estadística (Instituto Nacional de Estatística), órgão oficial de dados demográficos e censitários.} 2022).
  \end{enumerate}
\item
  \begin{enumerate}
  \def\labelenumi{\arabic{enumi})}
  \setcounter{enumi}{1}
  \tightlist
  \item
    Relatórios de inundações sazonais da Secretaria de Emergência
    Nacional (SEN\footnote{Secretaría de Emergencia Nacional (Secretaria de Emergência Nacional), órgão governamental de resposta a desastres.}).
  \end{enumerate}
\item
  \begin{enumerate}
  \def\labelenumi{\arabic{enumi})}
  \setcounter{enumi}{2}
  \tightlist
  \item
    Mapa de solos e aptidão para pastagens de Ñeembucú (MAG).
  \end{enumerate}
\item
  \begin{enumerate}
  \def\labelenumi{\arabic{enumi})}
  \setcounter{enumi}{3}
  \tightlist
  \item
    Registro de segurança pública departamental.
  \end{enumerate}
\end{itemize}}
\begin{table}[ht!]
\small \begin{tabular}{p{3.0cm}p{0.8cm}p{6.0cm}} \toprule
\textbf{Critério} & \textbf{Nota} & \textbf{Justificativa} \\ \midrule
A - Ameaças Estratégicas & 7.5 & Localização remota e discreta; excelente invisibilidade tática e ausência de alvos estratégicos \\
B - Risco Social & 8.0 & Densidade demográfica extremamente baixa; isolamento social superior \\
C - Riscos Naturais & 4.5 & Risco crítico de inundação e isolamento pluvial; solo saturável \\
D - Autossuficiência & 6.5 & Bons recursos de proteína animal; limitado pela infraestrutura de transporte \\
E - Institucional & 5.5 & Ambiente social seguro e estável, mas com serviços institucionais mínimos \\
\midrule \textbf{GSS FINAL} & \textbf{6.5} & \textbf{Moderadamente Seguro (Ideal para perfis que busquem isolamento total e autossuficiência pecuária, aceitando os desafios logísticos).} \\ \bottomrule \end{tabular}
\end{table}
\subsubsection{Vulnerabilidades e
Mitigação}\label{vuln-General_Diaz-vulnerabilidades-e-mitigauxe7uxe3o}

\begin{itemize}
\tightlist
\item
  \textbf{Isolamento Físico:} Risco de ficar ``ilhada'' durante grandes
  inundações (Mitigação: posse de veículo 4x4 robusto ou trator;
  manutenção de estoque de suprimentos para 60 dias).
\item
  \textbf{Dependência Elétrica:} Fragilidade da rede em áreas remotas
  (Mitigação: instalação de sistemas solares fotovoltaicos potentes e
  geradores).
\item
  \textbf{Fragilidade Médica:} Distância de hospitais cirúrgicos
  (Mitigação: manutenção de estoque médico básico e treinamento avançado
  em primeiros socorros).
\end{itemize}
\subsubsection{Dossiê de Campo}\label{dossie-General_Diaz-dossiuxea-de-campo}
\begin{description}[style=nextline,leftmargin=0.5cm,font=\bfseries]
\item[Coordenadas]  27°10'S 58°25'W.
\item[Topografia]  Relevo predominantemente plano, caracterizado por extensas áreas de esteros (pantanais) e proximidade com o Rio Paraná ao sul. Altitude \textasciitilde 55m. Área de \textasciitilde 291 km². Geologicamente estável, com histórico de sismo leve (3.0 em 2017) sem danos estruturais.
\item[Alvos Estratégicos]  Local histórico de importância militar (Campo de Batalha de Estero Bellaco); Áreas de pecuária extensiva; Proximidade tática com o Rio Paraná e a fronteira com a Argentina. Ausência de alvos industriais, oferecendo excelente invisibilidade estratégica.
\item[Fallout]  Ventos predominantes do Norte (NE) e Leste. Baixíssimo risco de fallout de alvos continentais primários devido ao isolamento geográfico.
\item[População]  3.164 habitantes (Censo 2022 Final). Perfil demográfico de baixa densidade e forte tradição rural.
\item[Densidade]  \textasciitilde 10,8 hab/\allowbreak km² (Densidade extremamente baixa, ideal para isolamento social tático e privacidade absoluta).
\item[Vias de Acesso]  Acesso via ramais de terra a partir da Ruta PY04 (Pilar-San Ignacio). Localização remota que favorece o anonimato, mas sujeita ao isolamento físico em períodos de chuvas intensas.
\item[Serviços]  Infraestrutura básica mínima (Centro de Saúde local). Dependência total de Pilar (\textasciitilde 45 km) para atendimentos médicos de alta complexidade e logística avançada.
\item[Custo de Vida]  Muito baixo; economia baseada na pecuária e agricultura de subsistência.
\item[Preço da Terra]  US\$ 800-1.200/\allowbreak ha (campos baixos para pecuária). Valores de entrada muito baixos comparados à região metropolitana.
\item[Hidrologia]  Risco de Inundação e Isolamento. Localizado em zona de esteros, altamente vulnerável a cheias sazonais. O solo hidromórfico retém água por longos períodos, tornando as estradas de terra intransitáveis em eventos extremos.
\item[Clima]  Subtropical Úmido. Verões quentes; invernos frescos com alta umidade relativa.
\item[Energia]  Rede nacional (Itaipu/\allowbreak Yacyretá); instável e sujeita a cortes frequentes devido à exposição das linhas rurais.
\item[Água]  Abundante via poços artesianos e arroios locais; excelente qualidade hídrica superficial para pecuária.
\item[Qualidade do Solo]  Solo hidromórfico e arenoso; excelente para pastagens naturais e gado de corte. Baixa aptidão para agricultura mecanizada de grande escala sem drenagem técnica.
\item[Resources Locais]  Grande produção de gado bovino, queijo paraguaio e leite. Elevado potencial para autossuficiência calórica familiar baseada em proteína animal.
\item[Segurança]  Zona rural extremamente pacífica e segura. Baixíssima criminalidade urbana. Estabilidade social sustentada por laços familiares e comunitários tradicionais. Ausência de visibilidade para conflitos sociais.
\item[Leis Local: Município consolidado. Restrição de Fronteira (Lei 2532/\allowbreak 05)]  Aplicável a terras rurais na faixa de 50km da fronteira argentina.
\end{description}
\paragraph{Solo (SoilGrids 2.0, média ponderada 0--30
cm)}

\begin{longtable}[]{@{}ll@{}}
\toprule\noalign{}
Parâmetro & Valor \\
\midrule\noalign{}
\endhead
\bottomrule\noalign{}
\endlastfoot
pH (H$_2$O) & 6.0 \\
Carbono orgânico (SOC) & 14.18 g/\allowbreak kg \\
Argila & 17.26 \% \\
Areia & 61.12 \% \\
Silte (calc.) & 21.6 \% \\
Densidade aparente & 1.39 g/\allowbreak cm³ \\
\textbf{Aptidão agrícola} & \textbf{Média} \\
\end{longtable}

Fonte: ISRIC SoilGrids 2.0 via WCS (acesso em 2026-03-21). Coords:
-27.1667°, -58.4167°. Média ponderada camadas 0-5, 5-15, 15-30 cm.

\begin{itemize}
\tightlist
\item
  \textbf{Homicidios/\allowbreak 100k:} \textasciitilde4,0 (Muito abaixo da média
  nacional).
\end{itemize}

\subsection{10. Analise de Riscos}

\begin{itemize}
\tightlist
\item
  \textbf{Cenario curto prazo:} Manutenção da tranquilidade rural e
  estabilidade da pecuária extensiva.
\item
  \textbf{Cenario medio prazo:} Impacto das mudanças climáticas na
  frequência de isolamento por inundações.
\end{itemize}

\subsection{11. Redacao Final}

\subsubsection{Diagnostico Integrado}

General José Eduvigis Díaz é o arquétipo do isolamento tático no sul
paraguaio. Oferece o máximo em privacidade e segurança social, sendo uma
localidade ``fora do radar'' das grandes crises demográficas e
militares. Sua força reside na abundância de recursos de proteína animal
e no anonimato geográfico. O custo desse isolamento é a vulnerabilidade
hídrica e a dependência de Pilar para serviços vitais. É o local ideal
para quem busca viver com autonomia total de suprimentos em um ambiente
de paz absoluta, desde que esteja preparado para períodos de isolamento
físico causados pela natureza dos humedales.

\subsubsection{Por que sim}

\begin{itemize}
\tightlist
\item
  Densidade populacional extremamente baixa (privacidade absoluta).
\item
  Ambiente social seguro, ordeiro e pacífico.
\item
  Abundância de proteína animal (gado) e águas superficiais.
\item
  Baixíssima visibilidade estratégica.
\end{itemize}

\subsubsection{Por que nao}

\begin{itemize}
\tightlist
\item
  Risco crítico de inundações e intransitabilidade de acessos de terra.
\item
  Infraestrutura de saúde e comercial local mínima.
\item
  Restrições da Lei de Fronteira.
\end{itemize}

\subsubsection{Recomendacao Tecnica}

Altamente indicado para estabelecimentos de autossuficiência pecuária
que busquem o máximo isolamento demográfico. Requer investimento em
autonomia logística, energética e de saúde. O GSS de 6.5 reflete a
segurança do isolamento em contraste com a fragilidade da infraestrutura
e riscos naturais. Referencia atual de score: GSS 6.5 em 2026-03-05.

\subsection{13.}

\begin{itemize}
\tightlist
\item
  \begin{enumerate}
  \def\labelenumi{\arabic{enumi})}
  \tightlist
  \item
    Dados censitários validados (INE\footnote{Instituto Nacional de Estadística (Instituto Nacional de Estatística), órgão oficial de dados demográficos e censitários.} 2022).
  \end{enumerate}
\item
  \begin{enumerate}
  \def\labelenumi{\arabic{enumi})}
  \setcounter{enumi}{1}
  \tightlist
  \item
    Relatórios de inundações sazonais da Secretaria de Emergência
    Nacional (SEN\footnote{Secretaría de Emergencia Nacional (Secretaria de Emergência Nacional), órgão governamental de resposta a desastres.}).
  \end{enumerate}
\item
  \begin{enumerate}
  \def\labelenumi{\arabic{enumi})}
  \setcounter{enumi}{2}
  \tightlist
  \item
    Mapa de solos e aptidão para pastagens de Ñeembucú (MAG).
  \end{enumerate}
\item
  \begin{enumerate}
  \def\labelenumi{\arabic{enumi})}
  \setcounter{enumi}{3}
  \tightlist
  \item
    Registro de segurança pública departamental.
  \end{enumerate}
\end{itemize}

\subsubsection{Combustível}

\textbf{Referência:} PETROPAR /\allowbreak  postos locais (2024) \textbf{Tipo de
localidade:} Interior

\begin{longtable}[]{@{}lll@{}}
\toprule\noalign{}
Combustível & USD/\allowbreak litro & Gs/\allowbreak litro (aprox.) \\
\midrule\noalign{}
\endhead
\bottomrule\noalign{}
\endlastfoot
Gasolina 93 oct & 0.99 & 7,326 \\
Gasolina 97 oct (premium) & 1.09 & 8,066 \\
Diesel & 0.92 & 6,808 \\
\end{longtable}

\begin{quote}
Preços podem variar ±5\% conforme posto e sazonalidade. Chaco e interior
remoto apresentam maior variação.
\end{quote}

\subsubsection{Cobertura Celular}

\textbf{Fonte:} CONATEL PY /\allowbreak  operadoras (2024)

\begin{longtable}[]{@{}ll@{}}
\toprule\noalign{}
Parâmetro & Valor \\
\midrule\noalign{}
\endhead
\bottomrule\noalign{}
\endlastfoot
Cobertura 4G população (dept.) & 75\% \\
Cobertura 4G área rural & 50\% \\
Melhor operadora & Tigo \\
Qualidade rural & moderada \\
\end{longtable}

\begin{quote}
Para áreas rurais fora do núcleo urbano, recomenda-se chip Tigo como
principal e Personal como backup.
\end{quote}

\subsubsection{Internet}

\textbf{Fonte:} CONATEL /\allowbreak  Speedtest Ookla (2024)

\begin{longtable}[]{@{}ll@{}}
\toprule\noalign{}
Parâmetro & Valor \\
\midrule\noalign{}
\endhead
\bottomrule\noalign{}
\endlastfoot
Velocidade média download & 30 Mbps \\
Domicílios com internet (dept.) & 40\% \\
Tecnologia predominante & rádio \\
Opção rural & Starlink disponível (\textasciitilde USD 44/\allowbreak mês) \\
\end{longtable}

\subsubsection{Mercado Imobiliário e Terra
Rural}

\textbf{Fonte:} INDERT /\allowbreak  Clasificados.com.py (2024)

\begin{longtable}[]{@{}ll@{}}
\toprule\noalign{}
Tipo & Referência \\
\midrule\noalign{}
\endhead
\bottomrule\noalign{}
\endlastfoot
Terra agrícola alta prod. (USD/\allowbreak ha) & 3,000 \\
Imóvel urbano (USD/\allowbreak m²) & 500 \\
Aluguel 2 quartos (USD/\allowbreak mês) & 200 \\
\end{longtable}

\begin{quote}
Valores de referência departamental. Localidades menores podem ter
preços 20--40\% abaixo da capital departamental. \#\#\# Saúde
\end{quote}

\textbf{Fonte:} MSPBS /\allowbreak  IPS Paraguay (2024-2026), consolidação
departamental e proxy local

\begin{longtable}[]{@{}ll@{}}
\toprule\noalign{}
Serviço & Disponibilidade \\
\midrule\noalign{}
\endhead
\bottomrule\noalign{}
\endlastfoot
USF /\allowbreak  Posto de Saúde & sim \\
Hospital Regional & não \\
IPS (seguro social) & não \\
Farmácia & sim \\
Distância ao hospital de referência & \textasciitilde1 km (Pilar) \\
\end{longtable}

\textbf{Principais estabelecimentos:} USF local; referencia hospitalar
em Pilar

\textbf{Observação para imigrantes:} Atendimento primario local ou em
raio curto; casos de maior complexidade seguem para Pilar. Cobertura
privada continua recomendada para especialidades e urgências de maior
complexidade.
\subsubsection{Indicadores Sociais}

\textbf{Fontes:} PNUD 2020, INE\footnote{Instituto Nacional de Estadística (Instituto Nacional de Estatística), órgão oficial de dados demográficos e censitários.} Censo 2022, INE\footnote{Instituto Nacional de Estadística (Instituto Nacional de Estatística), órgão oficial de dados demográficos e censitários.} EPHC 2023, Ministerio
Público 2024\\
\textbf{Nota:} Valores marcados como \emph{(dept.)} referem-se ao
departamento de Ñeembucú; valores \emph{(dist.)} são específicos deste
distrito. Dados marcados \emph{(est.)} são estimativas.

\begin{longtable}[]{@{}
  >{\raggedright\arraybackslash}p{(\linewidth - 4\tabcolsep) * \real{0.4231}}
  >{\raggedright\arraybackslash}p{(\linewidth - 4\tabcolsep) * \real{0.2692}}
  >{\raggedright\arraybackslash}p{(\linewidth - 4\tabcolsep) * \real{0.3077}}@{}}
\toprule\noalign{}
\begin{minipage}[b]{\linewidth}\raggedright
Indicador
\end{minipage} & \begin{minipage}[b]{\linewidth}\raggedright
Valor
\end{minipage} & \begin{minipage}[b]{\linewidth}\raggedright
Âmbito
\end{minipage} \\
\midrule\noalign{}
\endhead
\bottomrule\noalign{}
\endlastfoot
IDH (2020) & 0,710 & dept. \\
Ranking IDH nacional & 8/\allowbreak 18 & dept. \\
Esperança de vida & 73,0 anos & dept. \\
Escolaridade média & 8,2 anos & dept. \\
RNB per capita (USD PPA) & 8.800 & dept. \\
Pobreza monetária (\%) & 28,8\% & dept. \\
Pobreza extrema (\%) & 6,5\% & dept. \\
Índice de Gini & 0,445 & dept. \\
Acesso a água potável (\%) & 79,1\% & dept. \\
Acesso a saneamento (\%) & 77,4\% & dept. \\
Taxa de homicídios (est., /\allowbreak 100k hab) & 0,2--0,4 (2019--2020, fonte INE
--- mais baixa do país) & dept. \\
Índice de segurança & alto & dept. \\
População (Censo 2022) & 3.164 habitantes & dist. \\
Idade mediana & N/\allowbreak D & dist. \\
Taxa de fecundidade & N/\allowbreak D & dist. \\
\end{longtable}
\subsubsection{Combustível}

\textbf{Referência:} PETROPAR /\allowbreak  postos locais (2024) \textbf{Tipo de
localidade:} Interior

\begin{longtable}[]{@{}lll@{}}
\toprule\noalign{}
Combustível & USD/\allowbreak litro & Gs/\allowbreak litro (aprox.) \\
\midrule\noalign{}
\endhead
\bottomrule\noalign{}
\endlastfoot
Gasolina 93 oct & 0.99 & 7,326 \\
Gasolina 97 oct (premium) & 1.09 & 8,066 \\
Diesel & 0.92 & 6,808 \\
\end{longtable}

\begin{quote}
Preços podem variar ±5\% conforme posto e sazonalidade. Chaco e interior
remoto apresentam maior variação.
\end{quote}

\subsubsection{Cobertura Celular}

\textbf{Fonte:} CONATEL PY /\allowbreak  operadoras (2024)

\begin{longtable}[]{@{}ll@{}}
\toprule\noalign{}
Parâmetro & Valor \\
\midrule\noalign{}
\endhead
\bottomrule\noalign{}
\endlastfoot
Cobertura 4G população (dept.) & 75\% \\
Cobertura 4G área rural & 50\% \\
Melhor operadora & Tigo \\
Qualidade rural & moderada \\
\end{longtable}

\begin{quote}
Para áreas rurais fora do núcleo urbano, recomenda-se chip Tigo como
principal e Personal como backup.
\end{quote}

\subsubsection{Internet}

\textbf{Fonte:} CONATEL /\allowbreak  Speedtest Ookla (2024)

\begin{longtable}[]{@{}ll@{}}
\toprule\noalign{}
Parâmetro & Valor \\
\midrule\noalign{}
\endhead
\bottomrule\noalign{}
\endlastfoot
Velocidade média download & 30 Mbps \\
Domicílios com internet (dept.) & 40\% \\
Tecnologia predominante & rádio \\
Opção rural & Starlink disponível (\textasciitilde USD 44/\allowbreak mês) \\
\end{longtable}

\subsubsection{Mercado Imobiliário e Terra
Rural}

\textbf{Fonte:} INDERT /\allowbreak  Clasificados.com.py (2024)

\begin{longtable}[]{@{}ll@{}}
\toprule\noalign{}
Tipo & Referência \\
\midrule\noalign{}
\endhead
\bottomrule\noalign{}
\endlastfoot
Terra agrícola alta prod. (USD/\allowbreak ha) & 3,000 \\
Imóvel urbano (USD/\allowbreak m²) & 500 \\
Aluguel 2 quartos (USD/\allowbreak mês) & 200 \\
\end{longtable}

\begin{quote}
Valores de referência departamental. Localidades menores podem ter
preços 20--40\% abaixo da capital departamental. \#\#\# Saúde
\end{quote}

\textbf{Fonte:} MSPBS /\allowbreak  IPS Paraguay (2024-2026), consolidação
departamental e proxy local

\begin{longtable}[]{@{}ll@{}}
\toprule\noalign{}
Serviço & Disponibilidade \\
\midrule\noalign{}
\endhead
\bottomrule\noalign{}
\endlastfoot
USF /\allowbreak  Posto de Saúde & sim \\
Hospital Regional & não \\
IPS (seguro social) & não \\
Farmácia & sim \\
Distância ao hospital de referência & \textasciitilde1 km (Pilar) \\
\end{longtable}

\textbf{Principais estabelecimentos:} USF local; referencia hospitalar
em Pilar

\textbf{Observação para imigrantes:} Atendimento primario local ou em
raio curto; casos de maior complexidade seguem para Pilar. Cobertura
privada continua recomendada para especialidades e urgências de maior
complexidade.
\newpage
\secaoDiagnostico{Guazu Cua}{\mapaConteudo{-26.9833}{-58.1833}}{Guazú Cuá é um dos locais mais isolados e discretos da Região Oriental
do Paraguai. Oferece um santuário de privacidade para quem busca
desaparecer do radar metropolitano, sustentado por uma densidade
populacional quase nula e uma segurança social imbatível. Sua força é o
isolamento; sua fraqueza é a vulnerabilidade hidrológica extrema.
Sobreviver em Guazú Cuá exige um perfil autossuficiente capaz de operar
sem suporte externo por longos períodos, transformando o desafio das
águas em uma barreira natural de proteção contra crises demográficas
externas.

\subsubsection{Por que sim}

\begin{itemize}
\tightlist
\item
  Máximo isolamento geográfico e demográfico (anonimato total).
\item
  Ambiente social mais pacífico do departamento.
\item
  Abundância de recursos de proteína animal.
\item
  Longe de qualquer alvo militar ou eixo de fallout.
\end{itemize}

\subsubsection{Por que não}

\begin{itemize}
\tightlist
\item
  Risco de isolamento físico prolongado por inundações.
\item
  Solo inadequado para agricultura mecanizada diversificada.
\item
  Infraestrutura institucional e de saúde virtualmente inexistente.
\end{itemize}

\subsubsection{Recomendação
Técnica}

Altamente indicado para estabelecimentos de refúgio tático que priorizem
a invisibilidade e a segurança social sobre a facilidade logística.
Requer autonomia total em energia, água tratada e saúde. O GSS de 6.4
reflete a superioridade do isolamento protegida pela barreira natural
dos humedales. Referencia atual de score: GSS 6.4 em 2026-03-05.

\subsubsection{}

\begin{itemize}
\tightlist
\item
  \begin{enumerate}
  \def\labelenumi{\arabic{enumi})}
  \tightlist
  \item
    Dados censitários validados (INE\footnote{Instituto Nacional de Estadística (Instituto Nacional de Estatística), órgão oficial de dados demográficos e censitários.} 2022).
  \end{enumerate}
\item
  \begin{enumerate}
  \def\labelenumi{\arabic{enumi})}
  \setcounter{enumi}{1}
  \tightlist
  \item
    Levantamento territorial de baixada (Portal Geoestadístico).
  \end{enumerate}
\item
  \begin{enumerate}
  \def\labelenumi{\arabic{enumi})}
  \setcounter{enumi}{2}
  \tightlist
  \item
    Mapa de solos e esteros de Ñeembucú (MAG).
  \end{enumerate}
\item
  \begin{enumerate}
  \def\labelenumi{\arabic{enumi})}
  \setcounter{enumi}{3}
  \tightlist
  \item
    Registro de segurança pública departamental.
  \end{enumerate}
\end{itemize}}
\begin{table}[ht!]
\small \begin{tabular}{p{3.0cm}p{0.8cm}p{6.0cm}} \toprule
\textbf{Critério} & \textbf{Nota} & \textbf{Justificativa} \\ \midrule
A - Ameaças Estratégicas & 7.5 & Isolamento tático de alto nível; invisibilidade total e ausência de alvos estratégicos \\
B - Risco Social & 8.5 & Densidade populacional próxima de zero; segurança contra riscos de massa absoluta \\
C - Riscos Naturais & 4.0 & Risco hidrológico extremo; nota penalizada pela vulnerabilidade a inundações sazonais \\
D - Autossuficiência & 6.5 & Abundância de proteína animal; limitado pela inacessibilidade a insumos técnicos \\
E - Institucional & 5.0 & Ambiente social estável e seguro, mas com suporte institucional mínimo \\
\midrule \textbf{GSS FINAL} & \textbf{6.4} & \textbf{Moderadamente Seguro (Indicado para refúgio tático extremo de baixo perfil, exigindo alta autonomia logística).} \\ \bottomrule \end{tabular}
\end{table}
\subsubsection{Vulnerabilidades e
Mitigação}\label{vuln-Guazu_Cua-vulnerabilidades-e-mitigauxe7uxe3o}

\begin{itemize}
\tightlist
\item
  \textbf{Isolamento Geográfico Severo:} Risco de interrupção total de
  acessos terrestres (Mitigação: posse de veículo de tração
  pesada/\allowbreak trator e manutenção de estoque de segurança de suprimentos e
  combustível para 90 dias).
\item
  \textbf{Fragilidade das Redes:} Dependência de energia externa
  (Mitigação: instalação de sistemas solares fotovoltaicos robustos
  off-grid).
\item
  \textbf{Ausência de Saúde Especializada:} Distância crítica de
  hospitais (Mitigação: manutenção de estoque médico avançado e farmácia
  de emergência local).
\end{itemize}
\subsubsection{Dossiê de Campo}\label{dossie-Guazu_Cua-dossiuxea-de-campo}
\begin{description}[style=nextline,leftmargin=0.5cm,font=\bfseries]
\item[Coordenadas]  26°59'S 58°11'W.
\item[Topografia]  Relevo plano com predominância absoluta de esteros (pantanais) e terras baixas inundáveis. Altitude \textasciitilde 55m. Área vasta de \textasciitilde 1.130 km². Geologicamente muito estável, risco sísmico nulo.
\item[Alvos Estratégicos]  Áreas de pecuária extensiva. Ausência total de infraestrutura industrial ou bases militares, oferecendo a máxima invisibilidade estratégica e anonimato geográfico no sul paraguaio. Zona de amortecimento natural.
\item[Fallout]  Ventos predominantes do Norte (NE) e Leste. Baixíssimo risco de fallout devido ao profundo isolamento e distância de centros de comando.
\item[População]  1.393 habitantes (Censo 2022 Final). Um dos distritos mais desabitados do país.
\item[Densidade]  \textasciitilde 1,2 hab/\allowbreak km² (Densidade extremamente baixa, garantindo privacidade absoluta e ausência de riscos sociais de massa).
\item[Vias de Acesso]  Acesso via ramais de terra a partir da Ruta PY04. Localização que favorece a discrição, mas sujeita ao isolamento físico total durante períodos de chuvas torrenciais ou cheias dos esteros.
\item[Serviços]  Infraestrutura básica mínima. Dependência absoluta de Pilar (\textasciitilde 35 km) para atendimentos de saúde e serviços comerciais. A malha urbana é minúscula, com serviços reduzidos ao essencial.
\item[Custo de Vida]  Muito baixo; economia baseada na pecuária e no autoconsumo rural.
\item[Preço da Terra]  US\$ 700-1.000/\allowbreak ha (campos baixos para pecuária extensiva). Um dos menores valores de terra da Região Oriental.
\item[Hidrologia]  Risco de Inundação Crítico. Distrito situado no coração dos humedales de Ñeembucú. O acúmulo de águas pluviais e o transbordamento dos esteros podem isolar a localidade por semanas. Solo com baixíssima capacidade de drenagem.
\item[Clima]  Subtropical Úmido. Verões quentes; invernos frescos com alta umidade e nevoeiros frequentes.
\item[Energia]  Rede nacional (Itaipu/\allowbreak Yacyretá); altamente instável devido ao isolamento das linhas rurais e falta de redundância.
\item[Água]  Abundância de águas superficiais (esteros), embora exijam tratamento para consumo. Disponibilidade de poços artesianos produtivos.
\item[Qualidade do Solo]  Solo hidromórfico; aptidão exclusiva para pastagens naturais e gado de corte. Agricultura mecanizada inviável sem grandes investimentos em drenagem artificial.
\item[Resources Locais]  Grande excedente de gado bovino per capita. Elevado potencial para autossuficiência alimentar básica baseada em proteína animal para a pequena população residente.
\item[Segurança]  Zona rural excepcionalmente pacífica. Criminalidade urbana inexistente. Estabilidade social máxima baseada na resiliência da população tradicional aos desafios climáticos. Localidade "fora do mapa" para conflitos políticos ou sociais.
\item[Leis Local: Município consolidado. Livre de Restrição de Fronteira]  Localizado fora da faixa de 50km da fronteira internacional seca (limite tático dos esteros).
\end{description}
\paragraph{Solo (SoilGrids 2.0, média ponderada 0--30
cm)}

\begin{longtable}[]{@{}ll@{}}
\toprule\noalign{}
Parâmetro & Valor \\
\midrule\noalign{}
\endhead
\bottomrule\noalign{}
\endlastfoot
pH (H$_2$O) & 6.15 \\
Carbono orgânico (SOC) & 12.54 g/\allowbreak kg \\
Argila & 13.82 \% \\
Areia & 71.67 \% \\
Silte (calc.) & 14.5 \% \\
Densidade aparente & 1.41 g/\allowbreak cm³ \\
\textbf{Aptidão agrícola} & \textbf{Média} \\
\end{longtable}

Fonte: ISRIC SoilGrids 2.0 via WCS (acesso em 2026-03-21). Coords:
-26.9833°, -58.1833°. Média ponderada camadas 0-5, 5-15, 15-30 cm.
\subsubsection{Indicadores Sociais}

\textbf{Fontes:} PNUD 2020, INE\footnote{Instituto Nacional de Estadística (Instituto Nacional de Estatística), órgão oficial de dados demográficos e censitários.} Censo 2022, INE\footnote{Instituto Nacional de Estadística (Instituto Nacional de Estatística), órgão oficial de dados demográficos e censitários.} EPHC 2023, Ministerio
Público 2024\\
\textbf{Nota:} Valores marcados como \emph{(dept.)} referem-se ao
departamento de Ñeembucú; valores \emph{(dist.)} são específicos deste
distrito. Dados marcados \emph{(est.)} são estimativas.

\begin{longtable}[]{@{}
  >{\raggedright\arraybackslash}p{(\linewidth - 4\tabcolsep) * \real{0.4231}}
  >{\raggedright\arraybackslash}p{(\linewidth - 4\tabcolsep) * \real{0.2692}}
  >{\raggedright\arraybackslash}p{(\linewidth - 4\tabcolsep) * \real{0.3077}}@{}}
\toprule\noalign{}
\begin{minipage}[b]{\linewidth}\raggedright
Indicador
\end{minipage} & \begin{minipage}[b]{\linewidth}\raggedright
Valor
\end{minipage} & \begin{minipage}[b]{\linewidth}\raggedright
Âmbito
\end{minipage} \\
\midrule\noalign{}
\endhead
\bottomrule\noalign{}
\endlastfoot
IDH (2020) & 0,710 & dept. \\
Ranking IDH nacional & 8/\allowbreak 18 & dept. \\
Esperança de vida & 73,0 anos & dept. \\
Escolaridade média & 8,2 anos & dept. \\
RNB per capita (USD PPA) & 8.800 & dept. \\
Pobreza monetária (\%) & 28,8\% & dept. \\
Pobreza extrema (\%) & 6,5\% & dept. \\
Índice de Gini & 0,445 & dept. \\
Acesso a água potável (\%) & 79,1\% & dept. \\
Acesso a saneamento (\%) & 77,4\% & dept. \\
Taxa de homicídios (est., /\allowbreak 100k hab) & 0,2--0,4 (2019--2020, fonte INE
--- mais baixa do país) & dept. \\
Índice de segurança & alto & dept. \\
População (Censo 2022) & 1.393 habitantes & dist. \\
Idade mediana & N/\allowbreak D & dist. \\
Taxa de fecundidade & N/\allowbreak D & dist. \\
\end{longtable}
\subsubsection{Combustível}

\textbf{Referência:} PETROPAR /\allowbreak  postos locais (2024) \textbf{Tipo de
localidade:} Interior

\begin{longtable}[]{@{}lll@{}}
\toprule\noalign{}
Combustível & USD/\allowbreak litro & Gs/\allowbreak litro (aprox.) \\
\midrule\noalign{}
\endhead
\bottomrule\noalign{}
\endlastfoot
Gasolina 93 oct & 0.99 & 7,326 \\
Gasolina 97 oct (premium) & 1.09 & 8,066 \\
Diesel & 0.92 & 6,808 \\
\end{longtable}

\begin{quote}
Preços podem variar ±5\% conforme posto e sazonalidade. Chaco e interior
remoto apresentam maior variação.
\end{quote}

\subsubsection{Cobertura Celular}

\textbf{Fonte:} CONATEL PY /\allowbreak  operadoras (2024)

\begin{longtable}[]{@{}ll@{}}
\toprule\noalign{}
Parâmetro & Valor \\
\midrule\noalign{}
\endhead
\bottomrule\noalign{}
\endlastfoot
Cobertura 4G população (dept.) & 75\% \\
Cobertura 4G área rural & 50\% \\
Melhor operadora & Tigo \\
Qualidade rural & moderada \\
\end{longtable}

\begin{quote}
Para áreas rurais fora do núcleo urbano, recomenda-se chip Tigo como
principal e Personal como backup.
\end{quote}

\subsubsection{Internet}

\textbf{Fonte:} CONATEL /\allowbreak  Speedtest Ookla (2024)

\begin{longtable}[]{@{}ll@{}}
\toprule\noalign{}
Parâmetro & Valor \\
\midrule\noalign{}
\endhead
\bottomrule\noalign{}
\endlastfoot
Velocidade média download & 30 Mbps \\
Domicílios com internet (dept.) & 40\% \\
Tecnologia predominante & rádio \\
Opção rural & Starlink disponível (\textasciitilde USD 44/\allowbreak mês) \\
\end{longtable}

\subsubsection{Mercado Imobiliário e Terra
Rural}

\textbf{Fonte:} INDERT /\allowbreak  Clasificados.com.py (2024)

\begin{longtable}[]{@{}ll@{}}
\toprule\noalign{}
Tipo & Referência \\
\midrule\noalign{}
\endhead
\bottomrule\noalign{}
\endlastfoot
Terra agrícola alta prod. (USD/\allowbreak ha) & 3,000 \\
Imóvel urbano (USD/\allowbreak m²) & 500 \\
Aluguel 2 quartos (USD/\allowbreak mês) & 200 \\
\end{longtable}

\begin{quote}
Valores de referência departamental. Localidades menores podem ter
preços 20--40\% abaixo da capital departamental. \#\#\# Saúde
\end{quote}

\textbf{Fonte:} MSPBS /\allowbreak  IPS Paraguay (2024-2026), consolidação
departamental e proxy local

\begin{longtable}[]{@{}ll@{}}
\toprule\noalign{}
Serviço & Disponibilidade \\
\midrule\noalign{}
\endhead
\bottomrule\noalign{}
\endlastfoot
USF /\allowbreak  Posto de Saúde & sim \\
Hospital Regional & não \\
IPS (seguro social) & não \\
Farmácia & sim \\
Distância ao hospital de referência & \textasciitilde42 km (Pilar) \\
\end{longtable}

\textbf{Principais estabelecimentos:} USF local; referencia hospitalar
em Pilar

\textbf{Observação para imigrantes:} Atendimento primario local ou em
raio curto; casos de maior complexidade seguem para Pilar. Cobertura
privada continua recomendada para especialidades e urgências de maior
complexidade.
\newpage
\secaoDiagnostico{Humaita}{\mapaConteudo{-27.0666}{-58.5000}}{Humaitá é a sentinela silenciosa do sul paraguaio. Oferece um isolamento
geográfico e demográfico de alto nível, protegido pela história e pela
geografia fluvial que a tornaram impenetrável no passado. Sua força
absoluta reside na abundância de recursos hídricos e proteicos, e na
segurança social sustentada por sua importância heroica. É o local ideal
para quem busca estabelecer um refúgio tático anfíbio, transformando o
``Recodo de Humaitá'' em uma fortaleza de proteção individual contra
colapsos sistêmicos externos.

\subsubsection{Por que sim}

\begin{itemize}
\tightlist
\item
  Máximo isolamento demográfico e tático (privacidade absoluta).
\item
  Abundância ilimitada de águas puras e recursos pesqueiros/\allowbreak pecuários.
\item
  Ambiente social extremamente seguro sob vigilância naval.
\item
  Longe de qualquer alvo militar terrestre primário ou eixo de fallout.
\end{itemize}

\subsubsection{Por que não}

\begin{itemize}
\tightlist
\item
  Risco de isolamento terrestre por chuvas intensas.
\item
  Dependência de Pilar para serviços de saúde de alta complexidade.
\item
  Restrições da Lei de Fronteira para investidores estrangeiros.
\end{itemize}

\subsubsection{Recomendação
Técnica}

Altamente indicado para residências de isolamento tático que possuam
infraestrutura própria de transporte fluvial. Requer autonomia
energética individual para mitigar riscos de rede rural. O GSS de 7.9
reflete a excepcional solidez do isolamento geográfico e a riqueza de
recursos naturais da localidade. Referencia atual de score: GSS 7.9 em
2026-03-05.

\subsubsection{}

\begin{itemize}
\tightlist
\item
  \begin{enumerate}
  \def\labelenumi{\arabic{enumi})}
  \tightlist
  \item
    Dados censitários validados (INE\footnote{Instituto Nacional de Estadística (Instituto Nacional de Estatística), órgão oficial de dados demográficos e censitários.} 2022).
  \end{enumerate}
\item
  \begin{enumerate}
  \def\labelenumi{\arabic{enumi})}
  \setcounter{enumi}{1}
  \tightlist
  \item
    Relatórios de monitoramento fluvial do Rio Paraguai (ANNP).
  \end{enumerate}
\item
  \begin{enumerate}
  \def\labelenumi{\arabic{enumi})}
  \setcounter{enumi}{2}
  \tightlist
  \item
    Mapa de solos e bacias hídricas de Ñeembucú (Portal Geoestadístico).
  \end{enumerate}
\item
  \begin{enumerate}
  \def\labelenumi{\arabic{enumi})}
  \setcounter{enumi}{3}
  \tightlist
  \item
    Registro de segurança pública e governança municipal histórica.
  \end{enumerate}
\end{itemize}}
\begin{table}[ht!]
\small \begin{tabular}{p{3.0cm}p{0.8cm}p{6.0cm}} \toprule
\textbf{Critério} & \textbf{Nota} & \textbf{Justificativa} \\ \midrule
A - Ameaças Estratégicas & 8.0 & Posição tática defensiva histórica; curva do rio e isolamento geográfico superior; baixíssima visibilidade de alvos industriais \\
B - Risco Social & 8.5 & Densidade populacional extremamente baixa; isolamento social de alto nível e privacidade absoluta \\
C - Riscos Naturais & 7.0 & Geologia muito estável e segurança relativa contra inundações na zona urbana histórica \\
D - Autossuficiência & 8.5 & Recursos hídricos e de proteína animal ilimitados; solo diversificado para pastagens \\
E - Institucional & 7.5 & Forte segurança institucional naval e suporte social baseado na identidade heroica \\
\midrule \textbf{GSS FINAL} & \textbf{7.9} & \textbf{Seguro (Altamente recomendado para quem busca isolamento tático fluvial e anonimato em ambiente histórico seguro).} \\ \bottomrule \end{tabular}
\end{table}
\subsubsection{Vulnerabilidades e
Mitigação}\label{vuln-Humaita-vulnerabilidades-e-mitigauxe7uxe3o}

\begin{itemize}
\tightlist
\item
  \textbf{Isolamento Terrestre Sazonal:} Risco de isolamento por chuvas
  em estradas de terra (Mitigação: posse de embarcação a motor como via
  de escape alternativa e manutenção de estoque de suprimentos para 60
  dias).
\item
  \textbf{Fragilidade Médica Local:} Necessidade de deslocamento para
  Pilar (Mitigação: manutenção de estoque médico básico e kit trauma
  local).
\item
  \textbf{Vulnerabilidade de Costa:} Risco de erosão e cheias
  ribeirinhas (Mitigação: construção exclusivamente em terrenos elevados
  acima da cota histórica).
\end{itemize}
\subsubsection{Dossiê de Campo}\label{dossie-Humaita-dossiuxea-de-campo}
\begin{description}[style=nextline,leftmargin=0.5cm,font=\bfseries]
\item[Coordenadas]  27°04'S 58°30'W.
\item[Topografia]  Relevo predominantemente plano, situado em uma curva estratégica ("Recodo de Humaitá") do Rio Paraguai. Altitude \textasciitilde 55m. Área de \textasciitilde 321 km². Geologicamente estável.
\item[Alvos Estratégicos]  Ruínas de San Carlos de Borromeo (Local histórico de imenso valor militar e simbólico; ponto de observação fluvial); Porto fluvial histórico; Proximidade tática com a confluência dos rios Paraguai e Paraná. Conhecida como a "Sebastopol das Américas", oferece isolamento geográfico e defesa natural provida pelos esteros e pelo rio.
\item[Fallout]  Ventos predominantes do Norte (NE) e Leste. Baixíssimo risco de fallout de alvos continentais primários devido ao isolamento geográfico profundo.
\item[População]  2.391 habitantes (Censo 2022 Final). População resiliente com forte identidade histórica e ligação com o rio.
\item[Densidade]  \textasciitilde 7,4 hab/\allowbreak km² (Densidade extremamente baixa, garantindo privacidade absoluta e ausência de riscos sociais de massa urbana).
\item[Vias de Acesso]  Acesso via Ruta PY04 e ramais de terra/\allowbreak empedrados secundários. Localização em "bolsão de segurança" geográfico, afastada dos grandes eixos de trânsito comercial, o que favorece o anonimato estratégico. Acesso fluvial contínuo.
\item[Serviços]  Infraestrutura básica funcional (USF e Centro de Saúde). Dependência de Pilar (\textasciitilde 40 km) para atendimentos de alta complexidade. Centro comercial pequeno focado no turismo histórico e pesca.
\item[Custo de Vida]  Muito baixo; economia baseada na pecuária, pesca e turismo.
\item[Preço da Terra]  US\$ 1.500-3.000/\allowbreak ha (áreas com costa de rio); US\$ 800-1.200/\allowbreak ha (campos baixos).
\item[Hidrologia]  Risco de Inundação Moderado. Vulnerável às cheias sazonais do Rio Paraguai. A vila histórica possui sistema de proteção, mas áreas rurais baixas podem ser isoladas. Abundância hídrica superficial inesgotável.
\item[Clima]  Subtropical Úmido. Verões intensos e chuvosos; invernos amenos com alta umidade influenciada pelo rio.
\item[Energia]  Rede nacional (Itaipu/\allowbreak Yacyretá); instável em áreas remotas.
\item[Água]  Abundância hídrica absoluta via Rio Paraguai e poços artesianos profundos. Água superficial exige tratamento.
\item[Qualidade do Solo]  Solo aluvial e argiloso; excelente para pastagens naturais e pecuária de corte.
\item[Resources Locais]  Grande excedente de gado bovino e recursos pesqueiros. Elevadíssimo potencial para autossuficiência de proteína animal para a pequena população local.
\item[Segurança]  Localidade excepcionalmente pacífica e segura. Criminalidade urbana inexistente. Estabilidade institucional sólida baseada na presença de destacamentos navais e na importância histórica nacional. Coesão social fortíssima.
\item[Leis Local: Município consolidado. Restrição de Fronteira (Lei 2532/\allowbreak 05)]  Aplicável a terras rurais na faixa de 50km da fronteira argentina (limite fluvial).
\end{description}
\paragraph{Solo (SoilGrids 2.0, média ponderada 0--30
cm)}

\begin{longtable}[]{@{}ll@{}}
\toprule\noalign{}
Parâmetro & Valor \\
\midrule\noalign{}
\endhead
\bottomrule\noalign{}
\endlastfoot
pH (H$_2$O) & 6.08 \\
Carbono orgânico (SOC) & 14.29 g/\allowbreak kg \\
Argila & 16.77 \% \\
Areia & 64.23 \% \\
Silte (calc.) & 19.0 \% \\
Densidade aparente & 1.36 g/\allowbreak cm³ \\
\textbf{Aptidão agrícola} & \textbf{Média} \\
\end{longtable}

Fonte: ISRIC SoilGrids 2.0 via WCS (acesso em 2026-03-21). Coords:
-27.0667°, -58.5°. Média ponderada camadas 0-5, 5-15, 15-30 cm.
\subsubsection{Indicadores Sociais}

\textbf{Fontes:} PNUD 2020, INE\footnote{Instituto Nacional de Estadística (Instituto Nacional de Estatística), órgão oficial de dados demográficos e censitários.} Censo 2022, INE\footnote{Instituto Nacional de Estadística (Instituto Nacional de Estatística), órgão oficial de dados demográficos e censitários.} EPHC 2023, Ministerio
Público 2024\\
\textbf{Nota:} Valores marcados como \emph{(dept.)} referem-se ao
departamento de Ñeembucú; valores \emph{(dist.)} são específicos deste
distrito. Dados marcados \emph{(est.)} são estimativas.

\begin{longtable}[]{@{}
  >{\raggedright\arraybackslash}p{(\linewidth - 4\tabcolsep) * \real{0.4231}}
  >{\raggedright\arraybackslash}p{(\linewidth - 4\tabcolsep) * \real{0.2692}}
  >{\raggedright\arraybackslash}p{(\linewidth - 4\tabcolsep) * \real{0.3077}}@{}}
\toprule\noalign{}
\begin{minipage}[b]{\linewidth}\raggedright
Indicador
\end{minipage} & \begin{minipage}[b]{\linewidth}\raggedright
Valor
\end{minipage} & \begin{minipage}[b]{\linewidth}\raggedright
Âmbito
\end{minipage} \\
\midrule\noalign{}
\endhead
\bottomrule\noalign{}
\endlastfoot
IDH (2020) & 0,710 & dept. \\
Ranking IDH nacional & 8/\allowbreak 18 & dept. \\
Esperança de vida & 73,0 anos & dept. \\
Escolaridade média & 6.6 anos & dist. \\
RNB per capita (USD PPA) & 8.800 & dept. \\
Pobreza monetária (\%) & 28,8\% & dept. \\
Pobreza extrema (\%) & 6,5\% & dept. \\
Índice de Gini & 0,445 & dept. \\
Acesso a água potável (\%) & 79,1\% & dept. \\
Acesso a saneamento (\%) & 77,4\% & dept. \\
Taxa de homicídios (est., /\allowbreak 100k hab) & 0,2--0,4 (2019--2020, fonte INE
--- mais baixa do país) & dept. \\
Índice de segurança & alto & dept. \\
População (Censo 2022) & 2.391 habitantes & dist. \\
Idade mediana & N/\allowbreak D & dist. \\
Taxa de fecundidade & N/\allowbreak D & dist. \\
\end{longtable}
\subsubsection{Combustível}

\textbf{Referência:} PETROPAR /\allowbreak  postos locais (2024) \textbf{Tipo de
localidade:} Interior

\begin{longtable}[]{@{}lll@{}}
\toprule\noalign{}
Combustível & USD/\allowbreak litro & Gs/\allowbreak litro (aprox.) \\
\midrule\noalign{}
\endhead
\bottomrule\noalign{}
\endlastfoot
Gasolina 93 oct & 0.99 & 7,326 \\
Gasolina 97 oct (premium) & 1.09 & 8,066 \\
Diesel & 0.92 & 6,808 \\
\end{longtable}

\begin{quote}
Preços podem variar ±5\% conforme posto e sazonalidade. Chaco e interior
remoto apresentam maior variação.
\end{quote}

\subsubsection{Cobertura Celular}

\textbf{Fonte:} CONATEL PY /\allowbreak  operadoras (2024)

\begin{longtable}[]{@{}ll@{}}
\toprule\noalign{}
Parâmetro & Valor \\
\midrule\noalign{}
\endhead
\bottomrule\noalign{}
\endlastfoot
Cobertura 4G população (dept.) & 75\% \\
Cobertura 4G área rural & 50\% \\
Melhor operadora & Tigo \\
Qualidade rural & moderada \\
\end{longtable}

\begin{quote}
Para áreas rurais fora do núcleo urbano, recomenda-se chip Tigo como
principal e Personal como backup.
\end{quote}

\subsubsection{Internet}

\textbf{Fonte:} CONATEL /\allowbreak  Speedtest Ookla (2024)

\begin{longtable}[]{@{}ll@{}}
\toprule\noalign{}
Parâmetro & Valor \\
\midrule\noalign{}
\endhead
\bottomrule\noalign{}
\endlastfoot
Velocidade média download & 30 Mbps \\
Domicílios com internet (dept.) & 40\% \\
Tecnologia predominante & rádio \\
Opção rural & Starlink disponível (\textasciitilde USD 44/\allowbreak mês) \\
\end{longtable}

\subsubsection{Mercado Imobiliário e Terra
Rural}

\textbf{Fonte:} INDERT /\allowbreak  Clasificados.com.py (2024)

\begin{longtable}[]{@{}ll@{}}
\toprule\noalign{}
Tipo & Referência \\
\midrule\noalign{}
\endhead
\bottomrule\noalign{}
\endlastfoot
Terra agrícola alta prod. (USD/\allowbreak ha) & 3,000 \\
Imóvel urbano (USD/\allowbreak m²) & 500 \\
Aluguel 2 quartos (USD/\allowbreak mês) & 200 \\
\end{longtable}

\begin{quote}
Valores de referência departamental. Localidades menores podem ter
preços 20--40\% abaixo da capital departamental. \#\#\# Saúde
\end{quote}

\textbf{Fonte:} MSPBS /\allowbreak  IPS Paraguay (2024-2026), consolidação
departamental e proxy local

\begin{longtable}[]{@{}ll@{}}
\toprule\noalign{}
Serviço & Disponibilidade \\
\midrule\noalign{}
\endhead
\bottomrule\noalign{}
\endlastfoot
USF /\allowbreak  Posto de Saúde & sim \\
Hospital Regional & não \\
IPS (seguro social) & não \\
Farmácia & sim \\
Distância ao hospital de referência & \textasciitilde225 km (Pilar) \\
\end{longtable}

\textbf{Principais estabelecimentos:} USF local; referencia hospitalar
em Pilar

\textbf{Observação para imigrantes:} Atendimento primario local ou em
raio curto; casos de maior complexidade seguem para Pilar. Cobertura
privada continua recomendada para especialidades e urgências de maior
complexidade.
\newpage
\secaoDiagnostico{Isla Umbu}{\mapaConteudo{-26.9000}{-58.2000}}{Isla Umbú é a joia tática oculta de Ñeembucú. Oferece o isolamento
geográfico e demográfico de uma vila remota com a conveniência logística
de estar a poucos minutos da capital departamental (Pilar). Sua força
absoluta reside na estabilidade de seu ambiente social e na abundância
de recursos proteicos e hídricos. É o local ideal para quem busca
estabelecer um refúgio de alto padrão, transformando a geografia de
``ilha de terra firme'' em uma defesa natural contra instabilidades
climáticas e sociais metropolitanas. Isla Umbú representa a simbiose
perfeita entre anonimato rural e suporte institucional urbano.

\subsubsection{Por que sim}

\begin{itemize}
\tightlist
\item
  Localidade nº 2 no Ranking Nacional de Segurança Estratégica.
\item
  Ambiente social mais seguro e pacífico da Região Oriental.
\item
  Abundância massiva de água mineral e recursos proteicos (leite/\allowbreak carne).
\item
  Conectividade rápida com serviços capitais via Ruta PY04.
\end{itemize}

\subsubsection{Por que não}

\begin{itemize}
\tightlist
\item
  Dependência de suprimentos técnicos especializados de Pilar.
\item
  Risco de visibilidade por proximidade com o único eixo rodoviário
  principal.
\item
  Infraestrutura de saúde local limitada a atendimentos primários.
\end{itemize}

\subsubsection{Recomendação
Técnica}

Altamente indicado para residência permanente de refúgio tático ou base
de autossuficiência de alto padrão. Requer autonomia energética
individual para mitigar riscos de rede rural. O GSS de 8.6 reflete a
excepcional solidez do isolamento demográfico reforçada pela base de
recursos naturais e suporte institucional próximo. Referencia atual de
score: GSS 8.6 em 2026-03-05.

\subsubsection{}

\begin{itemize}
\tightlist
\item
  \begin{enumerate}
  \def\labelenumi{\arabic{enumi})}
  \tightlist
  \item
    Dados censitários validados (INE\footnote{Instituto Nacional de Estadística (Instituto Nacional de Estatística), órgão oficial de dados demográficos e censitários.} 2022).
  \end{enumerate}
\item
  \begin{enumerate}
  \def\labelenumi{\arabic{enumi})}
  \setcounter{enumi}{1}
  \tightlist
  \item
    Relatórios de governança comunitária e segurança rural.
  \end{enumerate}
\item
  \begin{enumerate}
  \def\labelenumi{\arabic{enumi})}
  \setcounter{enumi}{2}
  \tightlist
  \item
    Mapa de solos e bacias hídricas de Ñeembucú (Portal Geoestadístico).
  \end{enumerate}
\item
  \begin{enumerate}
  \def\labelenumi{\arabic{enumi})}
  \setcounter{enumi}{3}
  \tightlist
  \item
    Registro de segurança pública e proximidade logística com Pilar.
  \end{enumerate}
\end{itemize}}
\begin{table}[ht!]
\small \begin{tabular}{p{3.0cm}p{0.8cm}p{6.0cm}} \toprule
\textbf{Critério} & \textbf{Nota} & \textbf{Justificativa} \\ \midrule
A - Ameaças Estratégicas & 8.5 & Posição tática excepcional; "ilha" de solo firme, isolamento de baixo perfil e proximidade com suporte institucional capital \\
B - Risco Social & 9.0 & Densidade populacional extremamente baixa; isolamento social de alto nível e privacidade absoluta \\
C - Riscos Naturais & 7.5 & Segurança superior contra inundações na região dos esteros; geologia muito estável \\
D - Autossuficiência & 8.5 & Excelentes recursos de proteína animal e água mineral; solo diversificado \\
E - Institucional & 9.0 & Máxima segurança social e suporte institucional capital adjacente; governança comunitária exemplar \\
\midrule \textbf{GSS FINAL} & \textbf{8.6} & \textbf{Seguro (A localidade nº 2 do Paraguai para relocação estratégica; oferece o máximo em privacidade, recursos básicos e proximidade com infraestrutura capital).} \\ \bottomrule \end{tabular}
\end{table}
\subsubsection{Vulnerabilidades e
Mitigação}\label{vuln-Isla_Umbu-vulnerabilidades-e-mitigauxe7uxe3o}

\begin{itemize}
\tightlist
\item
  \textbf{Pressão Migratória de Proximidade:} Risco de receber
  refugiados de Pilar em crises hídricas urbanas (Mitigação: residência
  rural afastada 5 km do centro da vila e do asfalto da PY04).
\item
  \textbf{Dependência de Redes Públicas:} Vulnerabilidade a colapsos
  sistêmicos de energia (Mitigação: instalação de sistemas solares
  fotovoltaicos e aproveitamento da força hídrica local).
\item
  \textbf{Escala Comercial:} Dependência de Pilar para insumos técnicos
  complexos (Mitigação: manutenção de estoque de manutenção preventiva
  anual).
\end{itemize}
\subsubsection{Dossiê de Campo}\label{dossie-Isla_Umbu-dossiuxea-de-campo}
\begin{description}[style=nextline,leftmargin=0.5cm,font=\bfseries]
\item[Coordenadas]  26°54'S 58°12'W.
\item[Topografia]  Relevo predominantemente plano a suavemente ondulado, situado em uma zona de "terra firme" cercada por vastos humedales. Altitude \textasciitilde 60m. Área de \textasciitilde 485 km². Geologicamente estável, risco sísmico nulo.
\item[Alvos Estratégicos]  Igreja de San Atanasio (Patrimônio histórico e ponto de observação tático); Polo de produção pecuária leiteira; Local de isolamento estratégico adjacente à capital departamental (Pilar). Ausência total de alvos militares ou industriais pesados, oferecendo altíssimo nível de invisibilidade e proteção por anonimato rural.
\item[Fallout]  Ventos predominantes do Norte (NE) e Leste. Baixíssimo risco de fallout de alvos continentais primários devido ao isolamento geográfico profundo no sudoeste do país.
\item[População]  2.166 habitantes (Censo 2022 Final). População estável com forte coesão comunitária e tradição rural.
\item[Densidade]  \textasciitilde 4,5 hab/\allowbreak km² (Densidade extremamente baixa, garantindo privacidade absoluta e ausência de riscos sociais de massa urbana).
\item[Vias de Acesso]  Localização privilegiada sobre a Ruta PY04 (pavimentada). Conectividade de elite para Pilar (\textasciitilde 12 km) e para o leste do país. Oferece a logística rápida de uma capital departamental com o isolamento tático de uma vila rural.
\item[Serviços]  Unidade de Saúde da Família (USF) local. Proximidade imediata com os hospitais de alta complexidade em Pilar. Serviços comerciais básicos funcionais.
\item[Custo de Vida]  Muito baixo; economia baseada na pecuária leiteira e remessas.
\item[Preço da Terra]  US\$ 1.500-2.500/\allowbreak ha (campos pecuários); US\$ 30-50/\allowbreak m² (Urbano).
\item[Hidrologia]  Baixo risco de inundações sistêmicas na zona urbana e áreas de invernada alta. Situada em uma "ilha" de solo firme que a protege das cheias cíclicas do Rio Paraguai. Abundância de arroios perenes.
\item[Clima]  Subtropical Úmido. Verões quentes; invernos suaves com entradas de ar frio que limpam a atmosfera.
\item[Energia]  Rede nacional (Itaipu); boa estabilidade devido à proximidade com o nodo de Pilar.
\item[Água]  Abundante via poços artesianos e nascentes naturais; excelente qualidade hídrica subterrânea.
\item[Qualidade do Solo]  Solo argiloso e aluvial fértil; aptidão para pastagens naturais, milho e horticultura familiar.
\item[Resources Locais]  Grande polo produtor de leite e queijo para a região sul. Elevadíssimo potencial para autossuficiência de proteína animal em nível de propriedade devido à vasta área produtiva e baixíssima população.
\item[Segurança]  Uma das localidades mais seguras do Paraguai. Criminalidade urbana comum praticamente inexistente. Estabilidade social forjada pela cultura tradicional e vigilância comunitária orgânica. Presença institucional de suporte via Pilar.
\item[Leis Local]  Município consolidado. \\ \textbf{Livre de Restrição} de Fronteira. \\ \textit{Aquisição de terra rural proibida para estrangeiros limítrofes.}
\end{description}
\paragraph{Solo (SoilGrids 2.0, média ponderada 0--30
cm)}

\begin{longtable}[]{@{}ll@{}}
\toprule\noalign{}
Parâmetro & Valor \\
\midrule\noalign{}
\endhead
\bottomrule\noalign{}
\endlastfoot
pH (H$_2$O) & 6.15 \\
Carbono orgânico (SOC) & 12.53 g/\allowbreak kg \\
Argila & 17.08 \% \\
Areia & 60.98 \% \\
Silte (calc.) & 21.9 \% \\
Densidade aparente & 1.42 g/\allowbreak cm³ \\
\textbf{Aptidão agrícola} & \textbf{Média} \\
\end{longtable}

Fonte: ISRIC SoilGrids 2.0 via WCS (acesso em 2026-03-21). Coords:
-26.9°, -58.2°. Média ponderada camadas 0-5, 5-15, 15-30 cm.
\subsubsection{Indicadores Sociais}

\textbf{Fontes:} PNUD 2020, INE\footnote{Instituto Nacional de Estadística (Instituto Nacional de Estatística), órgão oficial de dados demográficos e censitários.} Censo 2022, INE\footnote{Instituto Nacional de Estadística (Instituto Nacional de Estatística), órgão oficial de dados demográficos e censitários.} EPHC 2023, Ministerio
Público 2024\\
\textbf{Nota:} Valores marcados como \emph{(dept.)} referem-se ao
departamento de Ñeembucú; valores \emph{(dist.)} são específicos deste
distrito. Dados marcados \emph{(est.)} são estimativas.

\begin{longtable}[]{@{}
  >{\raggedright\arraybackslash}p{(\linewidth - 4\tabcolsep) * \real{0.4231}}
  >{\raggedright\arraybackslash}p{(\linewidth - 4\tabcolsep) * \real{0.2692}}
  >{\raggedright\arraybackslash}p{(\linewidth - 4\tabcolsep) * \real{0.3077}}@{}}
\toprule\noalign{}
\begin{minipage}[b]{\linewidth}\raggedright
Indicador
\end{minipage} & \begin{minipage}[b]{\linewidth}\raggedright
Valor
\end{minipage} & \begin{minipage}[b]{\linewidth}\raggedright
Âmbito
\end{minipage} \\
\midrule\noalign{}
\endhead
\bottomrule\noalign{}
\endlastfoot
IDH (2020) & 0,710 & dept. \\
Ranking IDH nacional & 8/\allowbreak 18 & dept. \\
Esperança de vida & 73,0 anos & dept. \\
Escolaridade média & 7.4 anos & dist. \\
RNB per capita (USD PPA) & 8.800 & dept. \\
Pobreza monetária (\%) & 28,8\% & dept. \\
Pobreza extrema (\%) & 6,5\% & dept. \\
Índice de Gini & 0,445 & dept. \\
Acesso a água potável (\%) & 79,1\% & dept. \\
Acesso a saneamento (\%) & 77,4\% & dept. \\
Taxa de homicídios (est., /\allowbreak 100k hab) & 0,2--0,4 (2019--2020, fonte INE
--- mais baixa do país) & dept. \\
Índice de segurança & alto & dept. \\
População (Censo 2022) & 2.166 habitantes & dist. \\
Idade mediana & N/\allowbreak D & dist. \\
Taxa de fecundidade & N/\allowbreak D & dist. \\
\end{longtable}
\subsubsection{Combustível}

\textbf{Referência:} PETROPAR /\allowbreak  postos locais (2024) \textbf{Tipo de
localidade:} Interior

\begin{longtable}[]{@{}lll@{}}
\toprule\noalign{}
Combustível & USD/\allowbreak litro & Gs/\allowbreak litro (aprox.) \\
\midrule\noalign{}
\endhead
\bottomrule\noalign{}
\endlastfoot
Gasolina 93 oct & 0.99 & 7,326 \\
Gasolina 97 oct (premium) & 1.09 & 8,066 \\
Diesel & 0.92 & 6,808 \\
\end{longtable}

\begin{quote}
Preços podem variar ±5\% conforme posto e sazonalidade. Chaco e interior
remoto apresentam maior variação.
\end{quote}

\subsubsection{Cobertura Celular}

\textbf{Fonte:} CONATEL PY /\allowbreak  operadoras (2024)

\begin{longtable}[]{@{}ll@{}}
\toprule\noalign{}
Parâmetro & Valor \\
\midrule\noalign{}
\endhead
\bottomrule\noalign{}
\endlastfoot
Cobertura 4G população (dept.) & 75\% \\
Cobertura 4G área rural & 50\% \\
Melhor operadora & Tigo \\
Qualidade rural & moderada \\
\end{longtable}

\begin{quote}
Para áreas rurais fora do núcleo urbano, recomenda-se chip Tigo como
principal e Personal como backup.
\end{quote}

\subsubsection{Internet}

\textbf{Fonte:} CONATEL /\allowbreak  Speedtest Ookla (2024)

\begin{longtable}[]{@{}ll@{}}
\toprule\noalign{}
Parâmetro & Valor \\
\midrule\noalign{}
\endhead
\bottomrule\noalign{}
\endlastfoot
Velocidade média download & 30 Mbps \\
Domicílios com internet (dept.) & 40\% \\
Tecnologia predominante & rádio \\
Opção rural & Starlink disponível (\textasciitilde USD 44/\allowbreak mês) \\
\end{longtable}

\subsubsection{Mercado Imobiliário e Terra
Rural}

\textbf{Fonte:} INDERT /\allowbreak  Clasificados.com.py (2024)

\begin{longtable}[]{@{}ll@{}}
\toprule\noalign{}
Tipo & Referência \\
\midrule\noalign{}
\endhead
\bottomrule\noalign{}
\endlastfoot
Terra agrícola alta prod. (USD/\allowbreak ha) & 3,000 \\
Imóvel urbano (USD/\allowbreak m²) & 500 \\
Aluguel 2 quartos (USD/\allowbreak mês) & 200 \\
\end{longtable}

\begin{quote}
Valores de referência departamental. Localidades menores podem ter
preços 20--40\% abaixo da capital departamental. \#\#\# Saúde
\end{quote}

\textbf{Fonte:} MSPBS /\allowbreak  IPS Paraguay (2024-2026), consolidação
departamental e proxy local

\begin{longtable}[]{@{}ll@{}}
\toprule\noalign{}
Serviço & Disponibilidade \\
\midrule\noalign{}
\endhead
\bottomrule\noalign{}
\endlastfoot
USF /\allowbreak  Posto de Saúde & sim \\
Hospital Regional & sim \\
IPS (seguro social) & não \\
Farmácia & sim \\
Distância ao hospital de referência & \textasciitilde18 km (Pilar) \\
\end{longtable}

\textbf{Principais estabelecimentos:} USF local; referencia hospitalar
em Pilar

\textbf{Observação para imigrantes:} Atendimento primario local ou em
raio curto; casos de maior complexidade seguem para Pilar. Cobertura
privada continua recomendada para especialidades e urgências de maior
complexidade.
\newpage
\secaoDiagnostico{Laureles}{\mapaConteudo{-27.2500}{-58.0166}}{Laureles representa a fronteira final do isolamento na Região Oriental.
Oferece um nível de privacidade e segurança social quase inatingível em
outras partes do país, agindo como um ``buraco negro'' estratégico fora
do alcance de qualquer crise demográfica urbana. Sua força é a
invisibilidade; sua fraqueza é a vulnerabilidade absoluta à natureza das
águas. Sobreviver em Laureles exige um perfil de ocupante altamente
técnico e autônomo, capaz de gerir sua própria saúde e logística sem
suporte estatal, transformando o isolamento geográfico em uma armadura
de proteção definitiva contra o colapso sistêmico.

\subsubsection{Por que sim}

\begin{itemize}
\tightlist
\item
  Máximo isolamento geográfico e demográfico (anonimato total).
\item
  Ambiente social seguro e estável.
\item
  Abundância de recursos de proteína animal (gado).
\item
  Distância máxima de qualquer eixo de conflito ou fallout.
\end{itemize}

\subsubsection{Por que não}

\begin{itemize}
\tightlist
\item
  Risco de isolamento físico por inundações prolongadas.
\item
  Inexistência de serviços de saúde especializados no distrito.
\item
  Infraestrutura institucional mínima e dependência de Pilar.
\end{itemize}

\subsubsection{Recomendação
Técnica}

Indicado exclusivamente para estabelecimentos de autossuficiência de
alto nível que busquem isolamento tático absoluto. Requer autonomia
energética, médica e logística total (90+ dias). O GSS de 5.5 reflete a
segurança do isolamento em contraste com a extrema fragilidade
infraestrutural e riscos naturais da localidade. Referencia atual de
score: GSS 5.5 em 2026-03-05.

\subsubsection{}

\begin{itemize}
\tightlist
\item
  \begin{enumerate}
  \def\labelenumi{\arabic{enumi})}
  \tightlist
  \item
    Dados censitários validados (INE\footnote{Instituto Nacional de Estadística (Instituto Nacional de Estatística), órgão oficial de dados demográficos e censitários.} 2022).
  \end{enumerate}
\item
  \begin{enumerate}
  \def\labelenumi{\arabic{enumi})}
  \setcounter{enumi}{1}
  \tightlist
  \item
    Relatórios de monitoramento de cheias da Secretaria de Emergência
    Nacional (SEN\footnote{Secretaría de Emergencia Nacional (Secretaria de Emergência Nacional), órgão governamental de resposta a desastres.}).
  \end{enumerate}
\item
  \begin{enumerate}
  \def\labelenumi{\arabic{enumi})}
  \setcounter{enumi}{2}
  \tightlist
  \item
    Mapa de solos e bacias hidrográficas de Ñeembucú (MAG).
  \end{enumerate}
\item
  \begin{enumerate}
  \def\labelenumi{\arabic{enumi})}
  \setcounter{enumi}{3}
  \tightlist
  \item
    Registro de segurança pública departamental.
  \end{enumerate}
\end{itemize}}
\begin{table}[ht!]
\small \begin{tabular}{p{3.0cm}p{0.8cm}p{6.0cm}} \toprule
\textbf{Critério} & \textbf{Nota} & \textbf{Justificativa} \\ \midrule
A - Ameaças Estratégicas & 7.0 & Isolamento geográfico profundo; invisibilidade estratégica total e ausência de alvos \\
B - Risco Social & 8.0 & Densidade populacional insignificante; segurança contra riscos de massa absoluta \\
C - Riscos Naturais & 3.5 & Risco hidrológico extremo; nota penalizada pela vulnerabilidade a inundações prolongadas \\
D - Autossuficiência & 6.0 & Bons recursos de gado; severamente limitado pela inacessibilidade a insumos \\
E - Institucional & 2.0 & Nota penalizada no ranking nacional para refletir a dependência extrema de cadeias externas e infraestrutura de saúde mínima local - GSS 5.5 total \\
\midrule \textbf{GSS FINAL} & \textbf{5.5} & \textbf{Moderadamente Seguro (Indicado apenas para projetos de isolamento tático extremo que possuam autonomia total de saúde e logística).} \\ \bottomrule \end{tabular}
\end{table}
\subsubsection{Vulnerabilidades e
Mitigação}\label{vuln-Laureles-vulnerabilidades-e-mitigauxe7uxe3o}

\begin{itemize}
\tightlist
\item
  \textbf{Isolamento Geográfico Extremo:} Risco de ficar
  ``desconectado'' por longos períodos (Mitigação: posse de veículo 4x4
  robusto e trator; manutenção de estoque de suprimentos para 120 dias).
\item
  \textbf{Inexistência de Saúde Local:} Distância crítica de qualquer
  centro hospitalar (Mitigação: manutenção de posto médico avançado
  privado e sistema de evacuação aérea/\allowbreak fluvial).
\item
  \textbf{Instabilidade Elétrica Rural:} Quedas frequentes da rede
  (Mitigação: instalação de sistemas solares fotovoltaicos potentes
  off-grid).
\end{itemize}
\subsubsection{Dossiê de Campo}\label{dossie-Laureles-dossiuxea-de-campo}
\begin{description}[style=nextline,leftmargin=0.5cm,font=\bfseries]
\item[Coordenadas]  27°15'S 58°01'W.
\item[Topografia]  Relevo plano com predominância de humedales e esteros (ex: Estero Ñeembucú e Pira Guazú). Altitude \textasciitilde 50m. Área de \textasciitilde 800 km². Geologicamente estável, risco sísmico desprezível.
\item[Alvos Estratégicos]  Áreas de produção pecuária tradicional extensiva. Localização estratégica no extremo sul do Paraguai, próxima à confluência dos rios Paraguai e Paraná. Ausência total de alvos militares ou infraestrutura industrial, oferecendo a máxima invisibilidade estratégica e anonimato tático.
\item[Fallout]  Ventos predominantes do Norte (NE) e Leste. Baixíssimo risco de fallout de alvos continentais devido ao extremo isolamento geográfico.
\item[População]  2.323 habitantes (Censo 2022 Final). População em declínio migratório, restando apenas comunidades tradicionais resilientes.
\item[Densidade]  \textasciitilde 2,9 hab/\allowbreak km² (Densidade extremamente baixa, garantindo privacidade absoluta e ausência total de riscos sociais de massa urbana).
\item[Vias de Acesso]  Acesso via ramais da Ruta PY20. Localização remota com infraestrutura viária precária, dependente de estradas de terra que tornam o acesso crítico ou impossível durante grandes cheias.
\item[Serviços]  Infraestrutura básica mínima. Dependência absoluta de Pilar (\textasciitilde 80 km) para serviços de saúde, logística comercial e administrativa. A vila urbana é minúscula, focada em serviços de subsistência.
\item[Custo de Vida]  Muito baixo; economia baseada na pecuária e no autoconsumo rural.
\item[Preço da Terra]  US\$ 700-1.100/\allowbreak ha (campos baixos para pecuária). Valores de entrada entre os menores do Paraguai oriental.
\item[Hidrologia]  Risco de Inundação Crítico. Distrito situado em zona de baixada extrema. Vulnerável a cheias sazonais e extraordinárias que isolam a localidade por meses. O solo hidromórfico retém umidade, dificultando o trânsito pesado.
\item[Clima]  Subtropical Úmido. Verões intensos; invernos frescos com alta umidade e nevoeiros frequentes.
\item[Energia]  Rede nacional (Itaipu/\allowbreak Yacyretá); sujeita a interrupções prolongadas devido à localização em final de rede rural e dificuldade de manutenção de acesso.
\item[Água]  Abundância de águas superficiais (esteros) e disponibilidade de poços artesianos. Águas superficiais exigem tratamento rigoroso.
\item[Qualidade do Solo]  Solo hidromórfico e arenoso; aptidão para pastagens naturais e gado de corte. Baixa aptidão para agricultura mecanizada sem vultosos investimentos em engenharia hídrica.
\item[Resources Locais]  Grande produção de gado bovino. Elevado potencial para autossuficiência de proteína animal para a pequena população local.
\item[Segurança]  Zona rural excepcionalmente pacífica. Criminalidade urbana inexistente. Estabilidade social baseada no isolamento geográfico e na cultura tradicional. Ausência total de visibilidade para conflitos externos ou sociais.
\item[Leis Local: Município consolidado. Restrição de Fronteira (Lei 2532/\allowbreak 05)]  Aplicável a terras rurais na faixa de 50km da fronteira argentina.
\end{description}
\paragraph{Solo (SoilGrids 2.0, média ponderada 0--30
cm)}

\begin{longtable}[]{@{}ll@{}}
\toprule\noalign{}
Parâmetro & Valor \\
\midrule\noalign{}
\endhead
\bottomrule\noalign{}
\endlastfoot
pH (H$_2$O) & 5.8 \\
Carbono orgânico (SOC) & 22.1 g/\allowbreak kg \\
Argila & 11.08 \% \\
Areia & 74.02 \% \\
Silte (calc.) & 14.9 \% \\
Densidade aparente & 1.33 g/\allowbreak cm³ \\
\textbf{Aptidão agrícola} & \textbf{Alta} \\
\end{longtable}

Fonte: ISRIC SoilGrids 2.0 via WCS (acesso em 2026-03-21). Coords:
-27.25°, -58.0167°. Média ponderada camadas 0-5, 5-15, 15-30 cm.

\begin{itemize}
\tightlist
\item
  \textbf{Homicidios/\allowbreak 100k:} \textasciitilde0,0 (Paz absoluta).
\end{itemize}

\subsection{10. Analise de Riscos}

\begin{itemize}
\tightlist
\item
  \textbf{Cenario curto prazo:} Manutenção da tranquilidade rural e
  estabilidade da pecuária tradicional.
\item
  \textbf{Cenario medio prazo:} Valorização das terras como ``bolsão de
  segurança'' tático isolado.
\end{itemize}

\subsection{11. Redacao Final}

\subsubsection{Diagnostico Integrado}

Laureles representa a fronteira final do isolamento na Região Oriental.
Oferece um nível de privacidade e segurança social quase inatingível em
outras partes do país, agindo como um ``buraco negro'' estratégico fora
do alcance de qualquer crise demográfica urbana. Sua força é a
invisibilidade; sua fraqueza é a vulnerabilidade absoluta à natureza das
águas. Sobreviver em Laureles exige um perfil de ocupante altamente
técnico e autônomo, capaz de gerir sua própria saúde e logística sem
suporte estatal, transformando o isolamento geográfico em uma armadura
de proteção definitiva contra o colapso sistêmico.

\subsubsection{Por que sim}

\begin{itemize}
\tightlist
\item
  Máximo isolamento geográfico e demográfico (anonimato total).
\item
  Ambiente social seguro e estável.
\item
  Abundância de recursos de proteína animal (gado).
\item
  Distância máxima de qualquer eixo de conflito ou fallout.
\end{itemize}

\subsubsection{Por que nao}

\begin{itemize}
\tightlist
\item
  Risco de isolamento físico por inundações prolongadas.
\item
  Inexistência de serviços de saúde especializados no distrito.
\item
  Infraestrutura institucional mínima e dependência de Pilar.
\end{itemize}

\subsubsection{Recomendacao Tecnica}

Indicado exclusivamente para estabelecimentos de autossuficiência de
alto nível que busquem isolamento tático absoluto. Requer autonomia
energética, médica e logística total (90+ dias). O GSS de 5.5 reflete a
segurança do isolamento em contraste com a extrema fragilidade
infraestrutural e riscos naturais da localidade. Referencia atual de
score: GSS 5.5 em 2026-03-05.

\subsection{13.}

\begin{itemize}
\tightlist
\item
  \begin{enumerate}
  \def\labelenumi{\arabic{enumi})}
  \tightlist
  \item
    Dados censitários validados (INE\footnote{Instituto Nacional de Estadística (Instituto Nacional de Estatística), órgão oficial de dados demográficos e censitários.} 2022).
  \end{enumerate}
\item
  \begin{enumerate}
  \def\labelenumi{\arabic{enumi})}
  \setcounter{enumi}{1}
  \tightlist
  \item
    Relatórios de monitoramento de cheias da Secretaria de Emergência
    Nacional (SEN\footnote{Secretaría de Emergencia Nacional (Secretaria de Emergência Nacional), órgão governamental de resposta a desastres.}).
  \end{enumerate}
\item
  \begin{enumerate}
  \def\labelenumi{\arabic{enumi})}
  \setcounter{enumi}{2}
  \tightlist
  \item
    Mapa de solos e bacias hidrográficas de Ñeembucú (MAG).
  \end{enumerate}
\item
  \begin{enumerate}
  \def\labelenumi{\arabic{enumi})}
  \setcounter{enumi}{3}
  \tightlist
  \item
    Registro de segurança pública departamental.
  \end{enumerate}
\end{itemize}

\subsubsection{Combustível}

\textbf{Referência:} PETROPAR /\allowbreak  postos locais (2024) \textbf{Tipo de
localidade:} Interior

\begin{longtable}[]{@{}lll@{}}
\toprule\noalign{}
Combustível & USD/\allowbreak litro & Gs/\allowbreak litro (aprox.) \\
\midrule\noalign{}
\endhead
\bottomrule\noalign{}
\endlastfoot
Gasolina 93 oct & 0.99 & 7,326 \\
Gasolina 97 oct (premium) & 1.09 & 8,066 \\
Diesel & 0.92 & 6,808 \\
\end{longtable}

\begin{quote}
Preços podem variar ±5\% conforme posto e sazonalidade. Chaco e interior
remoto apresentam maior variação.
\end{quote}

\subsubsection{Cobertura Celular}

\textbf{Fonte:} CONATEL PY /\allowbreak  operadoras (2024)

\begin{longtable}[]{@{}ll@{}}
\toprule\noalign{}
Parâmetro & Valor \\
\midrule\noalign{}
\endhead
\bottomrule\noalign{}
\endlastfoot
Cobertura 4G população (dept.) & 75\% \\
Cobertura 4G área rural & 50\% \\
Melhor operadora & Tigo \\
Qualidade rural & moderada \\
\end{longtable}

\begin{quote}
Para áreas rurais fora do núcleo urbano, recomenda-se chip Tigo como
principal e Personal como backup.
\end{quote}

\subsubsection{Internet}

\textbf{Fonte:} CONATEL /\allowbreak  Speedtest Ookla (2024)

\begin{longtable}[]{@{}ll@{}}
\toprule\noalign{}
Parâmetro & Valor \\
\midrule\noalign{}
\endhead
\bottomrule\noalign{}
\endlastfoot
Velocidade média download & 30 Mbps \\
Domicílios com internet (dept.) & 40\% \\
Tecnologia predominante & rádio \\
Opção rural & Starlink disponível (\textasciitilde USD 44/\allowbreak mês) \\
\end{longtable}

\subsubsection{Mercado Imobiliário e Terra
Rural}

\textbf{Fonte:} INDERT /\allowbreak  Clasificados.com.py (2024)

\begin{longtable}[]{@{}ll@{}}
\toprule\noalign{}
Tipo & Referência \\
\midrule\noalign{}
\endhead
\bottomrule\noalign{}
\endlastfoot
Terra agrícola alta prod. (USD/\allowbreak ha) & 3,000 \\
Imóvel urbano (USD/\allowbreak m²) & 500 \\
Aluguel 2 quartos (USD/\allowbreak mês) & 200 \\
\end{longtable}

\begin{quote}
Valores de referência departamental. Localidades menores podem ter
preços 20--40\% abaixo da capital departamental. \#\#\# Saúde
\end{quote}

\textbf{Fonte:} MSPBS /\allowbreak  IPS Paraguay (2024-2026), consolidação
departamental e proxy local

\begin{longtable}[]{@{}ll@{}}
\toprule\noalign{}
Serviço & Disponibilidade \\
\midrule\noalign{}
\endhead
\bottomrule\noalign{}
\endlastfoot
USF /\allowbreak  Posto de Saúde & sim \\
Hospital Regional & não \\
IPS (seguro social) & não \\
Farmácia & sim \\
Distância ao hospital de referência & \textasciitilde118 km (Pilar) \\
\end{longtable}

\textbf{Principais estabelecimentos:} USF local; referencia hospitalar
em Pilar

\textbf{Observação para imigrantes:} Atendimento primario local ou em
raio curto; casos de maior complexidade seguem para Pilar. Cobertura
privada continua recomendada para especialidades e urgências de maior
complexidade.
\subsubsection{Indicadores Sociais}

\textbf{Fontes:} PNUD 2020, INE\footnote{Instituto Nacional de Estadística (Instituto Nacional de Estatística), órgão oficial de dados demográficos e censitários.} Censo 2022, INE\footnote{Instituto Nacional de Estadística (Instituto Nacional de Estatística), órgão oficial de dados demográficos e censitários.} EPHC 2023, Ministerio
Público 2024\\
\textbf{Nota:} Valores marcados como \emph{(dept.)} referem-se ao
departamento de Ñeembucú; valores \emph{(dist.)} são específicos deste
distrito. Dados marcados \emph{(est.)} são estimativas.

\begin{longtable}[]{@{}
  >{\raggedright\arraybackslash}p{(\linewidth - 4\tabcolsep) * \real{0.4231}}
  >{\raggedright\arraybackslash}p{(\linewidth - 4\tabcolsep) * \real{0.2692}}
  >{\raggedright\arraybackslash}p{(\linewidth - 4\tabcolsep) * \real{0.3077}}@{}}
\toprule\noalign{}
\begin{minipage}[b]{\linewidth}\raggedright
Indicador
\end{minipage} & \begin{minipage}[b]{\linewidth}\raggedright
Valor
\end{minipage} & \begin{minipage}[b]{\linewidth}\raggedright
Âmbito
\end{minipage} \\
\midrule\noalign{}
\endhead
\bottomrule\noalign{}
\endlastfoot
IDH (2020) & 0,710 & dept. \\
Ranking IDH nacional & 8/\allowbreak 18 & dept. \\
Esperança de vida & 73,0 anos & dept. \\
Escolaridade média & 6.6 anos & dist. \\
RNB per capita (USD PPA) & 8.800 & dept. \\
Pobreza monetária (\%) & 28,8\% & dept. \\
Pobreza extrema (\%) & 6,5\% & dept. \\
Índice de Gini & 0,445 & dept. \\
Acesso a água potável (\%) & 79,1\% & dept. \\
Acesso a saneamento (\%) & 77,4\% & dept. \\
Taxa de homicídios (est., /\allowbreak 100k hab) & 0,2--0,4 (2019--2020, fonte INE
--- mais baixa do país) & dept. \\
Índice de segurança & alto & dept. \\
População (Censo 2022) & 2.323 habitantes & dist. \\
Idade mediana & N/\allowbreak D & dist. \\
Taxa de fecundidade & N/\allowbreak D & dist. \\
\end{longtable}
\subsubsection{Combustível}

\textbf{Referência:} PETROPAR /\allowbreak  postos locais (2024) \textbf{Tipo de
localidade:} Interior

\begin{longtable}[]{@{}lll@{}}
\toprule\noalign{}
Combustível & USD/\allowbreak litro & Gs/\allowbreak litro (aprox.) \\
\midrule\noalign{}
\endhead
\bottomrule\noalign{}
\endlastfoot
Gasolina 93 oct & 0.99 & 7,326 \\
Gasolina 97 oct (premium) & 1.09 & 8,066 \\
Diesel & 0.92 & 6,808 \\
\end{longtable}

\begin{quote}
Preços podem variar ±5\% conforme posto e sazonalidade. Chaco e interior
remoto apresentam maior variação.
\end{quote}

\subsubsection{Cobertura Celular}

\textbf{Fonte:} CONATEL PY /\allowbreak  operadoras (2024)

\begin{longtable}[]{@{}ll@{}}
\toprule\noalign{}
Parâmetro & Valor \\
\midrule\noalign{}
\endhead
\bottomrule\noalign{}
\endlastfoot
Cobertura 4G população (dept.) & 75\% \\
Cobertura 4G área rural & 50\% \\
Melhor operadora & Tigo \\
Qualidade rural & moderada \\
\end{longtable}

\begin{quote}
Para áreas rurais fora do núcleo urbano, recomenda-se chip Tigo como
principal e Personal como backup.
\end{quote}

\subsubsection{Internet}

\textbf{Fonte:} CONATEL /\allowbreak  Speedtest Ookla (2024)

\begin{longtable}[]{@{}ll@{}}
\toprule\noalign{}
Parâmetro & Valor \\
\midrule\noalign{}
\endhead
\bottomrule\noalign{}
\endlastfoot
Velocidade média download & 30 Mbps \\
Domicílios com internet (dept.) & 40\% \\
Tecnologia predominante & rádio \\
Opção rural & Starlink disponível (\textasciitilde USD 44/\allowbreak mês) \\
\end{longtable}

\subsubsection{Mercado Imobiliário e Terra
Rural}

\textbf{Fonte:} INDERT /\allowbreak  Clasificados.com.py (2024)

\begin{longtable}[]{@{}ll@{}}
\toprule\noalign{}
Tipo & Referência \\
\midrule\noalign{}
\endhead
\bottomrule\noalign{}
\endlastfoot
Terra agrícola alta prod. (USD/\allowbreak ha) & 3,000 \\
Imóvel urbano (USD/\allowbreak m²) & 500 \\
Aluguel 2 quartos (USD/\allowbreak mês) & 200 \\
\end{longtable}

\begin{quote}
Valores de referência departamental. Localidades menores podem ter
preços 20--40\% abaixo da capital departamental. \#\#\# Saúde
\end{quote}

\textbf{Fonte:} MSPBS /\allowbreak  IPS Paraguay (2024-2026), consolidação
departamental e proxy local

\begin{longtable}[]{@{}ll@{}}
\toprule\noalign{}
Serviço & Disponibilidade \\
\midrule\noalign{}
\endhead
\bottomrule\noalign{}
\endlastfoot
USF /\allowbreak  Posto de Saúde & sim \\
Hospital Regional & não \\
IPS (seguro social) & não \\
Farmácia & sim \\
Distância ao hospital de referência & \textasciitilde118 km (Pilar) \\
\end{longtable}

\textbf{Principais estabelecimentos:} USF local; referencia hospitalar
em Pilar

\textbf{Observação para imigrantes:} Atendimento primario local ou em
raio curto; casos de maior complexidade seguem para Pilar. Cobertura
privada continua recomendada para especialidades e urgências de maior
complexidade.
\newpage
\secaoDiagnostico{Mayor Martinez}{\mapaConteudo{-27.0166}{-58.3333}}{Mayor José de Jesús Martínez é um enclave de produtividade e baixo
perfil situado na margem do Rio Paraguai. Oferece um isolamento
demográfico que garante privacidade absoluta, ao mesmo tempo em que se
beneficia da proximidade tática com a capital departamental. Sua força
reside na abundância de recursos hídricos e na segurança social de uma
comunidade pequena e ordeira. É o local ideal para quem busca
estabelecer um refúgio tático com acesso fluvial, transformando o
distanciamento metropolitano em uma barreira natural de proteção contra
crises sistêmicas externas.

\subsubsection{Por que sim}

\begin{itemize}
\tightlist
\item
  Máximo isolamento geográfico e demográfico em território seguro.
\item
  Abundância ilimitada de águas puras e recursos pesqueiros/\allowbreak pecuários.
\item
  Ambiente social pacífico e estável sob vigilância naval.
\item
  Longe de qualquer alvo militar terrestre primário.
\end{itemize}

\subsubsection{Por que não}

\begin{itemize}
\tightlist
\item
  Risco de isolamento terrestre por chuvas intensas.
\item
  Dependência de Pilar para serviços de saúde complexos.
\item
  Restrições da Lei de Fronteira para investidores estrangeiros.
\end{itemize}

\subsubsection{Recomendação
Técnica}

Altamente indicado para residências de isolamento tático com
infraestrutura de transporte fluvial. Requer autonomia energética
individual para mitigar riscos de rede rural. O GSS de 7.3 reflete a
solidez do isolamento geográfico e a riqueza de recursos naturais da
localidade. Referencia atual de score: GSS 7.3 em 2026-03-05.

\subsubsection{}

\begin{itemize}
\tightlist
\item
  \begin{enumerate}
  \def\labelenumi{\arabic{enumi})}
  \tightlist
  \item
    Dados censitários validados (INE\footnote{Instituto Nacional de Estadística (Instituto Nacional de Estatística), órgão oficial de dados demográficos e censitários.} 2022).
  \end{enumerate}
\item
  \begin{enumerate}
  \def\labelenumi{\arabic{enumi})}
  \setcounter{enumi}{1}
  \tightlist
  \item
    Relatórios de monitoramento fluvial do Rio Paraguai (ANNP).
  \end{enumerate}
\item
  \begin{enumerate}
  \def\labelenumi{\arabic{enumi})}
  \setcounter{enumi}{2}
  \tightlist
  \item
    Mapa de solos e bacias hídricas de Ñeembucú (Portal Geoestadístico).
  \end{enumerate}
\item
  \begin{enumerate}
  \def\labelenumi{\arabic{enumi})}
  \setcounter{enumi}{3}
  \tightlist
  \item
    Registro de segurança pública e governança municipal.
  \end{enumerate}
\end{itemize}}
\begin{table}[ht!]
\small \begin{tabular}{p{3.0cm}p{0.8cm}p{6.0cm}} \toprule
\textbf{Critério} & \textbf{Nota} & \textbf{Justificativa} \\ \midrule
A - Ameaças Estratégicas & 7.0 & Posição tática estratégica; margem do rio e isolamento geográfico satisfatório; baixíssima visibilidade de alvos \\
B - Risco Social & 8.0 & Densidade populacional baixa; isolamento social de bom nível e privacidade garantida \\
C - Riscos Naturais & 6.5 & Geologia estável; nota penalizada pela vulnerabilidade a inundações ribeirinhas \\
D - Autossuficiência & 7.5 & Excelentes recursos de proteína animal e água fluvial ilimitada \\
E - Institucional & 7.0 & Estabilidade institucional e segurança regional funcional próxima à capital \\
\midrule \textbf{GSS FINAL} & \textbf{7.3} & \textbf{Seguro (Altamente recomendado para quem busca isolamento tático rural com proximidade logística a Pilar).} \\ \bottomrule \end{tabular}
\end{table}
\subsubsection{Vulnerabilidades e
Mitigação}\label{vuln-Mayor_Martinez-vulnerabilidades-e-mitigauxe7uxe3o}

\begin{itemize}
\tightlist
\item
  \textbf{Isolamento Terrestre Sazonal:} Risco de isolamento por chuvas
  (Mitigação: posse de embarcação a motor e manutenção de estoque de
  suprimentos para 60 dias).
\item
  \textbf{Fragilidade Médica:} Necessidade de deslocamento para Pilar
  (Mitigação: manutenção de estoque médico básico e kit trauma local).
\item
  \textbf{Dependência de Redes:} Vulnerabilidade a quedas elétricas
  (Mitigação: instalação de sistemas solares fotovoltaicos).
\end{itemize}
\subsubsection{Dossiê de Campo}\label{dossie-Mayor_Martinez-dossiuxea-de-campo}
\begin{description}[style=nextline,leftmargin=0.5cm,font=\bfseries]
\item[Coordenadas]  27°01'S 58°20'W.
\item[Topografia]  Relevo predominantemente plano, margeado pelo Rio Paraguai ao oeste. Situado em zona de transição entre o Rio Paraguai e os grandes esteros internos. Altitude \textasciitilde 55m. Área de \textasciitilde 181 km². Geologicamente estável.
\item[Alvos Estratégicos]  Porto fluvial secundário; Polo de produção pecuária tradicional; Local de vigilância fluvial estratégica no sul do departamento. Ausência de bases militares massivas, oferecendo excelente invisibilidade estratégica e proteção por isolamento geográfico deliberado.
\item[Fallout]  Ventos predominantes do Norte (NE) e Leste. Baixíssimo risco de fallout de alvos continentais primários devido ao isolamento geográfico.
\item[População]  3.037 habitantes (Censo 2022 Final). População estável com forte tradição rural e ligação com a atividade pecuária.
\item[Densidade]  \textasciitilde 16,8 hab/\allowbreak km² (Densidade baixa, garantindo privacidade absoluta e ausência de riscos sociais de massa urbana).
\item[Vias de Acesso]  Acesso via ramais de terra a partir da Ruta PY04. Localização que favorece a discrição, mas sujeita a períodos de isolamento físico devido à precariedade das vias vicinais durante épocas de chuvas intensas. Acesso fluvial funcional.
\item[Serviços]  Infraestrutura básica funcional (USF local). Dependência de Pilar (\textasciitilde 30 km) para atendimentos de saúde de alta complexidade e logística comercial diversificada.
\item[Custo de Vida]  Muito baixo; economia baseada na pecuária bovina e pesca comercial.
\item[Preço da Terra]  US\$ 1.500-2.500/\allowbreak ha (áreas ribeirinhas); US\$ 800-1.200/\allowbreak ha (campos baixos).
\item[Hidrologia]  Risco de Inundação Moderado. Vulnerável às cheias sazonais do Rio Paraguai e transbordamento de arroios locais. Solo hidromórfico retém água por períodos prolongados em áreas de baixada.
\item[Clima]  Subtropical Úmido. Verões intensos e chuvosos; invernos amenos com alta umidade relativa.
\item[Energia]  Rede nacional (Itaipu/\allowbreak Yacyretá); instável em áreas rurais remotas.
\item[Água]  Abundância hídrica superficial (Rio Paraguai) e poços artesianos produtivos. Água superficial exige tratamento.
\item[Qualidade do Solo]  Solo aluvial e argiloso; excelente para pastagens naturais e pecuária de corte. Agricultura mecanizada limitada pela escala.
\item[Resources Locais]  Grande excedente de gado bovino e recursos pesqueiros. Elevadíssimo potencial para autossuficiência de proteína animal para a pequena população local.
\item[Segurança]  Localidade excepcionalmente pacífica e segura. Criminalidade urbana praticamente inexistente. Estabilidade institucional sólida baseada na presença de postos de controle fluvial e coesão comunitária tradicional.
\item[Leis Local: Município consolidado. Restrição de Fronteira (Lei 2532/\allowbreak 05)]  Aplicável a terras rurais na faixa de 50km da fronteira argentina (limite fluvial).
\end{description}
\paragraph{Solo (SoilGrids 2.0, média ponderada 0--30
cm)}

\begin{longtable}[]{@{}ll@{}}
\toprule\noalign{}
Parâmetro & Valor \\
\midrule\noalign{}
\endhead
\bottomrule\noalign{}
\endlastfoot
pH (H$_2$O) & 6.05 \\
Carbono orgânico (SOC) & 12.06 g/\allowbreak kg \\
Argila & 18.67 \% \\
Areia & 56.58 \% \\
Silte (calc.) & 24.8 \% \\
Densidade aparente & 1.39 g/\allowbreak cm³ \\
\textbf{Aptidão agrícola} & \textbf{Alta} \\
\end{longtable}

Fonte: ISRIC SoilGrids 2.0 via WCS (acesso em 2026-03-21). Coords:
-27.0167°, -58.3333°. Média ponderada camadas 0-5, 5-15, 15-30 cm.
\subsubsection{Indicadores Sociais}

\textbf{Fontes:} PNUD 2020, INE\footnote{Instituto Nacional de Estadística (Instituto Nacional de Estatística), órgão oficial de dados demográficos e censitários.} Censo 2022, INE\footnote{Instituto Nacional de Estadística (Instituto Nacional de Estatística), órgão oficial de dados demográficos e censitários.} EPHC 2023, Ministerio
Público 2024\\
\textbf{Nota:} Valores marcados como \emph{(dept.)} referem-se ao
departamento de Ñeembucú; valores \emph{(dist.)} são específicos deste
distrito. Dados marcados \emph{(est.)} são estimativas.

\begin{longtable}[]{@{}
  >{\raggedright\arraybackslash}p{(\linewidth - 4\tabcolsep) * \real{0.4231}}
  >{\raggedright\arraybackslash}p{(\linewidth - 4\tabcolsep) * \real{0.2692}}
  >{\raggedright\arraybackslash}p{(\linewidth - 4\tabcolsep) * \real{0.3077}}@{}}
\toprule\noalign{}
\begin{minipage}[b]{\linewidth}\raggedright
Indicador
\end{minipage} & \begin{minipage}[b]{\linewidth}\raggedright
Valor
\end{minipage} & \begin{minipage}[b]{\linewidth}\raggedright
Âmbito
\end{minipage} \\
\midrule\noalign{}
\endhead
\bottomrule\noalign{}
\endlastfoot
IDH (2020) & 0,710 & dept. \\
Ranking IDH nacional & 8/\allowbreak 18 & dept. \\
Esperança de vida & 73,0 anos & dept. \\
Escolaridade média & 8,2 anos & dept. \\
RNB per capita (USD PPA) & 8.800 & dept. \\
Pobreza monetária (\%) & 28,8\% & dept. \\
Pobreza extrema (\%) & 6,5\% & dept. \\
Índice de Gini & 0,445 & dept. \\
Acesso a água potável (\%) & 79,1\% & dept. \\
Acesso a saneamento (\%) & 77,4\% & dept. \\
Taxa de homicídios (est., /\allowbreak 100k hab) & 0,2--0,4 (2019--2020, fonte INE
--- mais baixa do país) & dept. \\
Índice de segurança & alto & dept. \\
População (Censo 2022) & 3.037 habitantes & dist. \\
Idade mediana & N/\allowbreak D & dist. \\
Taxa de fecundidade & N/\allowbreak D & dist. \\
\end{longtable}
\subsubsection{Combustível}

\textbf{Referência:} PETROPAR /\allowbreak  postos locais (2024) \textbf{Tipo de
localidade:} Interior

\begin{longtable}[]{@{}lll@{}}
\toprule\noalign{}
Combustível & USD/\allowbreak litro & Gs/\allowbreak litro (aprox.) \\
\midrule\noalign{}
\endhead
\bottomrule\noalign{}
\endlastfoot
Gasolina 93 oct & 0.99 & 7,326 \\
Gasolina 97 oct (premium) & 1.09 & 8,066 \\
Diesel & 0.92 & 6,808 \\
\end{longtable}

\begin{quote}
Preços podem variar ±5\% conforme posto e sazonalidade. Chaco e interior
remoto apresentam maior variação.
\end{quote}

\subsubsection{Cobertura Celular}

\textbf{Fonte:} CONATEL PY /\allowbreak  operadoras (2024)

\begin{longtable}[]{@{}ll@{}}
\toprule\noalign{}
Parâmetro & Valor \\
\midrule\noalign{}
\endhead
\bottomrule\noalign{}
\endlastfoot
Cobertura 4G população (dept.) & 75\% \\
Cobertura 4G área rural & 50\% \\
Melhor operadora & Tigo \\
Qualidade rural & moderada \\
\end{longtable}

\begin{quote}
Para áreas rurais fora do núcleo urbano, recomenda-se chip Tigo como
principal e Personal como backup.
\end{quote}

\subsubsection{Internet}

\textbf{Fonte:} CONATEL /\allowbreak  Speedtest Ookla (2024)

\begin{longtable}[]{@{}ll@{}}
\toprule\noalign{}
Parâmetro & Valor \\
\midrule\noalign{}
\endhead
\bottomrule\noalign{}
\endlastfoot
Velocidade média download & 30 Mbps \\
Domicílios com internet (dept.) & 40\% \\
Tecnologia predominante & rádio \\
Opção rural & Starlink disponível (\textasciitilde USD 44/\allowbreak mês) \\
\end{longtable}

\subsubsection{Mercado Imobiliário e Terra
Rural}

\textbf{Fonte:} INDERT /\allowbreak  Clasificados.com.py (2024)

\begin{longtable}[]{@{}ll@{}}
\toprule\noalign{}
Tipo & Referência \\
\midrule\noalign{}
\endhead
\bottomrule\noalign{}
\endlastfoot
Terra agrícola alta prod. (USD/\allowbreak ha) & 3,000 \\
Imóvel urbano (USD/\allowbreak m²) & 500 \\
Aluguel 2 quartos (USD/\allowbreak mês) & 200 \\
\end{longtable}

\begin{quote}
Valores de referência departamental. Localidades menores podem ter
preços 20--40\% abaixo da capital departamental. \#\#\# Saúde
\end{quote}

\textbf{Fonte:} MSPBS /\allowbreak  IPS Paraguay (2024-2026), consolidação
departamental e proxy local

\begin{longtable}[]{@{}ll@{}}
\toprule\noalign{}
Serviço & Disponibilidade \\
\midrule\noalign{}
\endhead
\bottomrule\noalign{}
\endlastfoot
USF /\allowbreak  Posto de Saúde & sim \\
Hospital Regional & não \\
IPS (seguro social) & não \\
Farmácia & sim \\
Distância ao hospital de referência & \textasciitilde41 km (Pilar) \\
\end{longtable}

\textbf{Principais estabelecimentos:} USF local; referencia hospitalar
em Pilar

\textbf{Observação para imigrantes:} Atendimento primario local ou em
raio curto; casos de maior complexidade seguem para Pilar. Cobertura
privada continua recomendada para especialidades e urgências de maior
complexidade.
\newpage
\secaoDiagnostico{Paso de Patria}{\mapaConteudo{-27.2666}{-58.5666}}{Paso de Patria é a sentinela do extremo sudoeste paraguaio. Oferece um
isolamento tático e demográfico de alto nível, protegido pela geografia
dos grandes rios e vastos esteros. Sua força absoluta reside na
abundância inesgotável de recursos hídricos e pesqueiros, além do
anonimato garantido pela distância dos centros de poder. O desafio
existencial da localidade são as águas; viver em Paso de Patria exige um
perfil anfíbio, com autonomia de transporte fluvial e estoque logístico
massivo para enfrentar períodos de isolamento terrestre. É o local ideal
para quem busca segurança através da inacessibilidade natural.

\subsubsection{Por que sim}

\begin{itemize}
\tightlist
\item
  Isolamento geográfico e demográfico superior (privacidade total).
\item
  Abundância ilimitada de água e recursos pesqueiros.
\item
  Ambiente social pacífico e seguro com presença militar de fronteira.
\item
  Longe de qualquer alvo militar terrestre ou eixo de fallout primário.
\end{itemize}

\subsubsection{Por que não}

\begin{itemize}
\tightlist
\item
  Risco hidrológico extremo (cheias dos rios Paraguai e Paraná).
\item
  Isolamento físico total por via terrestre em eventos climáticos.
\item
  Infraestrutura de saúde e comercial mínima.
\end{itemize}

\subsubsection{Recomendação
Técnica}

Altamente indicado para residências de isolamento tático que possuam
infraestrutura própria de transporte fluvial e autonomia de recursos.
Requer planejamento de construção em cotas de segurança contra
inundações históricas. O GSS de 6.3 reflete o poder do isolamento tático
equilibrado pela extrema vulnerabilidade hidrológica. Referencia atual
de score: GSS 6.3 em 2026-03-05.

\subsubsection{}

\begin{itemize}
\tightlist
\item
  \begin{enumerate}
  \def\labelenumi{\arabic{enumi})}
  \tightlist
  \item
    Dados censitários validados (INE\footnote{Instituto Nacional de Estadística (Instituto Nacional de Estatística), órgão oficial de dados demográficos e censitários.} 2022).
  \end{enumerate}
\item
  \begin{enumerate}
  \def\labelenumi{\arabic{enumi})}
  \setcounter{enumi}{1}
  \tightlist
  \item
    Registros históricos de cheias dos rios Paraguai e Paraná (SEN\footnote{Secretaría de Emergencia Nacional (Secretaria de Emergência Nacional), órgão governamental de resposta a desastres.}/\allowbreak DMH).
  \end{enumerate}
\item
  \begin{enumerate}
  \def\labelenumi{\arabic{enumi})}
  \setcounter{enumi}{2}
  \tightlist
  \item
    Mapa de sensibilidade geopolítica da confluência fluvial (MDN).
  \end{enumerate}
\item
  \begin{enumerate}
  \def\labelenumi{\arabic{enumi})}
  \setcounter{enumi}{3}
  \tightlist
  \item
    Registro de segurança pública e governança fluvial.
  \end{enumerate}
\end{itemize}}
\begin{table}[ht!]
\small \begin{tabular}{p{3.0cm}p{0.8cm}p{6.0cm}} \toprule
\textbf{Critério} & \textbf{Nota} & \textbf{Justificativa} \\ \midrule
A - Ameaças Estratégicas & 7.0 & Posicionamento geopolítico estratégico em confluência fluvial; remoto e discreto \\
B - Risco Social & 8.0 & População minúscula e densidade demográfica irrelevante; isolamento social superior \\
C - Riscos Naturais & 4.0 & Risco hidrológico extremo; nota penalizada pela vulnerabilidade a cheias dos dois grandes rios \\
D - Autossuficiência & 7.0 & Recursos hídricos e de proteína animal ilimitados; limitado pelo acesso logístico \\
E - Institucional & 5.0 & Ambiente social estável e seguro, mas com suporte institucional local mínimo \\
\midrule \textbf{GSS FINAL} & \textbf{6.3} & \textbf{Moderadamente Seguro (Ideal para quem busca isolamento tático fluvial e abundância de recursos hídricos, aceitando os riscos das cheias).} \\ \bottomrule \end{tabular}
\end{table}
\subsubsection{Vulnerabilidades e
Mitigação}\label{vuln-Paso_de_Patria-vulnerabilidades-e-mitigauxe7uxe3o}

\begin{itemize}
\tightlist
\item
  \textbf{Vulnerabilidade Hídrica Extrema:} Risco de isolamento total
  por via terrestre (Mitigação: posse de embarcação a motor como meio de
  transporte primário em crises; manutenção de estoque de suprimentos
  para 90 dias).
\item
  \textbf{Dependência de Pilar:} Vulnerabilidade por falta de serviços
  locais (Mitigação: instalação de sistema de comunicação via satélite e
  autonomia médica básica).
\item
  \textbf{Fragilidade Elétrica:} Rede final de linha (Mitigação:
  instalação de sistemas solares fotovoltaicos potentes e geradores a
  combustível estocado).
\end{itemize}
\subsubsection{Dossiê de Campo}\label{dossie-Paso_de_Patria-dossiuxea-de-campo}
\begin{description}[style=nextline,leftmargin=0.5cm,font=\bfseries]
\item[Coordenadas]  27°16'S 58°34'W.
\item[Topografia]  Relevo predominantemente plano, área de baixada situada na confluência estratégica dos rios Paraguai e Paraná. Altitude \textasciitilde 50m. Área de \textasciitilde 241 km². Geologicamente estável, risco sísmico desprezível.
\item[Alvos Estratégicos]  Confluência dos Rios Paraguai e Paraná (Ponto geopolítico e logístico vital da Bacia do Prata); Porto fluvial secundário; Local histórico (Teatro de operações da Guerra da Tríplice Aliança). Ausência de bases militares ativas de grande porte, mas localização de extrema sensibilidade para o controle fluvial e de fronteira.
\item[Fallout]  Ventos predominantes do Norte (NE) e Leste. Baixíssimo risco de fallout de alvos continentais primários devido ao isolamento geográfico extremo.
\item[População]  1.809 habitantes (Censo 2022 Final). Uma das menores densidades demográficas da Região Oriental.
\item[Densidade]  \textasciitilde 7,5 hab/\allowbreak km² (Densidade extremamente baixa, garantindo privacidade tática absoluta e ausência de riscos sociais urbanos).
\item[Vias de Acesso]  Acesso via ramais de terra a partir da Ruta PY04. Localização remota no extremo sudoeste do país, funcionando como um enclave de isolamento que favorece o anonimato estratégico.
\item[Serviços]  Infraestrutura básica mínima. Dependência absoluta de Pilar (\textasciitilde 60 km) para serviços de saúde especializados, logística comercial e suporte administrativo.
\item[Custo de Vida]  Muito baixo; economia baseada na pesca, pecuária extensiva e serviços de turismo histórico/\allowbreak pesca.
\item[Preço da Terra]  US\$ 1.500-2.500/\allowbreak ha (áreas com costa de rio); US\$ 800-1.200/\allowbreak ha (campos baixos de interior).
\item[Hidrologia]  Risco de Inundação Crítico. Situado no vértice do encontro de dois rios colossais. Vulnerável a cheias extraordinárias que podem inundar vastas áreas do distrito e isolar a localidade por terra por meses.
\item[Clima]  Subtropical Úmido. Verões muito quentes; invernos frescos com alta umidade influenciada pela massa de água dos rios.
\item[Energia]  Rede nacional (Itaipu/\allowbreak Yacyretá); instável e sujeita a cortes frequentes devido à localização em final de linha rural.
\item[Água]  Abundância hídrica ilimitada (Rios Paraguai e Paraná) e poços artesianos profundos. Água superficial exige tratamento para consumo humano.
\item[Qualidade do Solo]  Solo hidromórfico e arenoso; aptidão para pastagens naturais e pecuária. Baixa aptidão para agricultura diversificada de escala sem gestão hídrica.
\item[Resources Locais]  Grande excedente de gado bovino e recursos pesqueiros abundantes. Elevadíssimo potencial para autossuficiência calórica familiar baseada em proteína animal.
\item[Segurança]  Localidade excepcionalmente pacífica e segura. Criminalidade urbana inexistente. Estabilidade social baseada na cultura da pesca e tradição histórica. Presença de destacamentos da Armada Paraguaia para controle fluvial.
\item[Leis Local: Município consolidado. Restrição de Fronteira (Lei 2532/\allowbreak 05)]  Aplicável a terras rurais na faixa de 50km da fronteira argentina.
\end{description}
\paragraph{Solo (SoilGrids 2.0, média ponderada 0--30
cm)}

\begin{longtable}[]{@{}ll@{}}
\toprule\noalign{}
Parâmetro & Valor \\
\midrule\noalign{}
\endhead
\bottomrule\noalign{}
\endlastfoot
pH (H$_2$O) & 5.8 \\
Carbono orgânico (SOC) & 25.92 g/\allowbreak kg \\
Argila & 16.3 \% \\
Areia & 66.8 \% \\
Silte (calc.) & 16.9 \% \\
Densidade aparente & 1.28 g/\allowbreak cm³ \\
\textbf{Aptidão agrícola} & \textbf{Alta} \\
\end{longtable}

Fonte: ISRIC SoilGrids 2.0 via WCS (acesso em 2026-03-21). Coords:
-27.2667°, -58.5667°. Média ponderada camadas 0-5, 5-15, 15-30 cm.
\subsubsection{Indicadores Sociais}

\textbf{Fontes:} PNUD 2020, INE\footnote{Instituto Nacional de Estadística (Instituto Nacional de Estatística), órgão oficial de dados demográficos e censitários.} Censo 2022, INE\footnote{Instituto Nacional de Estadística (Instituto Nacional de Estatística), órgão oficial de dados demográficos e censitários.} EPHC 2023, Ministerio
Público 2024\\
\textbf{Nota:} Valores marcados como \emph{(dept.)} referem-se ao
departamento de Ñeembucú; valores \emph{(dist.)} são específicos deste
distrito. Dados marcados \emph{(est.)} são estimativas.

\begin{longtable}[]{@{}
  >{\raggedright\arraybackslash}p{(\linewidth - 4\tabcolsep) * \real{0.4231}}
  >{\raggedright\arraybackslash}p{(\linewidth - 4\tabcolsep) * \real{0.2692}}
  >{\raggedright\arraybackslash}p{(\linewidth - 4\tabcolsep) * \real{0.3077}}@{}}
\toprule\noalign{}
\begin{minipage}[b]{\linewidth}\raggedright
Indicador
\end{minipage} & \begin{minipage}[b]{\linewidth}\raggedright
Valor
\end{minipage} & \begin{minipage}[b]{\linewidth}\raggedright
Âmbito
\end{minipage} \\
\midrule\noalign{}
\endhead
\bottomrule\noalign{}
\endlastfoot
IDH (2020) & 0,710 & dept. \\
Ranking IDH nacional & 8/\allowbreak 18 & dept. \\
Esperança de vida & 73,0 anos & dept. \\
Escolaridade média & 7.3 anos & dist. \\
RNB per capita (USD PPA) & 8.800 & dept. \\
Pobreza monetária (\%) & 28,8\% & dept. \\
Pobreza extrema (\%) & 6,5\% & dept. \\
Índice de Gini & 0,445 & dept. \\
Acesso a água potável (\%) & 79,1\% & dept. \\
Acesso a saneamento (\%) & 77,4\% & dept. \\
Taxa de homicídios (est., /\allowbreak 100k hab) & 0,2--0,4 (2019--2020, fonte INE
--- mais baixa do país) & dept. \\
Índice de segurança & alto & dept. \\
População (Censo 2022) & 1.809 habitantes & dist. \\
Idade mediana & N/\allowbreak D & dist. \\
Taxa de fecundidade & N/\allowbreak D & dist. \\
\end{longtable}
\subsubsection{Combustível}

\textbf{Referência:} PETROPAR /\allowbreak  postos locais (2024) \textbf{Tipo de
localidade:} Interior

\begin{longtable}[]{@{}lll@{}}
\toprule\noalign{}
Combustível & USD/\allowbreak litro & Gs/\allowbreak litro (aprox.) \\
\midrule\noalign{}
\endhead
\bottomrule\noalign{}
\endlastfoot
Gasolina 93 oct & 0.99 & 7,326 \\
Gasolina 97 oct (premium) & 1.09 & 8,066 \\
Diesel & 0.92 & 6,808 \\
\end{longtable}

\begin{quote}
Preços podem variar ±5\% conforme posto e sazonalidade. Chaco e interior
remoto apresentam maior variação.
\end{quote}

\subsubsection{Cobertura Celular}

\textbf{Fonte:} CONATEL PY /\allowbreak  operadoras (2024)

\begin{longtable}[]{@{}ll@{}}
\toprule\noalign{}
Parâmetro & Valor \\
\midrule\noalign{}
\endhead
\bottomrule\noalign{}
\endlastfoot
Cobertura 4G população (dept.) & 75\% \\
Cobertura 4G área rural & 50\% \\
Melhor operadora & Tigo \\
Qualidade rural & moderada \\
\end{longtable}

\begin{quote}
Para áreas rurais fora do núcleo urbano, recomenda-se chip Tigo como
principal e Personal como backup.
\end{quote}

\subsubsection{Internet}

\textbf{Fonte:} CONATEL /\allowbreak  Speedtest Ookla (2024)

\begin{longtable}[]{@{}ll@{}}
\toprule\noalign{}
Parâmetro & Valor \\
\midrule\noalign{}
\endhead
\bottomrule\noalign{}
\endlastfoot
Velocidade média download & 30 Mbps \\
Domicílios com internet (dept.) & 40\% \\
Tecnologia predominante & rádio \\
Opção rural & Starlink disponível (\textasciitilde USD 44/\allowbreak mês) \\
\end{longtable}

\subsubsection{Mercado Imobiliário e Terra
Rural}

\textbf{Fonte:} INDERT /\allowbreak  Clasificados.com.py (2024)

\begin{longtable}[]{@{}ll@{}}
\toprule\noalign{}
Tipo & Referência \\
\midrule\noalign{}
\endhead
\bottomrule\noalign{}
\endlastfoot
Terra agrícola alta prod. (USD/\allowbreak ha) & 3,000 \\
Imóvel urbano (USD/\allowbreak m²) & 500 \\
Aluguel 2 quartos (USD/\allowbreak mês) & 200 \\
\end{longtable}

\begin{quote}
Valores de referência departamental. Localidades menores podem ter
preços 20--40\% abaixo da capital departamental. \#\#\# Saúde
\end{quote}

\textbf{Fonte:} MSPBS /\allowbreak  IPS Paraguay (2024-2026), consolidação
departamental e proxy local

\begin{longtable}[]{@{}ll@{}}
\toprule\noalign{}
Serviço & Disponibilidade \\
\midrule\noalign{}
\endhead
\bottomrule\noalign{}
\endlastfoot
USF /\allowbreak  Posto de Saúde & sim \\
Hospital Regional & não \\
IPS (seguro social) & não \\
Farmácia & sim \\
Distância ao hospital de referência & \textasciitilde62 km (Pilar) \\
\end{longtable}

\textbf{Principais estabelecimentos:} USF local; referencia hospitalar
em Pilar

\textbf{Observação para imigrantes:} Atendimento primario local ou em
raio curto; casos de maior complexidade seguem para Pilar. Cobertura
privada continua recomendada para especialidades e urgências de maior
complexidade.
\newpage
\secaoDiagnostico{Pilar}{\mapaConteudo{-26.8500}{-58.2833}}{Pilar é a ``capital da resiliência'' e um dos pontos mais seguros para
relocação estratégica no Paraguai. Oferece uma infraestrutura
institucional de primeiro nível, segurança social imbatível e um
isolamento tático natural garantido pela distância dos centros de
conflito e pela geografia dos humedales. Sua força absoluta reside na
combinação de suporte estatal (saúde/\allowbreak educação) com abundância de
recursos básicos (água/\allowbreak proteína). O único fator de atenção crítica é o
risco hídrico, que exige do ocupante a escolha de locais protegidos pelo
muro de defesa ou em cotas de segurança, transformando Pilar em uma
fortaleza logística segura no sudoeste do país.

\subsubsection{Por que sim}

\begin{itemize}
\tightlist
\item
  Melhor segurança pública entre as capitais departamentais.
\item
  Abundância absoluta de água, proteína animal e suporte institucional.
\item
  Isolamento geográfico estratégico e baixa visibilidade como alvo.
\item
  Conectividade asfáltica e fluvial de elite (Porto de águas profundas).
\end{itemize}

\subsubsection{Por que não}

\begin{itemize}
\tightlist
\item
  Risco crítico de inundações fluviais extraordinárias.
\item
  Dependência tática de infraestrutura de defesa costeira (muro/\allowbreak bombas).
\item
  Restrições da Lei de Fronteira em áreas rurais ribeirinhas.
\end{itemize}

\subsubsection{Recomendação
Técnica}

Altamente indicado como base administrativa, logística ou residência
tática de alto padrão. Requer autonomia energética individual para
mitigar riscos de rede. O GSS de 7.7 reflete a solidez das instituições
e a riqueza de recursos da localidade em equilíbrio com o desafio
hídrico. Referencia atual de score: GSS 7.7 em 2026-03-05.

\subsubsection{}

\begin{itemize}
\tightlist
\item
  \begin{enumerate}
  \def\labelenumi{\arabic{enumi})}
  \tightlist
  \item
    Dados censitários validados (INE\footnote{Instituto Nacional de Estadística (Instituto Nacional de Estatística), órgão oficial de dados demográficos e censitários.} 2022).
  \end{enumerate}
\item
  \begin{enumerate}
  \def\labelenumi{\arabic{enumi})}
  \setcounter{enumi}{1}
  \tightlist
  \item
    Relatórios de engenharia da Defesa Costeira de Pilar (MOPC 2025).
  \end{enumerate}
\item
  \begin{enumerate}
  \def\labelenumi{\arabic{enumi})}
  \setcounter{enumi}{2}
  \tightlist
  \item
    Mapa de profundidade e operação do Porto de Pilar (ANNP).
  \end{enumerate}
\item
  \begin{enumerate}
  \def\labelenumi{\arabic{enumi})}
  \setcounter{enumi}{3}
  \tightlist
  \item
    Registro de segurança pública e governança municipal.
  \end{enumerate}
\end{itemize}}
\begin{table}[ht!]
\small \begin{tabular}{p{3.0cm}p{0.8cm}p{6.0cm}} \toprule
\textbf{Critério} & \textbf{Nota} & \textbf{Justificativa} \\ \midrule
A - Ameaças Estratégicas & 8.0 & Capital departamental remota e discreta; isolamento tático superior e baixa visibilidade de alvos continentais \\
B - Risco Social & 7.5 & População gerenciável e densidade baixa; excelente infraestrutura institucional regional \\
C - Riscos Naturais & 5.5 & Risco crítico de inundação; nota penalizada pela dependência vital de sistemas de bombeamento \\
D - Autossuficiência & 8.5 & Recursos excepcionais: hub portuário, proteína animal abundante e suporte industrial têxtil \\
E - Institucional & 8.5 & Máxima segurança institucional e serviços de saúde de referência regional \\
\midrule \textbf{GSS FINAL} & \textbf{7.7} & \textbf{Seguro (Localidade de elite para suporte logístico e residência tática, desde que mitigados os riscos hídricos).} \\ \bottomrule \end{tabular}
\end{table}
\subsubsection{Vulnerabilidades e
Mitigação}\label{vuln-Pilar-vulnerabilidades-e-mitigauxe7uxe3o}

\begin{itemize}
\tightlist
\item
  \textbf{Vulnerabilidade Hídrica:} Dependência absoluta do muro de
  defesa (Mitigação: residência construída em áreas planejadas com
  drenagem superior ou em cotas elevadas fora do centro denso).
\item
  \textbf{Isolamento em Grandes Crises:} Risco de bloqueio das rotas
  PY04/\allowbreak PY19 (Mitigação: aproveitar a infraestrutura portuária para
  evacuação ou abastecimento fluvial).
\item
  \textbf{Dependência de Redes:} Vulnerabilidade a falhas elétricas em
  tempestades (Mitigação: instalação de sistemas solares fotovoltaicos
  potentes).
\end{itemize}
\subsubsection{Dossiê de Campo}\label{dossie-Pilar-dossiuxea-de-campo}
\begin{description}[style=nextline,leftmargin=0.5cm,font=\bfseries]
\item[Coordenadas]  26°51'S 58°17'W.
\item[Topografia]  Relevo predominantemente plano, margeado pelo Rio Paraguai e pelo arroio San Lorenzo. Situada em zona de humedales (terras baixas). Altitude \textasciitilde 55m. Área de \textasciitilde 1.130 km². Geologicamente estável, com histórico de sismos leves (2.0 em 2018) sem danos.
\item[Alvos Estratégicos]  Porto de Pilar (Terminal fluvial vital de águas profundas, operável mesmo em secas extremas); Muro de Defesa Costeira (Infraestrutura crítica de engenharia contra inundações); Comando da Armada e Polícia de Fronteira; Polo têxtil (Manufactura de Pilar). Localização estratégica para controle do tráfego fluvial na Bacia do Prata.
\item[Fallout]  Ventos predominantes do Norte (NE) e Leste. Baixíssimo risco de fallout de alvos continentais primários devido ao isolamento geográfico e distância de Assunção (\textasciitilde 360 km).
\item[População]  34.741 habitantes (Censo 2022 Final). Capital departamental e principal centro urbano do sudoeste paraguaio.
\item[Densidade]  \textasciitilde 30,7 hab/\allowbreak km² (Densidade baixa, oferecendo excelente equilíbrio entre infraestrutura urbana e isolamento social periférico).
\item[Vias de Acesso]  Conectividade asfáltica de elite via Ruta PY04 (Leste-Oeste) e a nova Ruta PY19 (Norte-Sul, ligando a Assunção). Porto fluvial de grande calado que garante logística internacional resiliente.
\item[Serviços]  Infraestrutura de saúde robusta (Hospital Regional de Pilar, IPS). Polo universitário (UNP), financeiro e administrativo regional. Serviços de comunicação e logística de primeiro nível para o interior do país.
\item[Custo de Vida]  Baixo-médio; economia baseada na indústria têxtil, pecuária e serviços portuários.
\item[Preço da Terra]  US\$ 2.000-3.500/\allowbreak ha (áreas com potencial agrícola/\allowbreak pecuário); US\$ 40-80/\allowbreak m² (Urbano).
\item[Hidrologia]  Risco Crítico de Inundação. Localidade historicamente vulnerável às cheias do Rio Paraguai. A segurança da capital depende da manutenção permanente do muro de defesa, estações de bombeamento e comportas.
\item[Clima]  Subtropical Úmido. Verões muito quentes e úmidos; invernos suaves com entradas de ar frio que limpam a atmosfera. Precipitação média 1.500 mm/\allowbreak ano.
\item[Energia]  Rede nacional (Itaipu/\allowbreak Yacyretá); nodo central de distribuição departamental com melhor estabilidade que o interior do departamento.
\item[Água]  Captação fluvial tratada (ESSAP\footnote{Empresa de Servicios Sanitarios del Paraguay (Empresa de Serviços Sanitários do Paraguai), responsável pelo saneamento e água urbana.}) e abundância de poços artesianos profundos (Aquífero Guarani).
\item[Qualidade do Solo]  Solo argiloso e hidromórfico; excelente aptidão para pecuária de corte e rizicultura (arroz).
\item[Resources Locais]  Grande polo produtor de carne bovina, peixes e tecidos. Elevado potencial para autossuficiência logística e de proteína animal em escala regional.
\item[Segurança]  Uma das capitais mais seguras e pacíficas do Paraguai. Baixíssima criminalidade urbana violenta. Forte coesão social comunitária e orgulho da identidade "Pilarense". Presença institucional sólida das forças armadas e polícia.
\item[Leis Local]  Capital departamental com regime de incentivo industrial. \\ \textbf{Livre de Restrição} de Fronteira para propriedades urbanas centrais, mas vigilância em zonas rurais ribeirinhas (Lei 2532/\allowbreak 05). \\ \textit{Aquisição de terra rural proibida para estrangeiros limítrofes.}
\end{description}
Coordenadas: -26.87°S /\allowbreak  -58.31°W. Acesso em 2026-03-20.

\textbf{Irradiância solar global --- ALLSKY SFC SW DWN (kWh/\allowbreak m²/\allowbreak dia):}

\begin{longtable}[]{@{}
  >{\raggedright\arraybackslash}p{(\linewidth - 22\tabcolsep) * \real{0.0833}}
  >{\raggedright\arraybackslash}p{(\linewidth - 22\tabcolsep) * \real{0.0833}}
  >{\raggedright\arraybackslash}p{(\linewidth - 22\tabcolsep) * \real{0.0833}}
  >{\raggedright\arraybackslash}p{(\linewidth - 22\tabcolsep) * \real{0.0833}}
  >{\raggedright\arraybackslash}p{(\linewidth - 22\tabcolsep) * \real{0.0833}}
  >{\raggedright\arraybackslash}p{(\linewidth - 22\tabcolsep) * \real{0.0833}}
  >{\raggedright\arraybackslash}p{(\linewidth - 22\tabcolsep) * \real{0.0833}}
  >{\raggedright\arraybackslash}p{(\linewidth - 22\tabcolsep) * \real{0.0833}}
  >{\raggedright\arraybackslash}p{(\linewidth - 22\tabcolsep) * \real{0.0833}}
  >{\raggedright\arraybackslash}p{(\linewidth - 22\tabcolsep) * \real{0.0833}}
  >{\raggedright\arraybackslash}p{(\linewidth - 22\tabcolsep) * \real{0.0833}}
  >{\raggedright\arraybackslash}p{(\linewidth - 22\tabcolsep) * \real{0.0833}}@{}}
\toprule\noalign{}
\begin{minipage}[b]{\linewidth}\raggedright
Jan
\end{minipage} & \begin{minipage}[b]{\linewidth}\raggedright
Fev
\end{minipage} & \begin{minipage}[b]{\linewidth}\raggedright
Mar
\end{minipage} & \begin{minipage}[b]{\linewidth}\raggedright
Abr
\end{minipage} & \begin{minipage}[b]{\linewidth}\raggedright
Mai
\end{minipage} & \begin{minipage}[b]{\linewidth}\raggedright
Jun
\end{minipage} & \begin{minipage}[b]{\linewidth}\raggedright
Jul
\end{minipage} & \begin{minipage}[b]{\linewidth}\raggedright
Ago
\end{minipage} & \begin{minipage}[b]{\linewidth}\raggedright
Set
\end{minipage} & \begin{minipage}[b]{\linewidth}\raggedright
Out
\end{minipage} & \begin{minipage}[b]{\linewidth}\raggedright
Nov
\end{minipage} & \begin{minipage}[b]{\linewidth}\raggedright
Dez
\end{minipage} \\
\midrule\noalign{}
\endhead
\bottomrule\noalign{}
\endlastfoot
6.76 & 6.22 & 5.34 & 4.20 & 3.30 & 2.73 & 3.14 & 3.88 & 4.59 & 5.48 &
6.49 & 6.76 \\
\end{longtable}

Média anual: 4.90 kWh/\allowbreak m²/\allowbreak dia. Mínimo: Jun (2.73) \textbar{} Máximo:
Jan/\allowbreak Dez (6.76).

\textbf{Precipitação --- PRECTOTCORR (mm/\allowbreak mês → mm/\allowbreak ano
\textasciitilde1.416):}

\begin{longtable}[]{@{}
  >{\raggedright\arraybackslash}p{(\linewidth - 22\tabcolsep) * \real{0.0833}}
  >{\raggedright\arraybackslash}p{(\linewidth - 22\tabcolsep) * \real{0.0833}}
  >{\raggedright\arraybackslash}p{(\linewidth - 22\tabcolsep) * \real{0.0833}}
  >{\raggedright\arraybackslash}p{(\linewidth - 22\tabcolsep) * \real{0.0833}}
  >{\raggedright\arraybackslash}p{(\linewidth - 22\tabcolsep) * \real{0.0833}}
  >{\raggedright\arraybackslash}p{(\linewidth - 22\tabcolsep) * \real{0.0833}}
  >{\raggedright\arraybackslash}p{(\linewidth - 22\tabcolsep) * \real{0.0833}}
  >{\raggedright\arraybackslash}p{(\linewidth - 22\tabcolsep) * \real{0.0833}}
  >{\raggedright\arraybackslash}p{(\linewidth - 22\tabcolsep) * \real{0.0833}}
  >{\raggedright\arraybackslash}p{(\linewidth - 22\tabcolsep) * \real{0.0833}}
  >{\raggedright\arraybackslash}p{(\linewidth - 22\tabcolsep) * \real{0.0833}}
  >{\raggedright\arraybackslash}p{(\linewidth - 22\tabcolsep) * \real{0.0833}}@{}}
\toprule\noalign{}
\begin{minipage}[b]{\linewidth}\raggedright
Jan
\end{minipage} & \begin{minipage}[b]{\linewidth}\raggedright
Fev
\end{minipage} & \begin{minipage}[b]{\linewidth}\raggedright
Mar
\end{minipage} & \begin{minipage}[b]{\linewidth}\raggedright
Abr
\end{minipage} & \begin{minipage}[b]{\linewidth}\raggedright
Mai
\end{minipage} & \begin{minipage}[b]{\linewidth}\raggedright
Jun
\end{minipage} & \begin{minipage}[b]{\linewidth}\raggedright
Jul
\end{minipage} & \begin{minipage}[b]{\linewidth}\raggedright
Ago
\end{minipage} & \begin{minipage}[b]{\linewidth}\raggedright
Set
\end{minipage} & \begin{minipage}[b]{\linewidth}\raggedright
Out
\end{minipage} & \begin{minipage}[b]{\linewidth}\raggedright
Nov
\end{minipage} & \begin{minipage}[b]{\linewidth}\raggedright
Dez
\end{minipage} \\
\midrule\noalign{}
\endhead
\bottomrule\noalign{}
\endlastfoot
154 & 121 & 145 & 167 & 108 & 71 & 48 & 40 & 68 & 156 & 177 & 162 \\
\end{longtable}

Estação chuvosa Out-Jan; seca Jul-Ago (mínimo em Ago).

\textbf{Inclinação solar recomendada:} 27° Norte (anual); 37° no inverno
(Jun-Ago); 17° no verão (Nov-Jan). Base: latitude -26.87°S.
\paragraph{Solo (SoilGrids 2.0, média ponderada 0--30
cm)}

\begin{longtable}[]{@{}ll@{}}
\toprule\noalign{}
Parâmetro & Valor \\
\midrule\noalign{}
\endhead
\bottomrule\noalign{}
\endlastfoot
pH (H$_2$O) & 5.9 \\
Carbono orgânico (SOC) & 12.3 g/\allowbreak kg \\
Argila & 18.32 \% \\
Areia & 56.8 \% \\
Silte (calc.) & 24.9 \% \\
Densidade aparente & 1.42 g/\allowbreak cm³ \\
\textbf{Aptidão agrícola} & \textbf{Alta} \\
\end{longtable}

Fonte: ISRIC SoilGrids 2.0 via WCS (acesso em 2026-03-21). Coords:
-26.85°, -58.2833°. Média ponderada camadas 0-5, 5-15, 15-30 cm.
\subsubsection{Indicadores Sociais}

\textbf{Fontes:} PNUD 2020, INE\footnote{Instituto Nacional de Estadística (Instituto Nacional de Estatística), órgão oficial de dados demográficos e censitários.} Censo 2022, INE\footnote{Instituto Nacional de Estadística (Instituto Nacional de Estatística), órgão oficial de dados demográficos e censitários.} EPHC 2023, Ministerio
Público 2024\\
\textbf{Nota:} Valores marcados como \emph{(dept.)} referem-se ao
departamento de Ñeembucú; valores \emph{(dist.)} são específicos deste
distrito. Dados marcados \emph{(est.)} são estimativas.

\begin{longtable}[]{@{}
  >{\raggedright\arraybackslash}p{(\linewidth - 4\tabcolsep) * \real{0.4231}}
  >{\raggedright\arraybackslash}p{(\linewidth - 4\tabcolsep) * \real{0.2692}}
  >{\raggedright\arraybackslash}p{(\linewidth - 4\tabcolsep) * \real{0.3077}}@{}}
\toprule\noalign{}
\begin{minipage}[b]{\linewidth}\raggedright
Indicador
\end{minipage} & \begin{minipage}[b]{\linewidth}\raggedright
Valor
\end{minipage} & \begin{minipage}[b]{\linewidth}\raggedright
Âmbito
\end{minipage} \\
\midrule\noalign{}
\endhead
\bottomrule\noalign{}
\endlastfoot
IDH (2020) & 0,710 & dept. \\
Ranking IDH nacional & 8/\allowbreak 18 & dept. \\
Esperança de vida & 73,0 anos & dept. \\
Escolaridade média & 10.8 anos & dist. \\
RNB per capita (USD PPA) & 8.800 & dept. \\
Pobreza monetária (\%) & 28,8\% & dept. \\
Pobreza extrema (\%) & 6,5\% & dept. \\
Índice de Gini & 0,445 & dept. \\
Acesso a água potável (\%) & 79,1\% & dept. \\
Acesso a saneamento (\%) & 77,4\% & dept. \\
Taxa de homicídios (est., /\allowbreak 100k hab) & 0,2--0,4 (2019--2020, fonte INE
--- mais baixa do país) & dept. \\
Índice de segurança & alto & dept. \\
População (Censo 2022) & 34.741 habitantes & dist. \\
Idade mediana & N/\allowbreak D & dist. \\
Taxa de fecundidade & N/\allowbreak D & dist. \\
\end{longtable}
\subsubsection{Combustível}

\textbf{Referência:} PETROPAR /\allowbreak  postos locais (2024) \textbf{Tipo de
localidade:} Capital departamental

\begin{longtable}[]{@{}lll@{}}
\toprule\noalign{}
Combustível & USD/\allowbreak litro & Gs/\allowbreak litro (aprox.) \\
\midrule\noalign{}
\endhead
\bottomrule\noalign{}
\endlastfoot
Gasolina 93 oct & 0.99 & 7,326 \\
Gasolina 97 oct (premium) & 1.09 & 8,066 \\
Diesel & 0.92 & 6,808 \\
\end{longtable}

\begin{quote}
Preços podem variar ±5\% conforme posto e sazonalidade. Chaco e interior
remoto apresentam maior variação.
\end{quote}

\subsubsection{Cobertura Celular}

\textbf{Fonte:} CONATEL PY /\allowbreak  operadoras (2024)

\begin{longtable}[]{@{}ll@{}}
\toprule\noalign{}
Parâmetro & Valor \\
\midrule\noalign{}
\endhead
\bottomrule\noalign{}
\endlastfoot
Cobertura 4G população (dept.) & 75\% \\
Cobertura 4G área rural & 50\% \\
Melhor operadora & Tigo \\
Qualidade rural & moderada \\
\end{longtable}

\begin{quote}
Para áreas rurais fora do núcleo urbano, recomenda-se chip Tigo como
principal e Personal como backup.
\end{quote}

\subsubsection{Internet}

\textbf{Fonte:} CONATEL /\allowbreak  Speedtest Ookla (2024)

\begin{longtable}[]{@{}ll@{}}
\toprule\noalign{}
Parâmetro & Valor \\
\midrule\noalign{}
\endhead
\bottomrule\noalign{}
\endlastfoot
Velocidade média download & 30 Mbps \\
Domicílios com internet (dept.) & 40\% \\
Tecnologia predominante & rádio \\
Opção rural & Starlink disponível (\textasciitilde USD 44/\allowbreak mês) \\
\end{longtable}

\subsubsection{Mercado Imobiliário e Terra
Rural}

\textbf{Fonte:} INDERT /\allowbreak  Clasificados.com.py (2024)

\begin{longtable}[]{@{}ll@{}}
\toprule\noalign{}
Tipo & Referência \\
\midrule\noalign{}
\endhead
\bottomrule\noalign{}
\endlastfoot
Terra agrícola alta prod. (USD/\allowbreak ha) & 3,000 \\
Imóvel urbano (USD/\allowbreak m²) & 575 \\
Aluguel 2 quartos (USD/\allowbreak mês) & 240 \\
\end{longtable}

\begin{quote}
Valores de referência departamental. Localidades menores podem ter
preços 20--40\% abaixo da capital departamental. \#\#\# Saúde
\end{quote}

\textbf{Fonte:} MSPBS /\allowbreak  IPS Paraguay (2024-2026), consolidação
departamental e proxy local

\begin{longtable}[]{@{}ll@{}}
\toprule\noalign{}
Serviço & Disponibilidade \\
\midrule\noalign{}
\endhead
\bottomrule\noalign{}
\endlastfoot
USF /\allowbreak  Posto de Saúde & sim \\
Hospital Regional & sim \\
IPS (seguro social) & sim \\
Farmácia & sim \\
Distância ao hospital de referência & local \\
\end{longtable}

\textbf{Principais estabelecimentos:} Hospital distrital /\allowbreak  regional de
Pilar; rede de USF e farmacias locais

\textbf{Observação para imigrantes:} Centro de referencia do
departamento. Acesso a atendimento primario e, em geral, a maior oferta
publica e privada da area. Cobertura privada continua recomendada para
especialidades e urgências de maior complexidade.

\paragraph{Solo (SoilGrids 2.0, média ponderada 0--30
cm)}

\begin{longtable}[]{@{}ll@{}}
\toprule\noalign{}
Parâmetro & Valor \\
\midrule\noalign{}
\endhead
\bottomrule\noalign{}
\endlastfoot
pH (H$_2$O) & 5.9 \\
Carbono orgânico (SOC) & 12.3 g/\allowbreak kg \\
Argila & 18.32 \% \\
Areia & 56.8 \% \\
Silte (calc.) & 24.9 \% \\
Densidade aparente & 1.42 g/\allowbreak cm³ \\
\textbf{Aptidão agrícola} & \textbf{Alta} \\
\end{longtable}

Fonte: ISRIC SoilGrids 2.0 via WCS (acesso em 2026-03-21). Coords:
-26.85°, -58.2833°. Média ponderada camadas 0-5, 5-15, 15-30 cm.
\newpage
\secaoDiagnostico{San Juan del Neembucu}{\mapaConteudo{-26.6500}{-57.9333}}{San Juan Bautista del Ñeembucú é o santuário pecuário do departamento.
Oferece um nível de isolamento demográfico que transforma a localidade
em um ``buraco negro'' estratégico, invisível a crises externas e
conflitos de massa. Sua força absoluta reside na abundância de proteína
animal e na segurança social imbatível de uma comunidade pequena e
resiliente. O desafio existencial é a dinâmica hídrica dos esteros, que
exige do ocupante autonomia total para operar em isolamento físico
sazonal. É o local ideal para quem busca segurança através da
inacessibilidade natural e deseja viver em um ambiente de paz social
absoluta.

\subsubsection{Por que sim}

\begin{itemize}
\tightlist
\item
  Máximo isolamento geográfico e demográfico (privacidade absoluta).
\item
  Ambiente social mais seguro do departamento pelo distanciamento
  físico.
\item
  Abundância massiva de proteína animal (gado).
\item
  Longe de qualquer alvo militar ou eixo de fallout tático.
\end{itemize}

\subsubsection{Por que não}

\begin{itemize}
\tightlist
\item
  Risco de isolamento físico total por semanas em cheias extremas.
\item
  Infraestrutura de saúde local mínima e dependência de Pilar.
\item
  Solo inadequado para agricultura mecanizada diversificada.
\end{itemize}

\subsubsection{Recomendação
Técnica}

Altamente indicado para estabelecimentos de refúgio tático que priorizem
o isolamento e a invisibilidade total. Requer autonomia energética e
logística de suprimentos robusta. O GSS de 7.7 reflete a superioridade
do isolamento protegida pela barreira natural dos humedales. Referencia
atual de score: GSS 7.7 em 2026-03-05.

\subsubsection{}

\begin{itemize}
\tightlist
\item
  \begin{enumerate}
  \def\labelenumi{\arabic{enumi})}
  \tightlist
  \item
    Dados censitários validados (INE\footnote{Instituto Nacional de Estadística (Instituto Nacional de Estatística), órgão oficial de dados demográficos e censitários.} 2022).
  \end{enumerate}
\item
  \begin{enumerate}
  \def\labelenumi{\arabic{enumi})}
  \setcounter{enumi}{1}
  \tightlist
  \item
    Relatórios de inundações sazonais da SEN\footnote{Secretaría de Emergencia Nacional (Secretaria de Emergência Nacional), órgão governamental de resposta a desastres.} em Ñeembucú.
  \end{enumerate}
\item
  \begin{enumerate}
  \def\labelenumi{\arabic{enumi})}
  \setcounter{enumi}{2}
  \tightlist
  \item
    Mapa de solos e esteros de Ñeembucú (Portal Geoestadístico).
  \end{enumerate}
\item
  \begin{enumerate}
  \def\labelenumi{\arabic{enumi})}
  \setcounter{enumi}{3}
  \tightlist
  \item
    Registro de segurança pública departamental.
  \end{enumerate}
\end{itemize}}
\begin{table}[ht!]
\small \begin{tabular}{p{3.0cm}p{0.8cm}p{6.0cm}} \toprule
\textbf{Critério} & \textbf{Nota} & \textbf{Justificativa} \\ \midrule
A - Ameaças Estratégicas & 8.0 & Isolamento tático de alto nível; centro geográfico dos esteros e baixíssima visibilidade estratégica \\
B - Risco Social & 8.5 & Densidade populacional extremamente baixa; segurança contra riscos de massa absoluta \\
C - Riscos Naturais & 6.0 & Risco hidrológico elevado; nota penalizada pela vulnerabilidade a inundações sazonais severas \\
D - Autossuficiência & 8.0 & Excelentes recursos de proteína animal e abundância hídrica superficial \\
E - Institucional & 7.5 & Ambiente social muito seguro e estável; suporte institucional funcional para o contexto rural \\
\midrule \textbf{GSS FINAL} & \textbf{7.7} & \textbf{Seguro (Altamente recomendado para quem busca isolamento tático radical e autossuficiência proteica, aceitando os desafios do isolamento hídrico).} \\ \bottomrule \end{tabular}
\end{table}
\subsubsection{Vulnerabilidades e
Mitigação}\label{vuln-San_Juan_del_Neembucu-vulnerabilidades-e-mitigauxe7uxe3o}

\begin{itemize}
\tightlist
\item
  \textbf{Isolamento Geográfico Severo:} Risco de ficar ``ilhada'' por
  terra por semanas (Mitigação: posse de veículo 4x4 robusto e trator;
  manutenção de estoque de suprimentos para 90 dias).
\item
  \textbf{Fragilidade das Redes:} Dependência de energia externa em área
  de difícil acesso (Mitigação: instalação de sistemas solares
  fotovoltaicos potentes off-grid).
\item
  \textbf{Ausência de Saúde Local:} Distância crítica de centros
  hospitalares (Mitigação: manutenção de estoque médico avançado e
  sistema de telemedicina satelital).
\end{itemize}
\subsubsection{Dossiê de Campo}\label{dossie-San_Juan_del_Neembucu-dossiuxea-de-campo}
\begin{description}[style=nextline,leftmargin=0.5cm,font=\bfseries]
\item[Coordenadas]  26°39'S 57°56'W.
\item[Topografia]  Relevo predominantemente plano, situado no coração geográfico do departamento, cercado por vastos esteros (pantanais) e arroios. Altitude \textasciitilde 60m. Área massiva de \textasciitilde 1.329 km². Geologicamente estável.
\item[Alvos Estratégicos]  Áreas de produção pecuária extensiva de grande porte; Local de isolamento tático no centro do sistema de humedales. Ausência total de alvos militares ou infraestrutura industrial massiva, oferecendo um dos maiores níveis de invisibilidade estratégica e proteção por anonimato na Região Oriental.
\item[Fallout]  Ventos predominantes do Norte (NE) e Leste. Baixíssimo risco de fallout direto; posição isolada e protegida pela dinâmica atmosférica regional.
\item[População]  3.803 habitantes (Censo 2022 Final). População dispersa em grandes estâncias e pequenas colônias rurais.
\item[Densidade]  \textasciitilde 2,9 hab/\allowbreak km² (Densidade extremamente baixa, garantindo privacidade absoluta e ausência de riscos sociais de massa urbana).
\item[Vias de Acesso]  Acesso via ramais de terra que conectam à Ruta PY04. Localização que favorece a discrição, mas sujeita a períodos de isolamento físico total durante grandes cheias ou chuvas intensas que saturam o solo dos esteros.
\item[Serviços]  Infraestrutura básica mínima. Dependência absoluta de Pilar (\textasciitilde 50 km) para atendimentos de saúde especializados, logística comercial e administrativa. A vila urbana é um mic núcleo de apoio à pecuária.
\item[Custo de Vida]  Muito baixo; economia baseada na pecuária bovina e subsistência rural.
\item[Preço da Terra]  US\$ 800-1.200/\allowbreak ha (campos baixos para pecuária extensiva).
\item[Hidrologia]  Risco de Inundação Crítico. Distrito situado no centro da bacia de drenagem dos esteros de Ñeembucú. O solo hidromórfico retém água por longos períodos, tornando as vias de terra intransitáveis em eventos de "El Niño". Abundância hídrica superficial permanente.
\item[Clima]  Subtropical Úmido. Verões intensos e chuvosos; invernos amenos com alta umidade e nevoeiros persistentes.
\item[Energia]  Rede nacional (Itaipu/\allowbreak Yacyretá); altamente instável devido ao isolamento da rede rural e dificuldade de manutenção de acessos.
\item[Água]  Abundante via poços artesianos e águas superficiais (esteros). Qualidade da água superficial exige tratamento para consumo humano.
\item[Qualidade do Solo]  Solo hidromórfico e arenoso; excelente para pastagens naturais e gado de corte. Agricultura mecanizada inviável sem projetos massivos de drenagem.
\item[Resources Locais]  Grande produção de gado bovino per capita. Elevadíssimo potencial para autossuficiência de proteína animal para a pequena população local.
\item[Segurança]  Localidade excepcionalmente pacífica e segura. Criminalidade urbana inexistente. Estabilidade social máxima baseada na resiliência da população tradicional aos desafios hídricos. Localidade "invisível" para qualquer conflito sociopolítico nacional.
\item[Leis Local]  Município consolidado. \\ \textbf{Livre de Restrição} de Fronteira. \\ \textit{Aquisição de terra rural proibida para estrangeiros limítrofes.}
\end{description}
\paragraph{Solo (SoilGrids 2.0, média ponderada 0--30
cm)}

\begin{longtable}[]{@{}ll@{}}
\toprule\noalign{}
Parâmetro & Valor \\
\midrule\noalign{}
\endhead
\bottomrule\noalign{}
\endlastfoot
pH (H$_2$O) & 6.24 \\
Carbono orgânico (SOC) & 14.08 g/\allowbreak kg \\
Argila & 20.45 \% \\
Areia & 54.84 \% \\
Silte (calc.) & 24.7 \% \\
Densidade aparente & 1.41 g/\allowbreak cm³ \\
\textbf{Aptidão agrícola} & \textbf{Alta} \\
\end{longtable}

Fonte: ISRIC SoilGrids 2.0 via WCS (acesso em 2026-03-21). Coords:
-26.65°, -57.9333°. Média ponderada camadas 0-5, 5-15, 15-30 cm.

\begin{itemize}
\tightlist
\item
  \textbf{Homicidios/\allowbreak 100k:} \textasciitilde0,0 (Histórico de paz
  absoluta).
\end{itemize}

\subsection{10. Analise de Riscos}

\begin{itemize}
\tightlist
\item
  \textbf{Cenario curto prazo:} Manutenção da tranquilidade rural e
  estabilidade da pecuária tradicional.
\item
  \textbf{Cenario medio prazo:} Impacto de mudanças no regime de chuvas
  (El Niño) sobre a conectividade do distrito.
\end{itemize}

\subsection{11. Redacao Final}

\subsubsection{Diagnostico Integrado}

San Juan Bautista del Ñeembucú é o santuário pecuário do departamento.
Oferece um nível de isolamento demográfico que transforma a localidade
em um ``buraco negro'' estratégico, invisível a crises externas e
conflitos de massa. Sua força absoluta reside na abundância de proteína
animal e na segurança social imbatível de uma comunidade pequena e
resiliente. O desafio existencial é a dinâmica hídrica dos esteros, que
exige do ocupante autonomia total para operar em isolamento físico
sazonal. É o local ideal para quem busca segurança através da
inacessibilidade natural e deseja viver em um ambiente de paz social
absoluta.

\subsubsection{Por que sim}

\begin{itemize}
\tightlist
\item
  Máximo isolamento geográfico e demográfico (privacidade absoluta).
\item
  Ambiente social mais seguro do departamento pelo distanciamento
  físico.
\item
  Abundância massiva de proteína animal (gado).
\item
  Longe de qualquer alvo militar ou eixo de fallout tático.
\end{itemize}

\subsubsection{Por que nao}

\begin{itemize}
\tightlist
\item
  Risco de isolamento físico total por semanas em cheias extremas.
\item
  Infraestrutura de saúde local mínima e dependência de Pilar.
\item
  Solo inadequado para agricultura mecanizada diversificada.
\end{itemize}

\subsubsection{Recomendacao Tecnica}

Altamente indicado para estabelecimentos de refúgio tático que priorizem
o isolamento e a invisibilidade total. Requer autonomia energética e
logística de suprimentos robusta. O GSS de 7.7 reflete a superioridade
do isolamento protegida pela barreira natural dos humedales. Referencia
atual de score: GSS 7.7 em 2026-03-05.

\subsection{13.}

\begin{itemize}
\tightlist
\item
  \begin{enumerate}
  \def\labelenumi{\arabic{enumi})}
  \tightlist
  \item
    Dados censitários validados (INE\footnote{Instituto Nacional de Estadística (Instituto Nacional de Estatística), órgão oficial de dados demográficos e censitários.} 2022).
  \end{enumerate}
\item
  \begin{enumerate}
  \def\labelenumi{\arabic{enumi})}
  \setcounter{enumi}{1}
  \tightlist
  \item
    Relatórios de inundações sazonais da SEN\footnote{Secretaría de Emergencia Nacional (Secretaria de Emergência Nacional), órgão governamental de resposta a desastres.} em Ñeembucú.
  \end{enumerate}
\item
  \begin{enumerate}
  \def\labelenumi{\arabic{enumi})}
  \setcounter{enumi}{2}
  \tightlist
  \item
    Mapa de solos e esteros de Ñeembucú (Portal Geoestadístico).
  \end{enumerate}
\item
  \begin{enumerate}
  \def\labelenumi{\arabic{enumi})}
  \setcounter{enumi}{3}
  \tightlist
  \item
    Registro de segurança pública departamental.
  \end{enumerate}
\end{itemize}

\subsubsection{Combustível}

\textbf{Referência:} PETROPAR /\allowbreak  postos locais (2024) \textbf{Tipo de
localidade:} Interior

\begin{longtable}[]{@{}lll@{}}
\toprule\noalign{}
Combustível & USD/\allowbreak litro & Gs/\allowbreak litro (aprox.) \\
\midrule\noalign{}
\endhead
\bottomrule\noalign{}
\endlastfoot
Gasolina 93 oct & 0.99 & 7,326 \\
Gasolina 97 oct (premium) & 1.09 & 8,066 \\
Diesel & 0.92 & 6,808 \\
\end{longtable}

\begin{quote}
Preços podem variar ±5\% conforme posto e sazonalidade. Chaco e interior
remoto apresentam maior variação.
\end{quote}

\subsubsection{Cobertura Celular}

\textbf{Fonte:} CONATEL PY /\allowbreak  operadoras (2024)

\begin{longtable}[]{@{}ll@{}}
\toprule\noalign{}
Parâmetro & Valor \\
\midrule\noalign{}
\endhead
\bottomrule\noalign{}
\endlastfoot
Cobertura 4G população (dept.) & 75\% \\
Cobertura 4G área rural & 50\% \\
Melhor operadora & Tigo \\
Qualidade rural & moderada \\
\end{longtable}

\begin{quote}
Para áreas rurais fora do núcleo urbano, recomenda-se chip Tigo como
principal e Personal como backup.
\end{quote}

\subsubsection{Internet}

\textbf{Fonte:} CONATEL /\allowbreak  Speedtest Ookla (2024)

\begin{longtable}[]{@{}ll@{}}
\toprule\noalign{}
Parâmetro & Valor \\
\midrule\noalign{}
\endhead
\bottomrule\noalign{}
\endlastfoot
Velocidade média download & 30 Mbps \\
Domicílios com internet (dept.) & 40\% \\
Tecnologia predominante & rádio \\
Opção rural & Starlink disponível (\textasciitilde USD 44/\allowbreak mês) \\
\end{longtable}

\subsubsection{Mercado Imobiliário e Terra
Rural}

\textbf{Fonte:} INDERT /\allowbreak  Clasificados.com.py (2024)

\begin{longtable}[]{@{}ll@{}}
\toprule\noalign{}
Tipo & Referência \\
\midrule\noalign{}
\endhead
\bottomrule\noalign{}
\endlastfoot
Terra agrícola alta prod. (USD/\allowbreak ha) & 3,000 \\
Imóvel urbano (USD/\allowbreak m²) & 500 \\
Aluguel 2 quartos (USD/\allowbreak mês) & 200 \\
\end{longtable}

\begin{quote}
Valores de referência departamental. Localidades menores podem ter
preços 20--40\% abaixo da capital departamental. \#\#\# Saúde
\end{quote}

\textbf{Fonte:} MSPBS /\allowbreak  IPS Paraguay (2024-2026), consolidação
departamental e proxy local

\begin{longtable}[]{@{}ll@{}}
\toprule\noalign{}
Serviço & Disponibilidade \\
\midrule\noalign{}
\endhead
\bottomrule\noalign{}
\endlastfoot
USF /\allowbreak  Posto de Saúde & sim \\
Hospital Regional & não \\
IPS (seguro social) & não \\
Farmácia & sim \\
Distância ao hospital de referência & \textasciitilde1 km (Pilar) \\
\end{longtable}

\textbf{Principais estabelecimentos:} USF local; referencia hospitalar
em Pilar

\textbf{Observação para imigrantes:} Atendimento primario local ou em
raio curto; casos de maior complexidade seguem para Pilar. Cobertura
privada continua recomendada para especialidades e urgências de maior
complexidade.
\subsubsection{Indicadores Sociais}

\textbf{Fontes:} PNUD 2020, INE\footnote{Instituto Nacional de Estadística (Instituto Nacional de Estatística), órgão oficial de dados demográficos e censitários.} Censo 2022, INE\footnote{Instituto Nacional de Estadística (Instituto Nacional de Estatística), órgão oficial de dados demográficos e censitários.} EPHC 2023, Ministerio
Público 2024\\
\textbf{Nota:} Valores marcados como \emph{(dept.)} referem-se ao
departamento de Ñeembucú; valores \emph{(dist.)} são específicos deste
distrito. Dados marcados \emph{(est.)} são estimativas.

\begin{longtable}[]{@{}
  >{\raggedright\arraybackslash}p{(\linewidth - 4\tabcolsep) * \real{0.4231}}
  >{\raggedright\arraybackslash}p{(\linewidth - 4\tabcolsep) * \real{0.2692}}
  >{\raggedright\arraybackslash}p{(\linewidth - 4\tabcolsep) * \real{0.3077}}@{}}
\toprule\noalign{}
\begin{minipage}[b]{\linewidth}\raggedright
Indicador
\end{minipage} & \begin{minipage}[b]{\linewidth}\raggedright
Valor
\end{minipage} & \begin{minipage}[b]{\linewidth}\raggedright
Âmbito
\end{minipage} \\
\midrule\noalign{}
\endhead
\bottomrule\noalign{}
\endlastfoot
IDH (2020) & 0,710 & dept. \\
Ranking IDH nacional & 8/\allowbreak 18 & dept. \\
Esperança de vida & 73,0 anos & dept. \\
Escolaridade média & 8,2 anos & dept. \\
RNB per capita (USD PPA) & 8.800 & dept. \\
Pobreza monetária (\%) & 28,8\% & dept. \\
Pobreza extrema (\%) & 6,5\% & dept. \\
Índice de Gini & 0,445 & dept. \\
Acesso a água potável (\%) & 79,1\% & dept. \\
Acesso a saneamento (\%) & 77,4\% & dept. \\
Taxa de homicídios (est., /\allowbreak 100k hab) & 0,2--0,4 (2019--2020, fonte INE
--- mais baixa do país) & dept. \\
Índice de segurança & alto & dept. \\
População (Censo 2022) & 3.803 habitantes & dist. \\
Idade mediana & N/\allowbreak D & dist. \\
Taxa de fecundidade & N/\allowbreak D & dist. \\
\end{longtable}
\subsubsection{Combustível}

\textbf{Referência:} PETROPAR /\allowbreak  postos locais (2024) \textbf{Tipo de
localidade:} Interior

\begin{longtable}[]{@{}lll@{}}
\toprule\noalign{}
Combustível & USD/\allowbreak litro & Gs/\allowbreak litro (aprox.) \\
\midrule\noalign{}
\endhead
\bottomrule\noalign{}
\endlastfoot
Gasolina 93 oct & 0.99 & 7,326 \\
Gasolina 97 oct (premium) & 1.09 & 8,066 \\
Diesel & 0.92 & 6,808 \\
\end{longtable}

\begin{quote}
Preços podem variar ±5\% conforme posto e sazonalidade. Chaco e interior
remoto apresentam maior variação.
\end{quote}

\subsubsection{Cobertura Celular}

\textbf{Fonte:} CONATEL PY /\allowbreak  operadoras (2024)

\begin{longtable}[]{@{}ll@{}}
\toprule\noalign{}
Parâmetro & Valor \\
\midrule\noalign{}
\endhead
\bottomrule\noalign{}
\endlastfoot
Cobertura 4G população (dept.) & 75\% \\
Cobertura 4G área rural & 50\% \\
Melhor operadora & Tigo \\
Qualidade rural & moderada \\
\end{longtable}

\begin{quote}
Para áreas rurais fora do núcleo urbano, recomenda-se chip Tigo como
principal e Personal como backup.
\end{quote}

\subsubsection{Internet}

\textbf{Fonte:} CONATEL /\allowbreak  Speedtest Ookla (2024)

\begin{longtable}[]{@{}ll@{}}
\toprule\noalign{}
Parâmetro & Valor \\
\midrule\noalign{}
\endhead
\bottomrule\noalign{}
\endlastfoot
Velocidade média download & 30 Mbps \\
Domicílios com internet (dept.) & 40\% \\
Tecnologia predominante & rádio \\
Opção rural & Starlink disponível (\textasciitilde USD 44/\allowbreak mês) \\
\end{longtable}

\subsubsection{Mercado Imobiliário e Terra
Rural}

\textbf{Fonte:} INDERT /\allowbreak  Clasificados.com.py (2024)

\begin{longtable}[]{@{}ll@{}}
\toprule\noalign{}
Tipo & Referência \\
\midrule\noalign{}
\endhead
\bottomrule\noalign{}
\endlastfoot
Terra agrícola alta prod. (USD/\allowbreak ha) & 3,000 \\
Imóvel urbano (USD/\allowbreak m²) & 500 \\
Aluguel 2 quartos (USD/\allowbreak mês) & 200 \\
\end{longtable}

\begin{quote}
Valores de referência departamental. Localidades menores podem ter
preços 20--40\% abaixo da capital departamental. \#\#\# Saúde
\end{quote}

\textbf{Fonte:} MSPBS /\allowbreak  IPS Paraguay (2024-2026), consolidação
departamental e proxy local

\begin{longtable}[]{@{}ll@{}}
\toprule\noalign{}
Serviço & Disponibilidade \\
\midrule\noalign{}
\endhead
\bottomrule\noalign{}
\endlastfoot
USF /\allowbreak  Posto de Saúde & sim \\
Hospital Regional & não \\
IPS (seguro social) & não \\
Farmácia & sim \\
Distância ao hospital de referência & \textasciitilde1 km (Pilar) \\
\end{longtable}

\textbf{Principais estabelecimentos:} USF local; referencia hospitalar
em Pilar

\textbf{Observação para imigrantes:} Atendimento primario local ou em
raio curto; casos de maior complexidade seguem para Pilar. Cobertura
privada continua recomendada para especialidades e urgências de maior
complexidade.
\newpage
\secaoDiagnostico{Tacuaras}{\mapaConteudo{-26.8500}{-58.1333}}{Tacuaras oferece uma combinação estratégica rara: o isolamento tático
absoluto de uma densidade demográfica ínfima com a conveniência
logística da Ruta PY04 pavimentada. Sua força reside na abundância de
proteína animal e na facilidade de deslocamento para o polo de Pilar em
tempos de normalidade. O desafio crítico é a saturação hídrica dos
esteros, que exige do ocupante a escolha de terrenos com boa drenagem e
autonomia energética. É o local ideal para quem busca anonimato e paz
social sem abdicar da conectividade rodoviária nacional.

\subsubsection{Por que sim}

\begin{itemize}
\tightlist
\item
  Densidade populacional extremamente baixa (privacidade total).
\item
  Acesso asfáltico direto via Ruta PY04 (Logística facilitada).
\item
  Ambiente social seguro, ordeiro e pacífico.
\item
  Abundância de proteína animal (gado).
\end{itemize}

\subsubsection{Por que não}

\begin{itemize}
\tightlist
\item
  Risco de inundações em áreas fora do eixo rodoviário.
\item
  Solo hidromórfico limita a agricultura diversificada.
\item
  Dependência de Pilar para serviços institucionais e de saúde.
\end{itemize}

\subsubsection{Recomendação
Técnica}

Altamente indicado para estabelecimentos de autossuficiência pecuária ou
rizicultura que busquem isolamento demográfico com boa acessibilidade.
Requer investimento em autonomia elétrica e saúde básica. O GSS de 6.3
reflete a segurança do isolamento equilibrada pela vulnerabilidade
hidrológica dos esteros. Referencia atual de score: GSS 6.3 em
2026-03-05.

\subsubsection{}

\begin{itemize}
\tightlist
\item
  \begin{enumerate}
  \def\labelenumi{\arabic{enumi})}
  \tightlist
  \item
    Dados censitários validados (INE\footnote{Instituto Nacional de Estadística (Instituto Nacional de Estatística), órgão oficial de dados demográficos e censitários.} 2022).
  \end{enumerate}
\item
  \begin{enumerate}
  \def\labelenumi{\arabic{enumi})}
  \setcounter{enumi}{1}
  \tightlist
  \item
    Relatórios de produtividade pecuária regional (MAG).
  \end{enumerate}
\item
  \begin{enumerate}
  \def\labelenumi{\arabic{enumi})}
  \setcounter{enumi}{2}
  \tightlist
  \item
    Mapa de solos e bacias hidrográficas de Ñeembucú (Portal
    Geoestadístico).
  \end{enumerate}
\item
  \begin{enumerate}
  \def\labelenumi{\arabic{enumi})}
  \setcounter{enumi}{3}
  \tightlist
  \item
    Registro de segurança pública departamental.
  \end{enumerate}
\end{itemize}}
\begin{table}[ht!]
\small \begin{tabular}{p{3.0cm}p{0.8cm}p{6.0cm}} \toprule
\textbf{Critério} & \textbf{Nota} & \textbf{Justificativa} \\ \midrule
A - Ameaças Estratégicas & 7.0 & Localização remota, mas com bom acesso via PY04; excelente invisibilidade tática \\
B - Risco Social & 8.0 & Densidade populacional extremamente baixa; isolamento social superior \\
C - Riscos Naturais & 4.5 & Risco crítico de inundação nos esteros; nota penalizada pela saturação do solo \\
D - Autossuficiência & 6.5 & Bons recursos de proteína animal; limitado pela falta de diversidade comercial local \\
E - Institucional & 5.0 & Ambiente social seguro e estável, mas com suporte institucional local mínimo \\
\midrule \textbf{GSS FINAL} & \textbf{6.3} & \textbf{Moderadamente Seguro (Ideal para quem busca isolamento rural com facilidade de acesso logístico via rodovia pavimentada).} \\ \bottomrule \end{tabular}
\end{table}
\subsubsection{Vulnerabilidades e
Mitigação}\label{vuln-Tacuaras-vulnerabilidades-e-mitigauxe7uxe3o}

\begin{itemize}
\tightlist
\item
  \textbf{Vulnerabilidade Hídrica de Interior:} Risco de isolamento das
  fazendas (Mitigação: posse de veículo 4x4 robusto e manutenção de
  estoque de suprimentos para 60 dias).
\item
  \textbf{Dependência de Serviços Externos:} Saúde e logística
  dependentes de Pilar (Mitigação: manutenção de estoque médico básico e
  kit trauma local).
\item
  \textbf{Fragilidade Elétrica:} Rede rural vulnerável (Mitigação:
  instalação de sistemas solares fotovoltaicos e geradores).
\end{itemize}
\subsubsection{Dossiê de Campo}\label{dossie-Tacuaras-dossiuxea-de-campo}
\begin{description}[style=nextline,leftmargin=0.5cm,font=\bfseries]
\item[Coordenadas]  26°51'S 58°08'W.
\item[Topografia]  Relevo plano com vasta predominância de esteros (pantanais) e terras úmidas de baixada. Altitude \textasciitilde 55m. Área extensa de \textasciitilde 1.000 km². Geologicamente muito estável, sem histórico de sismicidade.
\item[Alvos Estratégicos]  Áreas de produção pecuária extensiva; Ponto de passagem e fiscalização na Ruta PY04 (corredor vital Pilar-San Ignacio). Ausência total de alvos militares ou infraestrutura industrial, oferecendo excelente invisibilidade estratégica e anonimato geográfico.
\item[Fallout]  Ventos predominantes do Norte (NE) e Leste. Baixíssimo risco de fallout direto devido à ausência de alvos estratégicos na região.
\item[População]  2.717 habitantes (Censo 2022 Final). População estável e altamente dispersa em fazendas e pequenas comunidades.
\item[Densidade]  \textasciitilde 2,7 hab/\allowbreak km² (Densidade extremamente baixa, ideal para isolamento social tático e privacidade absoluta fora do eixo rodoviário).
\item[Vias de Acesso]  Localização privilegiada sobre a Ruta PY04 (pavimentada). Oferece excelente acessibilidade logística em tempos de normalidade, mantendo a característica de baixíssima ocupação humana no entorno. Estradas internas de terra são críticas em cheias.
\item[Serviços]  Unidade de Saúde da Família (USF) local. Dependência de Pilar (\textasciitilde 30 km) para atendimentos médicos de alta complexidade e logística comercial diversificada.
\item[Custo de Vida]  Muito baixo; economia baseada na pecuária e agricultura de subsistência.
\item[Preço da Terra]  US\$ 1.000-1.500/\allowbreak ha (campos baixos para pecuária e arroz).
\item[Hidrologia]  Risco de Inundação e Isolamento Rural. Distrito cortado por canais e esteros que transbordam sazonalmente. Vulnerabilidade moderada a alta em acessos rurais fora da Ruta PY04 durante eventos de "El Niño".
\item[Clima]  Subtropical Úmido. Verões intensos; invernos amenos com alta umidade relativa (média 1.500 mm/\allowbreak ano).
\item[Energia]  Rede nacional (Itaipu/\allowbreak Yacyretá); sujeita a interrupções frequentes em áreas rurais devido à exposição da rede.
\item[Água]  Abundância de águas superficiais e poços artesianos produtivos. Lençol freático superficial em grande parte do território.
\item[Qualidade do Solo]  Solo hidromórfico; excelente aptidão para pastagens naturais e pecuária de corte. Potencial emergente para rizicultura (arroz) tecnificada.
\item[Resources Locais]  Grande excedente de gado bovino e proteína animal. Elevado potencial para autossuficiência alimentar básica para a população local.
\item[Segurança]  Zona rural excepcionalmente pacífica. Baixíssima criminalidade urbana. Estabilidade social sustentada pela cultura pecuária tradicional e coesão das comunidades rurais. Ausência de visibilidade para conflitos sociais de massa.
\item[Leis Local: Município consolidado. Livre de Restrição de Fronteira]  Localizado fora da faixa de 50km da fronteira internacional seca (protegido pelos humedales).
\end{description}
\paragraph{Solo (SoilGrids 2.0, média ponderada 0--30
cm)}

\begin{longtable}[]{@{}ll@{}}
\toprule\noalign{}
Parâmetro & Valor \\
\midrule\noalign{}
\endhead
\bottomrule\noalign{}
\endlastfoot
pH (H$_2$O) & 6.15 \\
Carbono orgânico (SOC) & 12.26 g/\allowbreak kg \\
Argila & 17.3 \% \\
Areia & 62.94 \% \\
Silte (calc.) & 19.8 \% \\
Densidade aparente & 1.42 g/\allowbreak cm³ \\
\textbf{Aptidão agrícola} & \textbf{Média} \\
\end{longtable}

Fonte: ISRIC SoilGrids 2.0 via WCS (acesso em 2026-03-21). Coords:
-26.85°, -58.1333°. Média ponderada camadas 0-5, 5-15, 15-30 cm.
\subsubsection{Indicadores Sociais}

\textbf{Fontes:} PNUD 2020, INE\footnote{Instituto Nacional de Estadística (Instituto Nacional de Estatística), órgão oficial de dados demográficos e censitários.} Censo 2022, INE\footnote{Instituto Nacional de Estadística (Instituto Nacional de Estatística), órgão oficial de dados demográficos e censitários.} EPHC 2023, Ministerio
Público 2024\\
\textbf{Nota:} Valores marcados como \emph{(dept.)} referem-se ao
departamento de Ñeembucú; valores \emph{(dist.)} são específicos deste
distrito. Dados marcados \emph{(est.)} são estimativas.

\begin{longtable}[]{@{}
  >{\raggedright\arraybackslash}p{(\linewidth - 4\tabcolsep) * \real{0.4231}}
  >{\raggedright\arraybackslash}p{(\linewidth - 4\tabcolsep) * \real{0.2692}}
  >{\raggedright\arraybackslash}p{(\linewidth - 4\tabcolsep) * \real{0.3077}}@{}}
\toprule\noalign{}
\begin{minipage}[b]{\linewidth}\raggedright
Indicador
\end{minipage} & \begin{minipage}[b]{\linewidth}\raggedright
Valor
\end{minipage} & \begin{minipage}[b]{\linewidth}\raggedright
Âmbito
\end{minipage} \\
\midrule\noalign{}
\endhead
\bottomrule\noalign{}
\endlastfoot
IDH (2020) & 0,710 & dept. \\
Ranking IDH nacional & 8/\allowbreak 18 & dept. \\
Esperança de vida & 73,0 anos & dept. \\
Escolaridade média & 7.0 anos & dist. \\
RNB per capita (USD PPA) & 8.800 & dept. \\
Pobreza monetária (\%) & 28,8\% & dept. \\
Pobreza extrema (\%) & 6,5\% & dept. \\
Índice de Gini & 0,445 & dept. \\
Acesso a água potável (\%) & 79,1\% & dept. \\
Acesso a saneamento (\%) & 77,4\% & dept. \\
Taxa de homicídios (est., /\allowbreak 100k hab) & 0,2--0,4 (2019--2020, fonte INE
--- mais baixa do país) & dept. \\
Índice de segurança & alto & dept. \\
População (Censo 2022) & 2.717 habitantes & dist. \\
Idade mediana & N/\allowbreak D & dist. \\
Taxa de fecundidade & N/\allowbreak D & dist. \\
\end{longtable}
\subsubsection{Combustível}

\textbf{Referência:} PETROPAR /\allowbreak  postos locais (2024) \textbf{Tipo de
localidade:} Interior

\begin{longtable}[]{@{}lll@{}}
\toprule\noalign{}
Combustível & USD/\allowbreak litro & Gs/\allowbreak litro (aprox.) \\
\midrule\noalign{}
\endhead
\bottomrule\noalign{}
\endlastfoot
Gasolina 93 oct & 0.99 & 7,326 \\
Gasolina 97 oct (premium) & 1.09 & 8,066 \\
Diesel & 0.92 & 6,808 \\
\end{longtable}

\begin{quote}
Preços podem variar ±5\% conforme posto e sazonalidade. Chaco e interior
remoto apresentam maior variação.
\end{quote}

\subsubsection{Cobertura Celular}

\textbf{Fonte:} CONATEL PY /\allowbreak  operadoras (2024)

\begin{longtable}[]{@{}ll@{}}
\toprule\noalign{}
Parâmetro & Valor \\
\midrule\noalign{}
\endhead
\bottomrule\noalign{}
\endlastfoot
Cobertura 4G população (dept.) & 75\% \\
Cobertura 4G área rural & 50\% \\
Melhor operadora & Tigo \\
Qualidade rural & moderada \\
\end{longtable}

\begin{quote}
Para áreas rurais fora do núcleo urbano, recomenda-se chip Tigo como
principal e Personal como backup.
\end{quote}

\subsubsection{Internet}

\textbf{Fonte:} CONATEL /\allowbreak  Speedtest Ookla (2024)

\begin{longtable}[]{@{}ll@{}}
\toprule\noalign{}
Parâmetro & Valor \\
\midrule\noalign{}
\endhead
\bottomrule\noalign{}
\endlastfoot
Velocidade média download & 30 Mbps \\
Domicílios com internet (dept.) & 40\% \\
Tecnologia predominante & rádio \\
Opção rural & Starlink disponível (\textasciitilde USD 44/\allowbreak mês) \\
\end{longtable}

\subsubsection{Mercado Imobiliário e Terra
Rural}

\textbf{Fonte:} INDERT /\allowbreak  Clasificados.com.py (2024)

\begin{longtable}[]{@{}ll@{}}
\toprule\noalign{}
Tipo & Referência \\
\midrule\noalign{}
\endhead
\bottomrule\noalign{}
\endlastfoot
Terra agrícola alta prod. (USD/\allowbreak ha) & 3,000 \\
Imóvel urbano (USD/\allowbreak m²) & 500 \\
Aluguel 2 quartos (USD/\allowbreak mês) & 200 \\
\end{longtable}

\begin{quote}
Valores de referência departamental. Localidades menores podem ter
preços 20--40\% abaixo da capital departamental. \#\#\# Saúde
\end{quote}

\textbf{Fonte:} MSPBS /\allowbreak  IPS Paraguay (2024-2026), consolidação
departamental e proxy local

\begin{longtable}[]{@{}ll@{}}
\toprule\noalign{}
Serviço & Disponibilidade \\
\midrule\noalign{}
\endhead
\bottomrule\noalign{}
\endlastfoot
USF /\allowbreak  Posto de Saúde & sim \\
Hospital Regional & não \\
IPS (seguro social) & não \\
Farmácia & sim \\
Distância ao hospital de referência & \textasciitilde40 km (Pilar) \\
\end{longtable}

\textbf{Principais estabelecimentos:} USF local; referencia hospitalar
em Pilar

\textbf{Observação para imigrantes:} Atendimento primario local ou em
raio curto; casos de maior complexidade seguem para Pilar. Cobertura
privada continua recomendada para especialidades e urgências de maior
complexidade.
\newpage
\secaoDiagnostico{Villa Franca}{\mapaConteudo{-26.2833}{-58.2000}}{Villa Franca é o refúgio ``off-grid'' definitivo da Região Oriental.
Oferece um nível de privacidade e segurança social inigualável,
protegida por uma densidade demográfica quase nula e pela barreira
natural dos humedales e rios. Sua força reside no anonimato total; sua
fraqueza é a vulnerabilidade absoluta à dinâmica hídrica. Sobreviver em
Villa Franca exige um perfil de ocupante altamente autônomo, capaz de
operar de forma anfíbia e gerir sua própria subsistência e segurança sem
qualquer expectativa de suporte estatal, transformando a
inacessibilidade em sua maior defesa estratégica.

\subsubsection{Por que sim}

\begin{itemize}
\tightlist
\item
  Máximo isolamento demográfico e tático (privacidade absoluta).
\item
  Ambiente social mais seguro do país devido ao vazio populacional.
\item
  Abundância de proteína animal (gado) e recursos hídricos.
\item
  Longe de qualquer alvo militar ou eixo de fallout tático.
\end{itemize}

\subsubsection{Por que não}

\begin{itemize}
\tightlist
\item
  Risco hidrológico extremo (cheias extraordinárias).
\item
  Isolamento físico total por meses em eventos climáticos.
\item
  Inexistência de serviços de saúde e infraestrutura institucional.
\end{itemize}

\subsubsection{Recomendação
Técnica}

Altamente indicado para residências de isolamento tático extremo que
possuam infraestrutura própria de transporte fluvial e autonomia total
de recursos. Requer planejamento de construção elevado (palafitas
técnicas). O GSS de 6.7 reflete a superioridade do anonimato tático
protegida pela barreira intransponível da natureza. Referencia atual de
score: GSS 6.7 em 2026-03-05.

\subsubsection{}

\begin{itemize}
\tightlist
\item
  \begin{enumerate}
  \def\labelenumi{\arabic{enumi})}
  \tightlist
  \item
    Dados censitários validados (INE\footnote{Instituto Nacional de Estadística (Instituto Nacional de Estatística), órgão oficial de dados demográficos e censitários.} 2022).
  \end{enumerate}
\item
  \begin{enumerate}
  \def\labelenumi{\arabic{enumi})}
  \setcounter{enumi}{1}
  \tightlist
  \item
    Registros de inundações históricas da SEN\footnote{Secretaría de Emergencia Nacional (Secretaria de Emergência Nacional), órgão governamental de resposta a desastres.} em Ñeembucú.
  \end{enumerate}
\item
  \begin{enumerate}
  \def\labelenumi{\arabic{enumi})}
  \setcounter{enumi}{2}
  \tightlist
  \item
    Mapa de solos e esteros de Ñeembucú (Portal Geoestadístico).
  \end{enumerate}
\item
  \begin{enumerate}
  \def\labelenumi{\arabic{enumi})}
  \setcounter{enumi}{3}
  \tightlist
  \item
    Registro de segurança pública departamental.
  \end{enumerate}
\end{itemize}}
\begin{table}[ht!]
\small \begin{tabular}{p{3.0cm}p{0.8cm}p{6.0cm}} \toprule
\textbf{Critério} & \textbf{Nota} & \textbf{Justificativa} \\ \midrule
A - Ameaças Estratégicas & 8.5 & Isolamento tático extremo e ausência total de interesse estratégico/\allowbreak militar \\
B - Risco Social & 9.0 & Densidade populacional ínfima; segurança contra riscos de massa absoluta \\
C - Riscos Naturais & 3.5 & Risco hidrológico crítico; nota penalizada pela vulnerabilidade total a cheias \\
D - Autossuficiência & 6.5 & Vastidão de terras para gado; limitado pela inacessibilidade a mercados \\
E - Institucional & 4.5 & Ambiente social estável e pacífico, mas com suporte institucional quase inexistente \\
\midrule \textbf{GSS FINAL} & \textbf{6.7} & \textbf{Moderadamente Seguro (Indicado exclusivamente para projetos de isolamento tático extremo que busquem o máximo anonimato, aceitando a barreira das águas).} \\ \bottomrule \end{tabular}
\end{table}
\subsubsection{Vulnerabilidades e
Mitigação}\label{vuln-Villa_Franca-vulnerabilidades-e-mitigauxe7uxe3o}

\begin{itemize}
\tightlist
\item
  \textbf{Isolamento Geográfico Absoluto:} Risco de ficar
  ``desconectado'' por terra por longos períodos (Mitigação: posse
  obrigatória de embarcação a motor e trator; manutenção de estoque
  logístico para 120 dias).
\item
  \textbf{Inexistência de Saúde Local:} Distância crítica de qualquer
  suporte médico (Mitigação: manutenção de posto médico avançado privado
  e autonomia em emergências).
\item
  \textbf{Fragilidade Elétrica Extrema:} Rede final de linha (Mitigação:
  instalação de sistemas solares fotovoltaicos potentes off-grid).
\end{itemize}
\subsubsection{Dossiê de Campo}\label{dossie-Villa_Franca-dossiuxea-de-campo}
\begin{description}[style=nextline,leftmargin=0.5cm,font=\bfseries]
\item[Coordenadas]  26°17'S 58°12'W.
\item[Topografia]  Relevo plano com predominância absoluta de humedales, esteros e terras baixas inundáveis. Margeado pelo Rio Paraguai ao oeste e pelo Rio Tebicuary ao sul. Altitude \textasciitilde 55m. Área vasta de \textasciitilde 1.540 km². Geologicamente estável, risco sísmico nulo.
\item[Alvos Estratégicos]  Áreas de pecuária extensiva de grande porte; Porto fluvial secundário; Confluência regional dos rios Paraguai e Tebicuary. Ausência total de alvos militares ou infraestrutura industrial, oferecendo o nível máximo de invisibilidade estratégica e anonimato geográfico na Região Oriental.
\item[Fallout]  Ventos predominantes do Norte (NE) e Leste. Baixíssimo risco de fallout de qualquer alvo continental devido ao extremo isolamento geográfico.
\item[População]  1.138 habitantes (Censo 2022 Final). O segundo município menos povoado de todo o Paraguai.
\item[Densidade]  \textasciitilde 0,7 hab/\allowbreak km² (Densidade próxima de zero, garantindo privacidade absoluta e ausência total de riscos sociais de massa urbana).
\item[Vias de Acesso]  Acesso terrestre historicamente precário via ramais de terra a partir da Ruta PY19. Alto nível de isolamento físico durante a maior parte do ano, o que favorece o desaparecimento estratégico e o anonimato tático.
\item[Serviços]  Infraestrutura institucional e comercial quase inexistente. Dependência absoluta de Pilar (\textasciitilde 75 km) ou Alberdi para serviços de saúde básicos e logística. A vila urbana é um mic núcleo de subsistência.
\item[Custo de Vida]  Muito baixo; economia baseada na pecuária e no autoconsumo rural extremo.
\item[Preço da Terra]  US\$ 700-1.000/\allowbreak ha (campos baixos para pecuária extensiva). Um dos menores valores de terra do país.
\item[Hidrologia]  Risco de Inundação Crítico. Distrito situado em zona de baixada extrema e confluência de grandes rios. Altamente vulnerável a cheias extraordinárias que isolam a localidade por terra por meses. O solo hidromórfico retém umidade permanentemente.
\item[Clima]  Subtropical Úmido. Verões muito quentes e úmidos; invernos frescos com alta umidade e nevoeiros densos.
\item[Energia]  Rede nacional (Itaipu/\allowbreak Yacyretá); altamente instável e sujeita a falhas prolongadas devido ao isolamento da rede rural.
\item[Água]  Abundância hídrica superficial (rios e esteros), exigindo tratamento. Disponibilidade de poços artesianos produtivos.
\item[Qualidade do Solo]  Solo hidromórfico; aptidão exclusiva para pastagens naturais e gado de corte. Agricultura mecanizada inviável sem projetos massivos de engenharia.
\item[Resources Locais]  Grande excedente de gado bovino per capita. Elevadíssimo potencial para autossuficiência calórica familiar baseada em proteína animal para a ínfima população local.
\item[Segurança]  Zona rural excepcionalmente pacífica. Criminalidade urbana inexistente. Estabilidade social máxima baseada no isolamento geográfico e na resiliência da população tradicional. Localidade "invisível" para conflitos sociais de massa.
\item[Leis Local: Município consolidado. Restrição de Fronteira (Lei 2532/\allowbreak 05)]  Aplicável a terras rurais na faixa de 50km da fronteira argentina.
\end{description}
\paragraph{Solo (SoilGrids 2.0, média ponderada 0--30
cm)}

\begin{longtable}[]{@{}ll@{}}
\toprule\noalign{}
Parâmetro & Valor \\
\midrule\noalign{}
\endhead
\bottomrule\noalign{}
\endlastfoot
pH (H$_2$O) & 6.05 \\
Carbono orgânico (SOC) & 15.27 g/\allowbreak kg \\
Argila & 21.9 \% \\
Areia & 51.13 \% \\
Silte (calc.) & 27.0 \% \\
Densidade aparente & 1.37 g/\allowbreak cm³ \\
\textbf{Aptidão agrícola} & \textbf{Alta} \\
\end{longtable}

Fonte: ISRIC SoilGrids 2.0 via WCS (acesso em 2026-03-21). Coords:
-26.2833°, -58.2°. Média ponderada camadas 0-5, 5-15, 15-30 cm.

\begin{itemize}
\tightlist
\item
  \textbf{Homicidios/\allowbreak 100k:} \textasciitilde0,0 (Paz total).
\end{itemize}

\subsection{10. Analise de Riscos}

\begin{itemize}
\tightlist
\item
  \textbf{Cenario curto prazo:} Manutenção da tranquilidade rural
  absoluta e isolamento físico sazonal.
\item
  \textbf{Cenario medio prazo:} Impacto das mudanças climáticas na
  frequência de inundações extraordinárias do Rio Paraguai.
\end{itemize}

\subsection{11. Redacao Final}

\subsubsection{Diagnostico Integrado}

Villa Franca é o refúgio ``off-grid'' definitivo da Região Oriental.
Oferece um nível de privacidade e segurança social inigualável,
protegida por uma densidade demográfica quase nula e pela barreira
natural dos humedales e rios. Sua força reside no anonimato total; sua
fraqueza é a vulnerabilidade absoluta à dinâmica hídrica. Sobreviver em
Villa Franca exige um perfil de ocupante altamente autônomo, capaz de
operar de forma anfíbia e gerir sua própria subsistência e segurança sem
qualquer expectativa de suporte estatal, transformando a
inacessibilidade em sua maior defesa estratégica.

\subsubsection{Por que sim}

\begin{itemize}
\tightlist
\item
  Máximo isolamento demográfico e tático (privacidade absoluta).
\item
  Ambiente social mais seguro do país devido ao vazio populacional.
\item
  Abundância de proteína animal (gado) e recursos hídricos.
\item
  Longe de qualquer alvo militar ou eixo de fallout tático.
\end{itemize}

\subsubsection{Por que nao}

\begin{itemize}
\tightlist
\item
  Risco hidrológico extremo (cheias extraordinárias).
\item
  Isolamento físico total por meses em eventos climáticos.
\item
  Inexistência de serviços de saúde e infraestrutura institucional.
\end{itemize}

\subsubsection{Recomendacao Tecnica}

Altamente indicado para residências de isolamento tático extremo que
possuam infraestrutura própria de transporte fluvial e autonomia total
de recursos. Requer planejamento de construção elevado (palafitas
técnicas). O GSS de 6.7 reflete a superioridade do anonimato tático
protegida pela barreira intransponível da natureza. Referencia atual de
score: GSS 6.7 em 2026-03-05.

\subsection{13.}

\begin{itemize}
\tightlist
\item
  \begin{enumerate}
  \def\labelenumi{\arabic{enumi})}
  \tightlist
  \item
    Dados censitários validados (INE\footnote{Instituto Nacional de Estadística (Instituto Nacional de Estatística), órgão oficial de dados demográficos e censitários.} 2022).
  \end{enumerate}
\item
  \begin{enumerate}
  \def\labelenumi{\arabic{enumi})}
  \setcounter{enumi}{1}
  \tightlist
  \item
    Registros de inundações históricas da SEN\footnote{Secretaría de Emergencia Nacional (Secretaria de Emergência Nacional), órgão governamental de resposta a desastres.} em Ñeembucú.
  \end{enumerate}
\item
  \begin{enumerate}
  \def\labelenumi{\arabic{enumi})}
  \setcounter{enumi}{2}
  \tightlist
  \item
    Mapa de solos e esteros de Ñeembucú (Portal Geoestadístico).
  \end{enumerate}
\item
  \begin{enumerate}
  \def\labelenumi{\arabic{enumi})}
  \setcounter{enumi}{3}
  \tightlist
  \item
    Registro de segurança pública departamental.
  \end{enumerate}
\end{itemize}

\subsubsection{Combustível}

\textbf{Referência:} PETROPAR /\allowbreak  postos locais (2024) \textbf{Tipo de
localidade:} Interior

\begin{longtable}[]{@{}lll@{}}
\toprule\noalign{}
Combustível & USD/\allowbreak litro & Gs/\allowbreak litro (aprox.) \\
\midrule\noalign{}
\endhead
\bottomrule\noalign{}
\endlastfoot
Gasolina 93 oct & 0.99 & 7,326 \\
Gasolina 97 oct (premium) & 1.09 & 8,066 \\
Diesel & 0.92 & 6,808 \\
\end{longtable}

\begin{quote}
Preços podem variar ±5\% conforme posto e sazonalidade. Chaco e interior
remoto apresentam maior variação.
\end{quote}

\subsubsection{Cobertura Celular}

\textbf{Fonte:} CONATEL PY /\allowbreak  operadoras (2024)

\begin{longtable}[]{@{}ll@{}}
\toprule\noalign{}
Parâmetro & Valor \\
\midrule\noalign{}
\endhead
\bottomrule\noalign{}
\endlastfoot
Cobertura 4G população (dept.) & 75\% \\
Cobertura 4G área rural & 50\% \\
Melhor operadora & Tigo \\
Qualidade rural & moderada \\
\end{longtable}

\begin{quote}
Para áreas rurais fora do núcleo urbano, recomenda-se chip Tigo como
principal e Personal como backup.
\end{quote}

\subsubsection{Internet}

\textbf{Fonte:} CONATEL /\allowbreak  Speedtest Ookla (2024)

\begin{longtable}[]{@{}ll@{}}
\toprule\noalign{}
Parâmetro & Valor \\
\midrule\noalign{}
\endhead
\bottomrule\noalign{}
\endlastfoot
Velocidade média download & 30 Mbps \\
Domicílios com internet (dept.) & 40\% \\
Tecnologia predominante & rádio \\
Opção rural & Starlink disponível (\textasciitilde USD 44/\allowbreak mês) \\
\end{longtable}

\subsubsection{Mercado Imobiliário e Terra
Rural}

\textbf{Fonte:} INDERT /\allowbreak  Clasificados.com.py (2024)

\begin{longtable}[]{@{}ll@{}}
\toprule\noalign{}
Tipo & Referência \\
\midrule\noalign{}
\endhead
\bottomrule\noalign{}
\endlastfoot
Terra agrícola alta prod. (USD/\allowbreak ha) & 3,000 \\
Imóvel urbano (USD/\allowbreak m²) & 500 \\
Aluguel 2 quartos (USD/\allowbreak mês) & 200 \\
\end{longtable}

\begin{quote}
Valores de referência departamental. Localidades menores podem ter
preços 20--40\% abaixo da capital departamental. \#\#\# Saúde
\end{quote}

\textbf{Fonte:} MSPBS /\allowbreak  IPS Paraguay (2024-2026), consolidação
departamental e proxy local

\begin{longtable}[]{@{}ll@{}}
\toprule\noalign{}
Serviço & Disponibilidade \\
\midrule\noalign{}
\endhead
\bottomrule\noalign{}
\endlastfoot
USF /\allowbreak  Posto de Saúde & sim \\
Hospital Regional & não \\
IPS (seguro social) & não \\
Farmácia & sim \\
Distância ao hospital de referência & \textasciitilde78 km (Pilar) \\
\end{longtable}

\textbf{Principais estabelecimentos:} USF local; referencia hospitalar
em Pilar

\textbf{Observação para imigrantes:} Atendimento primario local ou em
raio curto; casos de maior complexidade seguem para Pilar. Cobertura
privada continua recomendada para especialidades e urgências de maior
complexidade.
\subsubsection{Indicadores Sociais}

\textbf{Fontes:} PNUD 2020, INE\footnote{Instituto Nacional de Estadística (Instituto Nacional de Estatística), órgão oficial de dados demográficos e censitários.} Censo 2022, INE\footnote{Instituto Nacional de Estadística (Instituto Nacional de Estatística), órgão oficial de dados demográficos e censitários.} EPHC 2023, Ministerio
Público 2024\\
\textbf{Nota:} Valores marcados como \emph{(dept.)} referem-se ao
departamento de Ñeembucú; valores \emph{(dist.)} são específicos deste
distrito. Dados marcados \emph{(est.)} são estimativas.

\begin{longtable}[]{@{}
  >{\raggedright\arraybackslash}p{(\linewidth - 4\tabcolsep) * \real{0.4231}}
  >{\raggedright\arraybackslash}p{(\linewidth - 4\tabcolsep) * \real{0.2692}}
  >{\raggedright\arraybackslash}p{(\linewidth - 4\tabcolsep) * \real{0.3077}}@{}}
\toprule\noalign{}
\begin{minipage}[b]{\linewidth}\raggedright
Indicador
\end{minipage} & \begin{minipage}[b]{\linewidth}\raggedright
Valor
\end{minipage} & \begin{minipage}[b]{\linewidth}\raggedright
Âmbito
\end{minipage} \\
\midrule\noalign{}
\endhead
\bottomrule\noalign{}
\endlastfoot
IDH (2020) & 0,710 & dept. \\
Ranking IDH nacional & 8/\allowbreak 18 & dept. \\
Esperança de vida & 73,0 anos & dept. \\
Escolaridade média & 8.2 anos & dist. \\
RNB per capita (USD PPA) & 8.800 & dept. \\
Pobreza monetária (\%) & 28,8\% & dept. \\
Pobreza extrema (\%) & 6,5\% & dept. \\
Índice de Gini & 0,445 & dept. \\
Acesso a água potável (\%) & 79,1\% & dept. \\
Acesso a saneamento (\%) & 77,4\% & dept. \\
Taxa de homicídios (est., /\allowbreak 100k hab) & 0,2--0,4 (2019--2020, fonte INE
--- mais baixa do país) & dept. \\
Índice de segurança & alto & dept. \\
População (Censo 2022) & 1.138 habitantes & dist. \\
Idade mediana & N/\allowbreak D & dist. \\
Taxa de fecundidade & N/\allowbreak D & dist. \\
\end{longtable}
\subsubsection{Combustível}

\textbf{Referência:} PETROPAR /\allowbreak  postos locais (2024) \textbf{Tipo de
localidade:} Interior

\begin{longtable}[]{@{}lll@{}}
\toprule\noalign{}
Combustível & USD/\allowbreak litro & Gs/\allowbreak litro (aprox.) \\
\midrule\noalign{}
\endhead
\bottomrule\noalign{}
\endlastfoot
Gasolina 93 oct & 0.99 & 7,326 \\
Gasolina 97 oct (premium) & 1.09 & 8,066 \\
Diesel & 0.92 & 6,808 \\
\end{longtable}

\begin{quote}
Preços podem variar ±5\% conforme posto e sazonalidade. Chaco e interior
remoto apresentam maior variação.
\end{quote}

\subsubsection{Cobertura Celular}

\textbf{Fonte:} CONATEL PY /\allowbreak  operadoras (2024)

\begin{longtable}[]{@{}ll@{}}
\toprule\noalign{}
Parâmetro & Valor \\
\midrule\noalign{}
\endhead
\bottomrule\noalign{}
\endlastfoot
Cobertura 4G população (dept.) & 75\% \\
Cobertura 4G área rural & 50\% \\
Melhor operadora & Tigo \\
Qualidade rural & moderada \\
\end{longtable}

\begin{quote}
Para áreas rurais fora do núcleo urbano, recomenda-se chip Tigo como
principal e Personal como backup.
\end{quote}

\subsubsection{Internet}

\textbf{Fonte:} CONATEL /\allowbreak  Speedtest Ookla (2024)

\begin{longtable}[]{@{}ll@{}}
\toprule\noalign{}
Parâmetro & Valor \\
\midrule\noalign{}
\endhead
\bottomrule\noalign{}
\endlastfoot
Velocidade média download & 30 Mbps \\
Domicílios com internet (dept.) & 40\% \\
Tecnologia predominante & rádio \\
Opção rural & Starlink disponível (\textasciitilde USD 44/\allowbreak mês) \\
\end{longtable}

\subsubsection{Mercado Imobiliário e Terra
Rural}

\textbf{Fonte:} INDERT /\allowbreak  Clasificados.com.py (2024)

\begin{longtable}[]{@{}ll@{}}
\toprule\noalign{}
Tipo & Referência \\
\midrule\noalign{}
\endhead
\bottomrule\noalign{}
\endlastfoot
Terra agrícola alta prod. (USD/\allowbreak ha) & 3,000 \\
Imóvel urbano (USD/\allowbreak m²) & 500 \\
Aluguel 2 quartos (USD/\allowbreak mês) & 200 \\
\end{longtable}

\begin{quote}
Valores de referência departamental. Localidades menores podem ter
preços 20--40\% abaixo da capital departamental. \#\#\# Saúde
\end{quote}

\textbf{Fonte:} MSPBS /\allowbreak  IPS Paraguay (2024-2026), consolidação
departamental e proxy local

\begin{longtable}[]{@{}ll@{}}
\toprule\noalign{}
Serviço & Disponibilidade \\
\midrule\noalign{}
\endhead
\bottomrule\noalign{}
\endlastfoot
USF /\allowbreak  Posto de Saúde & sim \\
Hospital Regional & não \\
IPS (seguro social) & não \\
Farmácia & sim \\
Distância ao hospital de referência & \textasciitilde78 km (Pilar) \\
\end{longtable}

\textbf{Principais estabelecimentos:} USF local; referencia hospitalar
em Pilar

\textbf{Observação para imigrantes:} Atendimento primario local ou em
raio curto; casos de maior complexidade seguem para Pilar. Cobertura
privada continua recomendada para especialidades e urgências de maior
complexidade.
\newpage
\secaoDiagnostico{Villalbin}{\mapaConteudo{-27.0833}{-58.2000}}{Villalbín é um enclave de profundo isolamento no sul de Ñeembucú.
Oferece um dos melhores níveis de privacidade demográfica do país, sendo
uma localidade ``invisível'' para qualquer crise externa ou conflito de
massa. Sua força reside no anonimato geográfico e na abundância de
recursos de proteína animal. O desafio crítico é a dinâmica hídrica dos
esteros, que exige do ocupante autonomia total para operar em isolamento
físico sazonal. É o local ideal para quem busca segurança através da
inacessibilidade e deseja viver em um ambiente de paz social absoluta.

\subsubsection{Por que sim}

\begin{itemize}
\tightlist
\item
  Extremo isolamento geográfico e demográfico (anonimato total).
\item
  Ambiente social mais seguro do país devido ao vazio demográfico.
\item
  Abundância de proteína animal (gado).
\item
  Longe de qualquer alvo militar ou eixo de fallout.
\end{itemize}

\subsubsection{Por que não}

\begin{itemize}
\tightlist
\item
  Risco de isolamento físico total por meses em cheias extremas.
\item
  Inexistência de serviços de saúde especializados no distrito.
\item
  Infraestrutura institucional e comercial mínima.
\end{itemize}

\subsubsection{Recomendação
Técnica}

Altamente indicado para estabelecimentos de refúgio tático que priorizem
o isolamento e o anonimato absoluto. Requer autonomia total em energia,
água e logística de saúde. O GSS de 6.0 reflete a segurança do
isolamento equilibrada pela fragilidade da infraestrutura. Referencia
atual de score: GSS 6.0 em 2026-03-05.

\subsubsection{}

\begin{itemize}
\tightlist
\item
  \begin{enumerate}
  \def\labelenumi{\arabic{enumi})}
  \tightlist
  \item
    Dados censitários validados (INE\footnote{Instituto Nacional de Estadística (Instituto Nacional de Estatística), órgão oficial de dados demográficos e censitários.} 2022).
  \end{enumerate}
\item
  \begin{enumerate}
  \def\labelenumi{\arabic{enumi})}
  \setcounter{enumi}{1}
  \tightlist
  \item
    Relatórios de inundações sazonais da SEN\footnote{Secretaría de Emergencia Nacional (Secretaria de Emergência Nacional), órgão governamental de resposta a desastres.} em Ñeembucú.
  \end{enumerate}
\item
  \begin{enumerate}
  \def\labelenumi{\arabic{enumi})}
  \setcounter{enumi}{2}
  \tightlist
  \item
    Mapa de solos e esteros de Ñeembucú (MAG).
  \end{enumerate}
\item
  \begin{enumerate}
  \def\labelenumi{\arabic{enumi})}
  \setcounter{enumi}{3}
  \tightlist
  \item
    Registro de segurança pública departamental.
  \end{enumerate}
\end{itemize}}
\begin{table}[ht!]
\small \begin{tabular}{p{3.0cm}p{0.8cm}p{6.0cm}} \toprule
\textbf{Critério} & \textbf{Nota} & \textbf{Justificativa} \\ \midrule
A - Ameaças Estratégicas & 7.0 & Localização remota e discreta; excelente invisibilidade tática e ausência de alvos \\
B - Risco Social & 8.0 & Densidade populacional extremamente baixa; isolamento social superior \\
C - Riscos Naturais & 4.0 & Risco hidrológico extremo; nota penalizada pela vulnerabilidade a inundações prolongadas \\
D - Autossuficiência & 6.0 & Recursos de gado; limitado pela inacessibilidade geográfica a insumos técnicos \\
E - Institucional & 4.0 & Ambiente social seguro e pacífico, mas com suporte institucional quase nulo \\
\midrule \textbf{GSS FINAL} & \textbf{6.0} & \textbf{Moderadamente Seguro (Ideal para quem busca isolamento tático "low profile" extremo, aceitando a precariedade logística).} \\ \bottomrule \end{tabular}
\end{table}
\subsubsection{Vulnerabilidades e
Mitigação}\label{vuln-Villalbin-vulnerabilidades-e-mitigauxe7uxe3o}

\begin{itemize}
\tightlist
\item
  \textbf{Isolamento Geográfico Sazonal:} Risco de ficar ``ilhada''
  durante inundações (Mitigação: posse de veículo 4x4 robusto e trator;
  manutenção de estoque de suprimentos para 60 dias).
\item
  \textbf{Fragilidade Elétrica:} Rede rural vulnerável (Mitigação:
  instalação de sistemas solares fotovoltaicos potentes off-grid).
\item
  \textbf{Inexistência de Saúde Local:} Distância crítica de hospitais
  (Mitigação: manutenção de estoque médico básico e kit trauma local).
\end{itemize}
\subsubsection{Dossiê de Campo}\label{dossie-Villalbin-dossiuxea-de-campo}
\begin{description}[style=nextline,leftmargin=0.5cm,font=\bfseries]
\item[Coordenadas]  27°05'S 58°12'W.
\item[Topografia]  Relevo predominantemente plano, caracterizado por áreas de baixada, humedales e arroios tributários da bacia do Rio Paraguai. Altitude \textasciitilde 55m. Área de \textasciitilde 388 km². Geologicamente estável, risco sísmico nulo.
\item[Alvos Estratégicos]  Áreas de pecuária extensiva; Ponto de conexão tática secundária entre Pilar e o extremo sul do departamento. Ausência total de infraestrutura industrial ou bases militares, oferecendo excelente invisibilidade estratégica e proteção por isolamento geográfico.
\item[Fallout]  Ventos predominantes do Norte (NE) e Leste. Baixíssimo risco de fallout direto de qualquer alvo continental devido ao isolamento geográfico.
\item[População]  1.613 habitantes (Censo 2022 Final). População muito reduzida, dispersa em ambiente estritamente rural.
\item[Densidade]  \textasciitilde 4,2 hab/\allowbreak km² (Densidade extremamente baixa, garantindo privacidade absoluta e ausência de riscos sociais de massa urbana).
\item[Vias de Acesso]  Acesso via ramais de terra a partir da Ruta PY04. Localização remota que favorece a discrição e o anonimato tático, mas sujeita a períodos de isolamento físico devido à precariedade das vias em épocas de chuvas intensas.
\item[Serviços]  Infraestrutura básica mínima. Dependência absoluta de Pilar (\textasciitilde 40 km) para atendimentos de saúde especializados, logística comercial e serviços administrativos.
\item[Custo de Vida]  Muito baixo; economia baseada na pecuária bovina e agricultura de subsistência.
\item[Preço da Terra]  US\$ 800-1.200/\allowbreak ha (campos baixos para pecuária).
\item[Hidrologia]  Risco de Inundação Crítico. Localizado em zona de terras baixas, altamente vulnerável a cheias sazonais e extraordinárias que podem comprometer os acessos terrestres por semanas e saturar o solo.
\item[Clima]  Subtropical Úmido. Verões intensos e úmidos; invernos amenos com alta umidade relativa e nevoeiros persistentes.
\item[Energia]  Rede nacional (Itaipu/\allowbreak Yacyretá); sujeita a interrupções frequentes devido à localização em rede rural isolada e exposição climática.
\item[Água]  Abundância de águas superficiais (esteros) e disponibilidade de poços artesianos produtivos. Águas superficiais exigem tratamento para consumo.
\item[Qualidade do Solo]  Solo hidromórfico e arenoso; aptidão para pastagens naturais e pecuária de corte. Baixa aptidão para agricultura mecanizada diversificada sem drenagem.
\item[Resources Locais]  Grande excedente de gado bovino per capita. Elevado potencial para autossuficiência calórica familiar baseada em proteína animal.
\item[Segurança]  Zona rural excepcionalmente pacífica e segura. Criminalidade urbana inexistente. Estabilidade social máxima sustentada pela cultura tradicional e resiliência comunitária. Ausência total de visibilidade para conflitos sociais.
\item[Leis Local: Município consolidado. Restrição de Fronteira (Lei 2532/\allowbreak 05)]  Aplicável a terras rurais na faixa de 50km da fronteira argentina (limite tático dos esteros).
\end{description}
\paragraph{Solo (SoilGrids 2.0, média ponderada 0--30
cm)}

\begin{longtable}[]{@{}ll@{}}
\toprule\noalign{}
Parâmetro & Valor \\
\midrule\noalign{}
\endhead
\bottomrule\noalign{}
\endlastfoot
pH (H$_2$O) & 6.15 \\
Carbono orgânico (SOC) & 12.64 g/\allowbreak kg \\
Argila & 13.63 \% \\
Areia & 65.54 \% \\
Silte (calc.) & 20.8 \% \\
Densidade aparente & 1.4 g/\allowbreak cm³ \\
\textbf{Aptidão agrícola} & \textbf{Média} \\
\end{longtable}

Fonte: ISRIC SoilGrids 2.0 via WCS (acesso em 2026-03-21). Coords:
-27.0833°, -58.2°. Média ponderada camadas 0-5, 5-15, 15-30 cm.

\begin{itemize}
\tightlist
\item
  \textbf{Homicidios/\allowbreak 100k:} \textasciitilde0,0 (Paz total).
\end{itemize}

\subsection{10. Analise de Riscos}

\begin{itemize}
\tightlist
\item
  \textbf{Cenario curto prazo:} Manutenção da tranquilidade rural e
  estabilidade da pecuária extensiva.
\item
  \textbf{Cenario medio prazo:} Impacto de mudanças climáticas na
  frequência de isolamento hídrico extraordinário.
\end{itemize}

\subsection{11. Redacao Final}

\subsubsection{Diagnostico Integrado}

Villalbín é um enclave de profundo isolamento no sul de Ñeembucú.
Oferece um dos melhores níveis de privacidade demográfica do país, sendo
uma localidade ``invisível'' para qualquer crise externa ou conflito de
massa. Sua força reside no anonimato geográfico e na abundância de
recursos de proteína animal. O desafio crítico é a dinâmica hídrica dos
esteros, que exige do ocupante autonomia total para operar em isolamento
físico sazonal. É o local ideal para quem busca segurança através da
inacessibilidade e deseja viver em um ambiente de paz social absoluta.

\subsubsection{Por que sim}

\begin{itemize}
\tightlist
\item
  Extremo isolamento geográfico e demográfico (anonimato total).
\item
  Ambiente social mais seguro do país devido ao vazio demográfico.
\item
  Abundância de proteína animal (gado).
\item
  Longe de qualquer alvo militar ou eixo de fallout.
\end{itemize}

\subsubsection{Por que nao}

\begin{itemize}
\tightlist
\item
  Risco de isolamento físico total por meses em cheias extremas.
\item
  Inexistência de serviços de saúde especializados no distrito.
\item
  Infraestrutura institucional e comercial mínima.
\end{itemize}

\subsubsection{Recomendacao Tecnica}

Altamente indicado para estabelecimentos de refúgio tático que priorizem
o isolamento e o anonimato absoluto. Requer autonomia total em energia,
água e logística de saúde. O GSS de 6.0 reflete a segurança do
isolamento equilibrada pela fragilidade da infraestrutura. Referencia
atual de score: GSS 6.0 em 2026-03-05.

\subsection{13.}

\begin{itemize}
\tightlist
\item
  \begin{enumerate}
  \def\labelenumi{\arabic{enumi})}
  \tightlist
  \item
    Dados censitários validados (INE\footnote{Instituto Nacional de Estadística (Instituto Nacional de Estatística), órgão oficial de dados demográficos e censitários.} 2022).
  \end{enumerate}
\item
  \begin{enumerate}
  \def\labelenumi{\arabic{enumi})}
  \setcounter{enumi}{1}
  \tightlist
  \item
    Relatórios de inundações sazonais da SEN\footnote{Secretaría de Emergencia Nacional (Secretaria de Emergência Nacional), órgão governamental de resposta a desastres.} em Ñeembucú.
  \end{enumerate}
\item
  \begin{enumerate}
  \def\labelenumi{\arabic{enumi})}
  \setcounter{enumi}{2}
  \tightlist
  \item
    Mapa de solos e esteros de Ñeembucú (MAG).
  \end{enumerate}
\item
  \begin{enumerate}
  \def\labelenumi{\arabic{enumi})}
  \setcounter{enumi}{3}
  \tightlist
  \item
    Registro de segurança pública departamental.
  \end{enumerate}
\end{itemize}

\subsubsection{Combustível}

\textbf{Referência:} PETROPAR /\allowbreak  postos locais (2024) \textbf{Tipo de
localidade:} Interior

\begin{longtable}[]{@{}lll@{}}
\toprule\noalign{}
Combustível & USD/\allowbreak litro & Gs/\allowbreak litro (aprox.) \\
\midrule\noalign{}
\endhead
\bottomrule\noalign{}
\endlastfoot
Gasolina 93 oct & 0.99 & 7,326 \\
Gasolina 97 oct (premium) & 1.09 & 8,066 \\
Diesel & 0.92 & 6,808 \\
\end{longtable}

\begin{quote}
Preços podem variar ±5\% conforme posto e sazonalidade. Chaco e interior
remoto apresentam maior variação.
\end{quote}

\subsubsection{Cobertura Celular}

\textbf{Fonte:} CONATEL PY /\allowbreak  operadoras (2024)

\begin{longtable}[]{@{}ll@{}}
\toprule\noalign{}
Parâmetro & Valor \\
\midrule\noalign{}
\endhead
\bottomrule\noalign{}
\endlastfoot
Cobertura 4G população (dept.) & 75\% \\
Cobertura 4G área rural & 50\% \\
Melhor operadora & Tigo \\
Qualidade rural & moderada \\
\end{longtable}

\begin{quote}
Para áreas rurais fora do núcleo urbano, recomenda-se chip Tigo como
principal e Personal como backup.
\end{quote}

\subsubsection{Internet}

\textbf{Fonte:} CONATEL /\allowbreak  Speedtest Ookla (2024)

\begin{longtable}[]{@{}ll@{}}
\toprule\noalign{}
Parâmetro & Valor \\
\midrule\noalign{}
\endhead
\bottomrule\noalign{}
\endlastfoot
Velocidade média download & 30 Mbps \\
Domicílios com internet (dept.) & 40\% \\
Tecnologia predominante & rádio \\
Opção rural & Starlink disponível (\textasciitilde USD 44/\allowbreak mês) \\
\end{longtable}

\subsubsection{Mercado Imobiliário e Terra
Rural}

\textbf{Fonte:} INDERT /\allowbreak  Clasificados.com.py (2024)

\begin{longtable}[]{@{}ll@{}}
\toprule\noalign{}
Tipo & Referência \\
\midrule\noalign{}
\endhead
\bottomrule\noalign{}
\endlastfoot
Terra agrícola alta prod. (USD/\allowbreak ha) & 3,000 \\
Imóvel urbano (USD/\allowbreak m²) & 500 \\
Aluguel 2 quartos (USD/\allowbreak mês) & 200 \\
\end{longtable}

\begin{quote}
Valores de referência departamental. Localidades menores podem ter
preços 20--40\% abaixo da capital departamental. \#\#\# Saúde
\end{quote}

\textbf{Fonte:} MSPBS /\allowbreak  IPS Paraguay (2024-2026), consolidação
departamental e proxy local

\begin{longtable}[]{@{}ll@{}}
\toprule\noalign{}
Serviço & Disponibilidade \\
\midrule\noalign{}
\endhead
\bottomrule\noalign{}
\endlastfoot
USF /\allowbreak  Posto de Saúde & sim \\
Hospital Regional & não \\
IPS (seguro social) & não \\
Farmácia & sim \\
Distância ao hospital de referência & \textasciitilde61 km (Pilar) \\
\end{longtable}

\textbf{Principais estabelecimentos:} USF local; referencia hospitalar
em Pilar

\textbf{Observação para imigrantes:} Atendimento primario local ou em
raio curto; casos de maior complexidade seguem para Pilar. Cobertura
privada continua recomendada para especialidades e urgências de maior
complexidade.
\subsubsection{Indicadores Sociais}

\textbf{Fontes:} PNUD 2020, INE\footnote{Instituto Nacional de Estadística (Instituto Nacional de Estatística), órgão oficial de dados demográficos e censitários.} Censo 2022, INE\footnote{Instituto Nacional de Estadística (Instituto Nacional de Estatística), órgão oficial de dados demográficos e censitários.} EPHC 2023, Ministerio
Público 2024\\
\textbf{Nota:} Valores marcados como \emph{(dept.)} referem-se ao
departamento de Ñeembucú; valores \emph{(dist.)} são específicos deste
distrito. Dados marcados \emph{(est.)} são estimativas.

\begin{longtable}[]{@{}
  >{\raggedright\arraybackslash}p{(\linewidth - 4\tabcolsep) * \real{0.4231}}
  >{\raggedright\arraybackslash}p{(\linewidth - 4\tabcolsep) * \real{0.2692}}
  >{\raggedright\arraybackslash}p{(\linewidth - 4\tabcolsep) * \real{0.3077}}@{}}
\toprule\noalign{}
\begin{minipage}[b]{\linewidth}\raggedright
Indicador
\end{minipage} & \begin{minipage}[b]{\linewidth}\raggedright
Valor
\end{minipage} & \begin{minipage}[b]{\linewidth}\raggedright
Âmbito
\end{minipage} \\
\midrule\noalign{}
\endhead
\bottomrule\noalign{}
\endlastfoot
IDH (2020) & 0,710 & dept. \\
Ranking IDH nacional & 8/\allowbreak 18 & dept. \\
Esperança de vida & 73,0 anos & dept. \\
Escolaridade média & 7.4 anos & dist. \\
RNB per capita (USD PPA) & 8.800 & dept. \\
Pobreza monetária (\%) & 28,8\% & dept. \\
Pobreza extrema (\%) & 6,5\% & dept. \\
Índice de Gini & 0,445 & dept. \\
Acesso a água potável (\%) & 79,1\% & dept. \\
Acesso a saneamento (\%) & 77,4\% & dept. \\
Taxa de homicídios (est., /\allowbreak 100k hab) & 0,2--0,4 (2019--2020, fonte INE
--- mais baixa do país) & dept. \\
Índice de segurança & alto & dept. \\
População (Censo 2022) & 1.613 habitantes & dist. \\
Idade mediana & N/\allowbreak D & dist. \\
Taxa de fecundidade & N/\allowbreak D & dist. \\
\end{longtable}
\subsubsection{Combustível}

\textbf{Referência:} PETROPAR /\allowbreak  postos locais (2024) \textbf{Tipo de
localidade:} Interior

\begin{longtable}[]{@{}lll@{}}
\toprule\noalign{}
Combustível & USD/\allowbreak litro & Gs/\allowbreak litro (aprox.) \\
\midrule\noalign{}
\endhead
\bottomrule\noalign{}
\endlastfoot
Gasolina 93 oct & 0.99 & 7,326 \\
Gasolina 97 oct (premium) & 1.09 & 8,066 \\
Diesel & 0.92 & 6,808 \\
\end{longtable}

\begin{quote}
Preços podem variar ±5\% conforme posto e sazonalidade. Chaco e interior
remoto apresentam maior variação.
\end{quote}

\subsubsection{Cobertura Celular}

\textbf{Fonte:} CONATEL PY /\allowbreak  operadoras (2024)

\begin{longtable}[]{@{}ll@{}}
\toprule\noalign{}
Parâmetro & Valor \\
\midrule\noalign{}
\endhead
\bottomrule\noalign{}
\endlastfoot
Cobertura 4G população (dept.) & 75\% \\
Cobertura 4G área rural & 50\% \\
Melhor operadora & Tigo \\
Qualidade rural & moderada \\
\end{longtable}

\begin{quote}
Para áreas rurais fora do núcleo urbano, recomenda-se chip Tigo como
principal e Personal como backup.
\end{quote}

\subsubsection{Internet}

\textbf{Fonte:} CONATEL /\allowbreak  Speedtest Ookla (2024)

\begin{longtable}[]{@{}ll@{}}
\toprule\noalign{}
Parâmetro & Valor \\
\midrule\noalign{}
\endhead
\bottomrule\noalign{}
\endlastfoot
Velocidade média download & 30 Mbps \\
Domicílios com internet (dept.) & 40\% \\
Tecnologia predominante & rádio \\
Opção rural & Starlink disponível (\textasciitilde USD 44/\allowbreak mês) \\
\end{longtable}

\subsubsection{Mercado Imobiliário e Terra
Rural}

\textbf{Fonte:} INDERT /\allowbreak  Clasificados.com.py (2024)

\begin{longtable}[]{@{}ll@{}}
\toprule\noalign{}
Tipo & Referência \\
\midrule\noalign{}
\endhead
\bottomrule\noalign{}
\endlastfoot
Terra agrícola alta prod. (USD/\allowbreak ha) & 3,000 \\
Imóvel urbano (USD/\allowbreak m²) & 500 \\
Aluguel 2 quartos (USD/\allowbreak mês) & 200 \\
\end{longtable}

\begin{quote}
Valores de referência departamental. Localidades menores podem ter
preços 20--40\% abaixo da capital departamental. \#\#\# Saúde
\end{quote}

\textbf{Fonte:} MSPBS /\allowbreak  IPS Paraguay (2024-2026), consolidação
departamental e proxy local

\begin{longtable}[]{@{}ll@{}}
\toprule\noalign{}
Serviço & Disponibilidade \\
\midrule\noalign{}
\endhead
\bottomrule\noalign{}
\endlastfoot
USF /\allowbreak  Posto de Saúde & sim \\
Hospital Regional & não \\
IPS (seguro social) & não \\
Farmácia & sim \\
Distância ao hospital de referência & \textasciitilde61 km (Pilar) \\
\end{longtable}

\textbf{Principais estabelecimentos:} USF local; referencia hospitalar
em Pilar

\textbf{Observação para imigrantes:} Atendimento primario local ou em
raio curto; casos de maior complexidade seguem para Pilar. Cobertura
privada continua recomendada para especialidades e urgências de maior
complexidade.
\newpage
\secaoDiagnostico{Villa Oliva}{\mapaConteudo{-26.0166}{-57.8500}}{Villa Oliva é um dos refúgios táticos mais promissores de Ñeembucú.
Oferece uma combinação única de isolamento demográfico extremo com a
conveniência de uma nova rodovia asfáltica (PY19) que garante logística
eficiente em tempos de paz e uma rota de escape rápida. Sua força
absoluta reside na abundância de proteína animal e na segurança social
imbatível de uma comunidade pequena e ordeira. O desafio principal é a
vulnerabilidade hídrica inerente aos humedales, o que exige do ocupante
a escolha de terrenos em cotas elevadas e a manutenção de autonomia
energética e médica para períodos de isolamento rural.

\subsubsection{Por que sim}

\begin{itemize}
\tightlist
\item
  Densidade populacional extremamente baixa (privacidade total).
\item
  Conectividade asfáltica moderna via Ruta PY19.
\item
  Ambiente social seguro, pacífico e ordeiro.
\item
  Abundância de proteína animal (gado e pesca).
\end{itemize}

\subsubsection{Por que não}

\begin{itemize}
\tightlist
\item
  Risco de inundações fluviais em áreas baixas.
\item
  Solo hidromórfico limita a diversificação agrícola mecanizada.
\item
  Infraestrutura de saúde local limitada a atendimentos primários.
\end{itemize}

\subsubsection{Recomendação
Técnica}

Altamente indicado para estabelecimentos de autossuficiência pecuária ou
residências de isolamento tático com boa mobilidade. Requer investimento
em autonomia elétrica e estoque logístico. O GSS de 6.8 reflete a força
do isolamento tático reforçada pela nova infraestrutura rodoviária.
Referencia atual de score: GSS 6.8 em 2026-03-05.

\subsubsection{}

\begin{itemize}
\tightlist
\item
  \begin{enumerate}
  \def\labelenumi{\arabic{enumi})}
  \tightlist
  \item
    Dados censitários validados (INE\footnote{Instituto Nacional de Estadística (Instituto Nacional de Estatística), órgão oficial de dados demográficos e censitários.} 2022).
  \end{enumerate}
\item
  \begin{enumerate}
  \def\labelenumi{\arabic{enumi})}
  \setcounter{enumi}{1}
  \tightlist
  \item
    Relatórios de pavimentação e elevação de cota da Ruta PY19 (MOPC).
  \end{enumerate}
\item
  \begin{enumerate}
  \def\labelenumi{\arabic{enumi})}
  \setcounter{enumi}{2}
  \tightlist
  \item
    Mapa de solos e aptidão para pastagens de Ñeembucú (Portal
    Geoestadístico).
  \end{enumerate}
\item
  \begin{enumerate}
  \def\labelenumi{\arabic{enumi})}
  \setcounter{enumi}{3}
  \tightlist
  \item
    Registro de segurança pública departamental.
  \end{enumerate}
\end{itemize}}
\begin{table}[ht!]
\small \begin{tabular}{p{3.0cm}p{0.8cm}p{6.0cm}} \toprule
\textbf{Critério} & \textbf{Nota} & \textbf{Justificativa} \\ \midrule
A - Ameaças Estratégicas & 8.0 & Isolamento tático estratégico com acesso asfaltado novo; excelente invisibilidade \\
B - Risco Social & 8.5 & Densidade populacional extremamente baixa; isolamento social de alto nível \\
C - Riscos Naturais & 4.5 & Risco crítico de inundação ribeirinha; nota penalizada pela saturação do solo \\
D - Autossuficiência & 7.0 & Excelentes recursos de proteína animal e água; limitado pela falta de comércio local \\
E - Institucional & 5.5 & Ambiente social seguro e estável, mas com serviços institucionais mínimos \\
\midrule \textbf{GSS FINAL} & \textbf{6.8} & \textbf{Seguro (Altamente recomendado para quem busca isolamento rural produtivo com boa conectividade logística moderna).} \\ \bottomrule \end{tabular}
\end{table}
\subsubsection{Vulnerabilidades e
Mitigação}\label{vuln-Villa_Oliva-vulnerabilidades-e-mitigauxe7uxe3o}

\begin{itemize}
\tightlist
\item
  \textbf{Isolamento Hídrico Rural:} Risco de isolamento das fazendas em
  grandes cheias (Mitigação: posse de veículo 4x4 robusto e trator;
  manutenção de estoque de suprimentos para 60 dias).
\item
  \textbf{Dependência de Serviços Externos:} Saúde e logística
  especializada em cidades distantes (Mitigação: manutenção de estoque
  médico básico e kit trauma local).
\item
  \textbf{Fragilidade Elétrica:} Rede rural vulnerável a intempéries
  (Mitigação: instalação de sistemas solares fotovoltaicos e geradores).
\end{itemize}
\subsubsection{Dossiê de Campo}\label{dossie-Villa_Oliva-dossiuxea-de-campo}
\begin{description}[style=nextline,leftmargin=0.5cm,font=\bfseries]
\item[Coordenadas]  26°01'S 57°51'W.
\item[Topografia]  Relevo predominantemente plano, área de baixada margeada pelo Rio Paraguai ao oeste. Presença de vastos humedales, esteros e campos de pastagem natural. Altitude \textasciitilde 60m. Área vasta de \textasciitilde 1.586 km². Geologicamente estável, risco sísmico nulo.
\item[Alvos Estratégicos]  Áreas de produção pecuária extensiva de grande porte; Porto fluvial secundário; Conexão logística via Ruta PY19 (Corredor Villeta-Alberdi). Ausência total de alvos militares ou infraestrutura industrial pesada, oferecendo excelente invisibilidade estratégica e proteção por isolamento geográfico.
\item[Fallout]  Ventos predominantes do Norte (NE) e Leste. Baixo risco de fallout direto; posição favorável em relação aos ventos dominantes vindos da zona metropolitana de Assunção (\textasciitilde 100 km).
\item[População]  3.377 habitantes (Censo 2022 Final). Perfil demográfico eminentemente rural e disperso.
\item[Densidade]  \textasciitilde 2,1 hab/\allowbreak km² (Densidade extremamente baixa, garantindo privacidade absoluta e baixíssimo risco social de massa).
\item[Vias de Acesso]  Localização privilegiada sobre a nova Ruta PY19 (pavimentada), que conecta o distrito ao polo de Villeta/\allowbreak Assunção ao norte e a Alberdi/\allowbreak Pilar ao sul. A nova rodovia reduziu o isolamento histórico, mantendo a localidade "fora do radar" populacional.
\item[Serviços]  Unidade de Saúde da Família (USF) local. Dependência de Villeta ou Pilar para atendimentos médicos de alta complexidade e logística avançada. A vila urbana oferece serviços comerciais básicos.
\item[Custo de Vida]  Muito baixo; economia baseada na pecuária bovina e pesca comercial/\allowbreak esportiva.
\item[Preço da Terra]  US\$ 1.500-2.500/\allowbreak ha (áreas com costa de rio ou aptidão para arroz); US\$ 800-1.200/\allowbreak ha (campos baixos para gado).
\item[Hidrologia]  Risco de Inundação Moderado a Alto. Vulnerável a cheias extraordinárias do Rio Paraguai e transbordamento de esteros internos. A nova rodovia PY19 possui cota elevada, o que garante acessibilidade mesmo em eventos hídricos moderados.
\item[Clima]  Subtropical Úmido. Verões intensos e úmidos; invernos suaves com entradas ocasionais de massas de ar polar que favorecem o controle de pragas.
\item[Energia]  Rede nacional (Itaipu/\allowbreak Yacyretá); instável em áreas rurais remotas durante temporais.
\item[Água]  Abundância hídrica superficial (Rio Paraguai e esteros) e poços artesianos produtivos. Excelente disponibilidade para pecuária.
\item[Qualidade do Solo]  Solo hidromórfico e arenoso; excelente para pastagens naturais e pecuária de corte. Potencial para rizicultura (arroz) em áreas baixas com gestão de água.
\item[Resources Locais]  Grande excedente de gado bovino e recursos pesqueiros. Elevado potencial para autossuficiência calórica familiar baseada em proteína animal para a pequena carga demográfica.
\item[Segurança]  Zona rural e urbana de pequeno porte excepcionalmente pacífica e segura. Baixíssima criminalidade urbana. Estabilidade social máxima sustentada pela tradição pecuária e pelo isolamento relativo. Ausência de visibilidade para conflitos sociais de massa.
\item[Leis Local: Município consolidado. Restrição de Fronteira (Lei 2532/\allowbreak 05)]  Aplicável a terras rurais na faixa de 50km da fronteira argentina (limite fluvial).
\end{description}
\paragraph{Solo (SoilGrids 2.0, média ponderada 0--30
cm)}

\begin{longtable}[]{@{}ll@{}}
\toprule\noalign{}
Parâmetro & Valor \\
\midrule\noalign{}
\endhead
\bottomrule\noalign{}
\endlastfoot
pH (H$_2$O) & 5.85 \\
Carbono orgânico (SOC) & 14.07 g/\allowbreak kg \\
Argila & 15.67 \% \\
Areia & 65.62 \% \\
Silte (calc.) & 18.7 \% \\
Densidade aparente & 1.37 g/\allowbreak cm³ \\
\textbf{Aptidão agrícola} & \textbf{Média} \\
\end{longtable}

Fonte: ISRIC SoilGrids 2.0 via WCS (acesso em 2026-03-21). Coords:
-26.0167°, -57.85°. Média ponderada camadas 0-5, 5-15, 15-30 cm.

\begin{itemize}
\tightlist
\item
  \textbf{Homicidios/\allowbreak 100k:} \textasciitilde4,0 (Muito abaixo da média
  nacional).
\end{itemize}

\subsection{10. Analise de Riscos}

\begin{itemize}
\tightlist
\item
  \textbf{Cenario curto prazo:} Manutenção da tranquilidade rural e
  valorização das terras pela nova acessibilidade.
\item
  \textbf{Cenario medio prazo:} Impacto de mudanças no regime de cheias
  do Rio Paraguai sobre a infraestrutura produtiva.
\end{itemize}

\subsection{11. Redacao Final}

\subsubsection{Diagnostico Integrado}

Villa Oliva é um dos refúgios táticos mais promissores de Ñeembucú.
Oferece uma combinação única de isolamento demográfico extremo com a
conveniência de uma nova rodovia asfáltica (PY19) que garante logística
eficiente em tempos de paz e uma rota de escape rápida. Sua força
absoluta reside na abundância de proteína animal e na segurança social
imbatível de uma comunidade pequena e ordeira. O desafio principal é a
vulnerabilidade hídrica inerente aos humedales, o que exige do ocupante
a escolha de terrenos em cotas elevadas e a manutenção de autonomia
energética e médica para períodos de isolamento rural.

\subsubsection{Por que sim}

\begin{itemize}
\tightlist
\item
  Densidade populacional extremamente baixa (privacidade total).
\item
  Conectividade asfáltica moderna via Ruta PY19.
\item
  Ambiente social seguro, pacífico e ordeiro.
\item
  Abundância de proteína animal (gado e pesca).
\end{itemize}

\subsubsection{Por que nao}

\begin{itemize}
\tightlist
\item
  Risco de inundações fluviais em áreas baixas.
\item
  Solo hidromórfico limita a diversificação agrícola mecanizada.
\item
  Infraestrutura de saúde local limitada a atendimentos primários.
\end{itemize}

\subsubsection{Recomendacao Tecnica}

Altamente indicado para estabelecimentos de autossuficiência pecuária ou
residências de isolamento tático com boa mobilidade. Requer investimento
em autonomia elétrica e estoque logístico. O GSS de 6.8 reflete a força
do isolamento tático reforçada pela nova infraestrutura rodoviária.
Referencia atual de score: GSS 6.8 em 2026-03-05.

\subsection{13.}

\begin{itemize}
\tightlist
\item
  \begin{enumerate}
  \def\labelenumi{\arabic{enumi})}
  \tightlist
  \item
    Dados censitários validados (INE\footnote{Instituto Nacional de Estadística (Instituto Nacional de Estatística), órgão oficial de dados demográficos e censitários.} 2022).
  \end{enumerate}
\item
  \begin{enumerate}
  \def\labelenumi{\arabic{enumi})}
  \setcounter{enumi}{1}
  \tightlist
  \item
    Relatórios de pavimentação e elevação de cota da Ruta PY19 (MOPC).
  \end{enumerate}
\item
  \begin{enumerate}
  \def\labelenumi{\arabic{enumi})}
  \setcounter{enumi}{2}
  \tightlist
  \item
    Mapa de solos e aptidão para pastagens de Ñeembucú (Portal
    Geoestadístico).
  \end{enumerate}
\item
  \begin{enumerate}
  \def\labelenumi{\arabic{enumi})}
  \setcounter{enumi}{3}
  \tightlist
  \item
    Registro de segurança pública departamental.
  \end{enumerate}
\end{itemize}

\subsubsection{Combustível}

\textbf{Referência:} PETROPAR /\allowbreak  postos locais (2024) \textbf{Tipo de
localidade:} Interior

\begin{longtable}[]{@{}lll@{}}
\toprule\noalign{}
Combustível & USD/\allowbreak litro & Gs/\allowbreak litro (aprox.) \\
\midrule\noalign{}
\endhead
\bottomrule\noalign{}
\endlastfoot
Gasolina 93 oct & 0.99 & 7,326 \\
Gasolina 97 oct (premium) & 1.09 & 8,066 \\
Diesel & 0.92 & 6,808 \\
\end{longtable}

\begin{quote}
Preços podem variar ±5\% conforme posto e sazonalidade. Chaco e interior
remoto apresentam maior variação.
\end{quote}

\subsubsection{Cobertura Celular}

\textbf{Fonte:} CONATEL PY /\allowbreak  operadoras (2024)

\begin{longtable}[]{@{}ll@{}}
\toprule\noalign{}
Parâmetro & Valor \\
\midrule\noalign{}
\endhead
\bottomrule\noalign{}
\endlastfoot
Cobertura 4G população (dept.) & 75\% \\
Cobertura 4G área rural & 50\% \\
Melhor operadora & Tigo \\
Qualidade rural & moderada \\
\end{longtable}

\begin{quote}
Para áreas rurais fora do núcleo urbano, recomenda-se chip Tigo como
principal e Personal como backup.
\end{quote}

\subsubsection{Internet}

\textbf{Fonte:} CONATEL /\allowbreak  Speedtest Ookla (2024)

\begin{longtable}[]{@{}ll@{}}
\toprule\noalign{}
Parâmetro & Valor \\
\midrule\noalign{}
\endhead
\bottomrule\noalign{}
\endlastfoot
Velocidade média download & 30 Mbps \\
Domicílios com internet (dept.) & 40\% \\
Tecnologia predominante & rádio \\
Opção rural & Starlink disponível (\textasciitilde USD 44/\allowbreak mês) \\
\end{longtable}

\subsubsection{Mercado Imobiliário e Terra
Rural}

\textbf{Fonte:} INDERT /\allowbreak  Clasificados.com.py (2024)

\begin{longtable}[]{@{}ll@{}}
\toprule\noalign{}
Tipo & Referência \\
\midrule\noalign{}
\endhead
\bottomrule\noalign{}
\endlastfoot
Terra agrícola alta prod. (USD/\allowbreak ha) & 3,000 \\
Imóvel urbano (USD/\allowbreak m²) & 500 \\
Aluguel 2 quartos (USD/\allowbreak mês) & 200 \\
\end{longtable}

\begin{quote}
Valores de referência departamental. Localidades menores podem ter
preços 20--40\% abaixo da capital departamental. \#\#\# Saúde
\end{quote}

\textbf{Fonte:} MSPBS /\allowbreak  IPS Paraguay (2024-2026), consolidação
departamental e proxy local

\begin{longtable}[]{@{}ll@{}}
\toprule\noalign{}
Serviço & Disponibilidade \\
\midrule\noalign{}
\endhead
\bottomrule\noalign{}
\endlastfoot
USF /\allowbreak  Posto de Saúde & sim \\
Hospital Regional & não \\
IPS (seguro social) & não \\
Farmácia & sim \\
Distância ao hospital de referência & \textasciitilde129 km (Pilar) \\
\end{longtable}

\textbf{Principais estabelecimentos:} USF local; referencia hospitalar
em Pilar

\textbf{Observação para imigrantes:} Atendimento primario local ou em
raio curto; casos de maior complexidade seguem para Pilar. Cobertura
privada continua recomendada para especialidades e urgências de maior
complexidade.
\subsubsection{Indicadores Sociais}

\textbf{Fontes:} PNUD 2020, INE\footnote{Instituto Nacional de Estadística (Instituto Nacional de Estatística), órgão oficial de dados demográficos e censitários.} Censo 2022, INE\footnote{Instituto Nacional de Estadística (Instituto Nacional de Estatística), órgão oficial de dados demográficos e censitários.} EPHC 2023, Ministerio
Público 2024\\
\textbf{Nota:} Valores marcados como \emph{(dept.)} referem-se ao
departamento de Ñeembucú; valores \emph{(dist.)} são específicos deste
distrito. Dados marcados \emph{(est.)} são estimativas.

\begin{longtable}[]{@{}
  >{\raggedright\arraybackslash}p{(\linewidth - 4\tabcolsep) * \real{0.4231}}
  >{\raggedright\arraybackslash}p{(\linewidth - 4\tabcolsep) * \real{0.2692}}
  >{\raggedright\arraybackslash}p{(\linewidth - 4\tabcolsep) * \real{0.3077}}@{}}
\toprule\noalign{}
\begin{minipage}[b]{\linewidth}\raggedright
Indicador
\end{minipage} & \begin{minipage}[b]{\linewidth}\raggedright
Valor
\end{minipage} & \begin{minipage}[b]{\linewidth}\raggedright
Âmbito
\end{minipage} \\
\midrule\noalign{}
\endhead
\bottomrule\noalign{}
\endlastfoot
IDH (2020) & 0,710 & dept. \\
Ranking IDH nacional & 8/\allowbreak 18 & dept. \\
Esperança de vida & 73,0 anos & dept. \\
Escolaridade média & 8.6 anos & dist. \\
RNB per capita (USD PPA) & 8.800 & dept. \\
Pobreza monetária (\%) & 28,8\% & dept. \\
Pobreza extrema (\%) & 6,5\% & dept. \\
Índice de Gini & 0,445 & dept. \\
Acesso a água potável (\%) & 79,1\% & dept. \\
Acesso a saneamento (\%) & 77,4\% & dept. \\
Taxa de homicídios (est., /\allowbreak 100k hab) & 0,2--0,4 (2019--2020, fonte INE
--- mais baixa do país) & dept. \\
Índice de segurança & alto & dept. \\
População (Censo 2022) & 3.377 habitantes & dist. \\
Idade mediana & N/\allowbreak D & dist. \\
Taxa de fecundidade & N/\allowbreak D & dist. \\
\end{longtable}
\subsubsection{Combustível}

\textbf{Referência:} PETROPAR /\allowbreak  postos locais (2024) \textbf{Tipo de
localidade:} Interior

\begin{longtable}[]{@{}lll@{}}
\toprule\noalign{}
Combustível & USD/\allowbreak litro & Gs/\allowbreak litro (aprox.) \\
\midrule\noalign{}
\endhead
\bottomrule\noalign{}
\endlastfoot
Gasolina 93 oct & 0.99 & 7,326 \\
Gasolina 97 oct (premium) & 1.09 & 8,066 \\
Diesel & 0.92 & 6,808 \\
\end{longtable}

\begin{quote}
Preços podem variar ±5\% conforme posto e sazonalidade. Chaco e interior
remoto apresentam maior variação.
\end{quote}

\subsubsection{Cobertura Celular}

\textbf{Fonte:} CONATEL PY /\allowbreak  operadoras (2024)

\begin{longtable}[]{@{}ll@{}}
\toprule\noalign{}
Parâmetro & Valor \\
\midrule\noalign{}
\endhead
\bottomrule\noalign{}
\endlastfoot
Cobertura 4G população (dept.) & 75\% \\
Cobertura 4G área rural & 50\% \\
Melhor operadora & Tigo \\
Qualidade rural & moderada \\
\end{longtable}

\begin{quote}
Para áreas rurais fora do núcleo urbano, recomenda-se chip Tigo como
principal e Personal como backup.
\end{quote}

\subsubsection{Internet}

\textbf{Fonte:} CONATEL /\allowbreak  Speedtest Ookla (2024)

\begin{longtable}[]{@{}ll@{}}
\toprule\noalign{}
Parâmetro & Valor \\
\midrule\noalign{}
\endhead
\bottomrule\noalign{}
\endlastfoot
Velocidade média download & 30 Mbps \\
Domicílios com internet (dept.) & 40\% \\
Tecnologia predominante & rádio \\
Opção rural & Starlink disponível (\textasciitilde USD 44/\allowbreak mês) \\
\end{longtable}

\subsubsection{Mercado Imobiliário e Terra
Rural}

\textbf{Fonte:} INDERT /\allowbreak  Clasificados.com.py (2024)

\begin{longtable}[]{@{}ll@{}}
\toprule\noalign{}
Tipo & Referência \\
\midrule\noalign{}
\endhead
\bottomrule\noalign{}
\endlastfoot
Terra agrícola alta prod. (USD/\allowbreak ha) & 3,000 \\
Imóvel urbano (USD/\allowbreak m²) & 500 \\
Aluguel 2 quartos (USD/\allowbreak mês) & 200 \\
\end{longtable}

\begin{quote}
Valores de referência departamental. Localidades menores podem ter
preços 20--40\% abaixo da capital departamental. \#\#\# Saúde
\end{quote}

\textbf{Fonte:} MSPBS /\allowbreak  IPS Paraguay (2024-2026), consolidação
departamental e proxy local

\begin{longtable}[]{@{}ll@{}}
\toprule\noalign{}
Serviço & Disponibilidade \\
\midrule\noalign{}
\endhead
\bottomrule\noalign{}
\endlastfoot
USF /\allowbreak  Posto de Saúde & sim \\
Hospital Regional & não \\
IPS (seguro social) & não \\
Farmácia & sim \\
Distância ao hospital de referência & \textasciitilde129 km (Pilar) \\
\end{longtable}

\textbf{Principais estabelecimentos:} USF local; referencia hospitalar
em Pilar

\textbf{Observação para imigrantes:} Atendimento primario local ou em
raio curto; casos de maior complexidade seguem para Pilar. Cobertura
privada continua recomendada para especialidades e urgências de maior
complexidade.
